% Les types de familles — Allemand 1ère A — Cours signé OnBuch+
% Programme officiel MINESEC. Compiler avec : tectonic ALA01-les-types-de-familles.tex
\newcommand{\DOCMATIERE}{Allemand}
\newcommand{\DOCNIVEAU}{1ère A}
\newcommand{\DOCMODULE}{Chapitre 1 : Vie familiale et sociale}
\newcommand{\DOCLECON}{ALA01}
\newcommand{\DOCTITRE}{Les types de familles}
% OnBuch+ — préambule « cours signés » ENRICHI (compilé avec tectonic / XeLaTeX)
% Système visuel « Pop » d'OnBuch+ : fond crème (jamais blanc), bordures encre
% épaisses, ombres « dures » décalées, titres Archivo Black en capitales.
% Version MATHÉMATIQUES : graphes (pgfplots/tikz), boîtes pédagogiques
% (définition, propriété, méthode, attention, le savais-tu, exercice résolu,
% exercice, TP, bilan), page de garde et signature OnBuch+ ancrée en bas.
%
% Macros attendues dans main.tex AVANT \input{preamble} :
%   \DOCMATIERE  \DOCNIVEAU  \DOCMODULE  \DOCLECON (numéro)  \DOCTITRE
\documentclass[11pt]{article}

\usepackage{fontspec}
\usepackage[ngerman,french]{babel} % français = langue principale ; l'allemand via \all{..} / textebox
\usepackage[a4paper,margin=2cm,top=3.4cm,bottom=2.8cm,headheight=1.4cm,footskip=1.3cm]{geometry}
\usepackage{amsmath,amssymb,amsfonts}
\usepackage{mathtools} % \xmapsto, \coloneqq... souvent utilisés par les modèles
\usepackage{cancel} % simplifications barrées (bilans de forces)
% \abs{z} (module, valeur absolue) et \norm{u} (norme), utilisés par les modèles
\providecommand{\abs}{}\let\abs\relax\DeclarePairedDelimiter{\abs}{\lvert}{\rvert}
\providecommand{\norm}{}\let\norm\relax\DeclarePairedDelimiter{\norm}{\lVert}{\rVert}
\usepackage[table]{xcolor} % [table] : fournit aussi \rowcolors (alternance de lignes)
\usepackage{pagecolor}
\usepackage{tikz}
\usetikzlibrary{arrows.meta,positioning,calc,shapes.geometric,shapes.misc,shapes.arrows,decorations.pathmorphing,decorations.pathreplacing,decorations.markings,patterns,shadows,mindmap,trees,fit,backgrounds,matrix,intersections,angles,quotes}
\usepackage{pgfplots}
\usepackage[version=4]{mhchem} % équations nucléaires \ce{^{235}_{92}U}
\usepackage[european,siunitx]{circuitikz} % schémas électriques
\pgfplotsset{compat=1.18}
\usepgfplotslibrary{fillbetween,groupplots} % aires entre courbes (calcul intégral)
\usepackage{enumitem}
\usepackage{fancyhdr}
% [most] charge toutes les bibliothèques tcolorbox (listings, minted, raster,
% documentation...) et gonfle inutilement la mémoire TeX sur un document long
% (~300 boîtes) : on ne charge que ce qui est réellement utilisé.
\usepackage[skins,breakable]{tcolorbox}
\usepackage{graphicx}
\usepackage{colortbl}
\usepackage{booktabs}
\usepackage{tabularx}
\usepackage{diagbox} % case d'en-tête diagonale des tableaux à double entrée
% Colonnes extensibles centrée / gauche / droite (C, L, R), souvent utilisées par les modèles
\newcolumntype{C}{>{\centering\arraybackslash}X}
\newcolumntype{L}{>{\raggedright\arraybackslash}X}
\newcolumntype{R}{>{\raggedleft\arraybackslash}X}
\usepackage{array}
\usepackage{multicol}
\usepackage{multirow}
\usepackage{siunitx}
% parse-numbers=false : accepte aussi \SI{2,5 \times 10^{-3}}{...}
\sisetup{locale=FR,output-decimal-marker={,},group-minimum-digits=4,parse-numbers=false}
\DeclareSIUnit\ohm{\ensuremath{\symup{\Omega}}} % U+2126 absent de la police
\DeclareSIUnit\division{div} % réglages d'oscilloscope (ms/div, V/div)
\DeclareSIUnit\tr{tr} % tours (vitesse de rotation, ex. \SI{1500}{\tr\per\minute})
\DeclareSIUnit\molar{M} % concentration molaire (ex. \SI{3}{\molar}, \SI{1}{\micro\molar})
\DeclareSIUnit\an{an} % année (le modèle confond parfois avec \year, un compteur TeX) — ex. \SI{5}{\kilo\meter\per\an}
\DeclareSIUnit\FCFA{FCFA} % franc CFA (ex. \SI{500}{\FCFA\per\kilo\gram})
\DeclareSIUnit\degreeBrix{\textdegree Brix} % teneur en sucre d'un sirop/jus (réfractométrie)
\DeclareSIUnit\month{mois} % le modèle utilise parfois \month, un compteur TeX, comme unité de durée
\usepackage{qrcode}
% \degree : commande fréquemment utilisée par les modèles pour un angle ou une
% température en dehors de \SI{}{} (non fournie par les paquets chargés).
\providecommand{\degree}{^{\circ}}
% \qed (fin de démonstration), utilisé par les modèles sans amsthm
\providecommand{\qed}{\hfill$\square$}
% \barre : double barre des valeurs interdites dans un tableau de variations (tkz-tab)
\providecommand{\barre}{\ensuremath{\|}}
% environnement proof (amsthm non chargé) utilisé par les modèles
\ifdefined\proof\else\newenvironment{proof}[1][Démonstration]{\par\noindent\textit{#1.}\ }{\qed\par}\fi
\usepackage{lastpage}
\usepackage{hyperref}
\hypersetup{hidelinks}

% ============================================================
% Palette « Pop » OnBuch+ — mêmes hex que la classe P (plus_ui.dart)
% ============================================================
\definecolor{poporange}{HTML}{F59321}
\definecolor{popdark}{HTML}{E07A0C}
\definecolor{popcream}{HTML}{FFF6E8}
\definecolor{popink}{HTML}{1C1714}
\definecolor{poppurple}{HTML}{7A5AE0}
\definecolor{popgreen}{HTML}{2E9E6A}
\definecolor{poppink}{HTML}{E0457B}
\definecolor{popblue}{HTML}{2F7DE1}
\definecolor{popgold}{HTML}{C98A12}
\definecolor{popmuted}{HTML}{8A7F76}
\definecolor{popline}{HTML}{EADFD2}
\definecolor{popgray}{HTML}{8A7F76}
\colorlet{popgrey}{popgray}
% Alias tolérants (le modèle nomme parfois les couleurs différemment)
\colorlet{popwhite}{white}
\colorlet{popblack}{popink}
\colorlet{popred}{poppink}
\colorlet{popyellow}{popgold}
% Teintes claires (fonds de tableaux/encadrés) — le modèle nomme parfois
% les couleurs avec un suffixe L (ex. popblueL) sans les avoir définies.
\colorlet{poporangeL}{poporange!20}
\colorlet{popdarkL}{popdark!20}
\colorlet{poppurpleL}{poppurple!20}
\colorlet{popgreenL}{popgreen!20}
\colorlet{poppinkL}{poppink!20}
\colorlet{popblueL}{popblue!20}
\colorlet{popgoldL}{popgold!20}
\colorlet{popredL}{popred!20}
\colorlet{popgrayL}{popgray!20}
% Teintes claires pour fonds de tableaux / zones
\colorlet{popblueL}{popblue!10!white}
\colorlet{poporangeL}{poporange!14!white}
\colorlet{popgreenL}{popgreen!12!white}
\colorlet{poppurpleL}{poppurple!12!white}
\colorlet{poppinkL}{poppink!12!white}
\colorlet{popgoldL}{popgold!14!white}

\pagecolor{popcream}

% ============================================================
% Typographie
% ============================================================
\defaultfontfeatures{Ligatures=TeX}
\setmainfont{PlusJakartaSans-Medium.ttf}[Path=../../../fonts/,BoldFont=PlusJakartaSans-Bold.ttf]
% Maths en sans-serif (Fira Math) avec chiffres et lettres droites de Plus
% Jakarta, pour que les formules se fondent dans le texte.
\usepackage[math-style=ISO,bold-style=ISO]{unicode-math}
\setmathfont{FiraMath-Regular.otf}[Path=../../../fonts/]
\setmathfont{PlusJakartaSans-Medium.ttf}[Path=../../../fonts/,range={up/num,up/{Latin,latin},\mathup}]
\usepackage{icomma} % virgule décimale sans espace en mode math
% Caractères mathématiques Unicode absents des polices de texte (Plus Jakarta,
% Archivo Black) : rendus par leur équivalent mathématique, en texte comme en
% formule — sinon ils s'affichent en carré vide (ex. « Arithmétique dans ℤ »).
\usepackage{newunicodechar}
\newunicodechar{ℝ}{\ensuremath{\mathbb{R}}}
\newunicodechar{ℤ}{\ensuremath{\mathbb{Z}}}
\newunicodechar{ℂ}{\ensuremath{\mathbb{C}}}
\newunicodechar{ℕ}{\ensuremath{\mathbb{N}}}
\newunicodechar{ℚ}{\ensuremath{\mathbb{Q}}}
\newunicodechar{𝔻}{\ensuremath{\mathbb{D}}}
\newunicodechar{α}{\ensuremath{\alpha}}
\newunicodechar{β}{\ensuremath{\beta}}
\newunicodechar{γ}{\ensuremath{\gamma}}
\newunicodechar{δ}{\ensuremath{\delta}}
\newunicodechar{ε}{\ensuremath{\varepsilon}}
\newunicodechar{θ}{\ensuremath{\theta}}
\newunicodechar{λ}{\ensuremath{\lambda}}
\newunicodechar{μ}{\ensuremath{\mu}}
\newunicodechar{σ}{\ensuremath{\sigma}}
\newunicodechar{τ}{\ensuremath{\tau}}
\newunicodechar{φ}{\ensuremath{\varphi}}
\newunicodechar{ψ}{\ensuremath{\psi}}
\newunicodechar{ω}{\ensuremath{\omega}}
\newunicodechar{Σ}{\ensuremath{\Sigma}}
% majuscules produites par \MakeUppercase dans les titres (α-aminés -> Α)
\newunicodechar{Α}{\ensuremath{\alpha}}
\newunicodechar{Β}{\ensuremath{\beta}}
\newunicodechar{Γ}{\ensuremath{\Gamma}}
\newunicodechar{Δ}{\ensuremath{\Delta}}
\newunicodechar{Ε}{\ensuremath{\varepsilon}}
\newunicodechar{Θ}{\ensuremath{\Theta}}
\newunicodechar{Λ}{\ensuremath{\Lambda}}
\newunicodechar{Μ}{\ensuremath{\mu}}
\newunicodechar{Τ}{\ensuremath{\tau}}
\newunicodechar{Φ}{\ensuremath{\Phi}}
\newunicodechar{Ψ}{\ensuremath{\Psi}}
\newunicodechar{Ω}{\ensuremath{\Omega}}
\newunicodechar{↦}{\ensuremath{\mapsto}}
\newunicodechar{∈}{\ensuremath{\in}}
\newunicodechar{∉}{\ensuremath{\notin}}
\newunicodechar{∧}{\ensuremath{\wedge}}
\newunicodechar{⇔}{\ensuremath{\Leftrightarrow}}
\newunicodechar{⟺}{\ensuremath{\Longleftrightarrow}}
\newunicodechar{⇒}{\ensuremath{\Rightarrow}}
\newunicodechar{⟹}{\ensuremath{\Longrightarrow}}
\newunicodechar{∘}{\ensuremath{\circ}}
\newunicodechar{∪}{\ensuremath{\cup}}
\newunicodechar{∩}{\ensuremath{\cap}}
\newunicodechar{≡}{\ensuremath{\equiv}}
\newunicodechar{⊂}{\ensuremath{\subset}}
\newunicodechar{⊥}{\ensuremath{\perp}}
\newunicodechar{∃}{\ensuremath{\exists}}
\newunicodechar{∀}{\ensuremath{\forall}}
\newunicodechar{∛}{\ensuremath{\sqrt[3]{\,}}}
\newunicodechar{∜}{\ensuremath{\sqrt[4]{\,}}}
\newunicodechar{⃗}{\ensuremath{{}^{\to}}}
\newunicodechar{ⁿ}{\ifmmode^{n}\else\textsuperscript{\ensuremath{n}}\fi}
\newunicodechar{ˣ}{\ifmmode^{x}\else\textsuperscript{\ensuremath{x}}\fi}
\newunicodechar{ᵇ}{\ifmmode^{b}\else\textsuperscript{\ensuremath{b}}\fi}
\newunicodechar{ʳ}{\ifmmode^{r}\else\textsuperscript{\ensuremath{r}}\fi}
\newunicodechar{ᵘ}{\ifmmode^{u}\else\textsuperscript{\ensuremath{u}}\fi}
\newunicodechar{ʸ}{\ifmmode^{y}\else\textsuperscript{\ensuremath{y}}\fi}
\newunicodechar{ᵃ}{\ifmmode^{a}\else\textsuperscript{\ensuremath{a}}\fi}
\newunicodechar{ᵉ}{\ifmmode^{e}\else\textsuperscript{\ensuremath{e}}\fi}
\newunicodechar{ᵏ}{\ifmmode^{k}\else\textsuperscript{\ensuremath{k}}\fi}
\newunicodechar{ᵖ}{\ifmmode^{p}\else\textsuperscript{\ensuremath{p}}\fi}
\newunicodechar{ᴺ}{\ifmmode^{N}\else\textsuperscript{\ensuremath{N}}\fi}
\newunicodechar{ᶜ}{\ifmmode^{c}\else\textsuperscript{\ensuremath{c}}\fi}
\newunicodechar{ʰ}{\ifmmode^{h}\else\textsuperscript{\ensuremath{h}}\fi}
\newunicodechar{ᵠ}{\ifmmode^{q}\else\textsuperscript{\ensuremath{q}}\fi}
\newunicodechar{ₙ}{\ifmmode_{n}\else\textsubscript{\ensuremath{n}}\fi}
\newunicodechar{ₐ}{\ifmmode_{a}\else\textsubscript{\ensuremath{a}}\fi}
\newunicodechar{ᵢ}{\ifmmode_{i}\else\textsubscript{\ensuremath{i}}\fi}
\newunicodechar{ᵣ}{\ifmmode_{r}\else\textsubscript{\ensuremath{r}}\fi}
\newunicodechar{ₖ}{\ifmmode_{k}\else\textsubscript{\ensuremath{k}}\fi}
\newunicodechar{ᵦ}{\ifmmode_{\beta}\else\textsubscript{\ensuremath{\beta}}\fi}
\newunicodechar{ₓ}{\ifmmode_{x}\else\textsubscript{\ensuremath{x}}\fi}
\newfontfamily\popdisplay{ArchivoBlack-Regular.ttf}[Path=../../../fonts/]
\newfontfamily\poplogo{SpaceGrotesk-Bold.ttf}[Path=../../../fonts/]
\newfontfamily\popsemi{PlusJakartaSans-SemiBold.ttf}[Path=../../../fonts/]
\newfontfamily\popxbold{PlusJakartaSans-ExtraBold.ttf}[Path=../../../fonts/]
\newfontfamily\popxboldls{PlusJakartaSans-ExtraBold.ttf}[Path=../../../fonts/,LetterSpace=3.0]

\setlist[enumerate]{leftmargin=1.7em,itemsep=5pt,topsep=4pt}
\setlist[itemize]{leftmargin=1.7em,itemsep=4pt,topsep=4pt,label={\color{poporange}\textbullet}}
\setlength{\parindent}{0pt}
\setlength{\parskip}{5pt}
\renewcommand{\arraystretch}{1.25}

% pgfplots : style Pop par défaut
\pgfplotsset{
  every axis/.append style={axis line style={popink,line width=1pt},tick style={popink},
    grid style={popline,line width=0.5pt},label style={font=\small},tick label style={font=\footnotesize}},
  popcurve/.style={poporange,line width=1.8pt,no markers,smooth},
}
% Même style disponible dans un \draw TikZ simple (hors pgfplots)
\tikzset{popcurve/.style={poporange,line width=1.8pt}}
\tikzset{popgrid/.style={popline,very thin}}
\colorlet{popcurve}{poporange} % le nom du style est parfois utilisé comme couleur

% ============================================================
% Pill / squiggle
% ============================================================
\newcommand{\poppill}[3]{%
  \tikz[baseline=-0.55ex]\node[draw=popink,line width=1.2pt,rounded corners=2.6pt,%
    fill=#2,inner xsep=6.5pt,inner ysep=3pt]{%
    \popxboldls\fontsize{7}{7}\selectfont\color{#3}\MakeUppercase{#1}};%
}
\newcommand{\popsquiggle}[1]{%
  \tikz[baseline=-0.4ex]{
    \draw[#1,line width=2.6pt,line cap=round]
      plot [smooth] coordinates {(0,0.09) (0.34,0.01) (0.68,0.21) (1.02,0.01) (1.36,0.10)};
  }%
}
% Mot-clé surligné (marqueur orange sous le texte)
\newcommand{\cle}[1]{{\popxbold\color{popdark}#1}}

% ============================================================
% En-tête / pied
% ============================================================
\pagestyle{fancy}
\fancyhf{}
\renewcommand{\headrulewidth}{0pt}
\renewcommand{\footrulewidth}{0pt}
\lhead{\raisebox{-0.15ex}{\poplogo\fontsize{13}{13}\selectfont\color{popink}OnBuch\color{poporange}+}}
\rhead{\poppill{\DOCMATIERE\ \textbullet\ \DOCNIVEAU\ \textbullet\ Leçon \DOCLECON}{white}{popink}}
\lfoot{\popxbold\fontsize{7.6}{7.6}\selectfont\color{popmuted}\MakeUppercase{Cours signé OnBuch+}\ \textbullet\ {\popsemi\DOCTITRE}}
\rfoot{\tikz[baseline=-0.6ex]\node[rounded corners=8.5pt,fill=popink,minimum height=17pt,minimum width=17pt,inner xsep=5pt,inner sep=0]{\popxbold\fontsize{8}{8}\selectfont\color{popcream}\thepage\,/\,\pageref*{LastPage}};}

% ============================================================
% Titres
% ============================================================
\newcommand{\chaptitre}[2]{%
  \begin{tcolorbox}[enhanced,colback=poporange,colframe=popink,boxrule=2.6pt,arc=11pt,
      drop shadow=popink,left=15pt,right=15pt,top=12pt,bottom=11pt,before skip=3pt,after skip=18pt]
    \poppill{#2}{popink}{popcream}\par\vspace{9pt}
    {\raggedright\hyphenpenalty=10000\exhyphenpenalty=10000\popdisplay\fontsize{22}{24}\selectfont\color{popcream}\MakeUppercase{#1}\par}
  \end{tcolorbox}
}
% Section numérotée : pastille orange avec le numéro + titre Archivo
\newcounter{popsec}
\newcommand{\coursec}[1]{%
  \stepcounter{popsec}\setcounter{popsub}{0}%
  \par\vspace{16pt}\noindent
  \begin{minipage}{\linewidth}
  \tikz[baseline=-0.7ex]\node[draw=popink,line width=1.6pt,fill=poporange,rounded corners=4pt,minimum size=22pt,inner sep=0]{\popdisplay\fontsize{12}{12}\selectfont\color{popcream}\thepopsec};%
  \hspace{8pt}{\popdisplay\fontsize{15}{17}\selectfont\color{popink}\MakeUppercase{#1}}\par
  \vspace{4pt}\noindent{\color{poporange}\rule{36pt}{3.6pt}}
  \end{minipage}\par\nopagebreak\vspace{8pt}
}
\newcounter{popsub}
\newcommand{\courssub}[1]{%
  \stepcounter{popsub}%
  \par\vspace{9pt}\noindent{\popxbold\fontsize{11.5}{13}\selectfont\color{popink}{\color{poporange}\thepopsec.\thepopsub}\hspace{6pt}#1}\par\nopagebreak\vspace{4pt}
}

% ============================================================
% Boîtes pédagogiques — même recette : fond blanc, bordure épaisse colorée,
% ombre dure encre, badge plein. Titre optionnel [#1].
% ============================================================
\tcbset{popbase/.style={breakable,enhanced,colback=white,boxrule=2.3pt,arc=9pt,
  shadow={3pt}{-3pt}{0pt}{popink},left=14pt,right=14pt,top=11pt,bottom=11pt,before skip=11pt,after skip=15pt,
  fontupper=\popsemi\fontsize{10.6}{14.5}\selectfont\color{popink},
  fontlower=\popsemi\fontsize{10.6}{14.5}\selectfont\color{popink}}}
% Coche dessinée (le glyphe ✓ n'existe pas dans les polices du document)
\newcommand{\popcheck}[1]{\tikz[baseline=-0.5ex]\draw[#1,line width=1.6pt,line cap=round,line join=round](-0.13,0.0)--(-0.04,-0.09)--(0.13,0.1);}
\providecommand{\coche}{\popcheck{popgreen}}
\providecommand{\resultat}[1]{\par\smallskip\noindent\fcolorbox{popgreen}{popgreenL}{\parbox{\dimexpr\linewidth-2\fboxsep-2\fboxrule\relax}{\textbf{Résultat.} #1}}\par\smallskip}
\newcommand{\popboxhead}[3]{\poppill{#1}{#2}{white}\ifx\relax#3\relax\else\hspace{6pt}{\popxbold\fontsize{10.6}{12}\selectfont\color{popink}#3}\fi\par\vspace{7pt}}

\newtcolorbox{aretenir}[1][]{popbase,colframe=popgreen,before upper={\popboxhead{À retenir}{popgreen}{#1}}}
% Texte en allemand (document de lecture, dialogue, extrait) : cadre
% d'étude. Un titre optionnel (source, auteur) s'affiche dans l'en-tête.
\newtcolorbox{textebox}[1][]{popbase,colframe=popblue,before upper={\popboxhead{Text}{popblue}{#1}\begin{otherlanguage}{ngerman}},after upper={\end{otherlanguage}}}
% Marques ✓ / ✗ (forme correcte / fautive) souvent utilisées par les modèles
\providecommand{\cross}{\textcolor{poppink}{\ensuremath{\times}}}
\providecommand{\xmark}{\cross}\providecommand{\crossmark}{\cross}\providecommand{\cmark}{\ensuremath{\checkmark}}
% Ligne de réponse / justification à compléter (inventée par les modèles)
\providecommand{\justification}[1]{\par\noindent\textit{Justification~:} \dotfill\par}
% \textipa (package tipa non chargé) : la police ne couvre pas l'API -> simple passage
\providecommand{\textipa}[1]{#1}
% Soulignement (ulem non chargé) : \uline, \ul
\providecommand{\uline}[1]{\underline{#1}}\providecommand{\ul}[1]{\underline{#1}}
% \alert (beamer) inventée par les modèles : texte mis en évidence
\providecommand{\alert}[1]{\textbf{\textcolor{poppink}{#1}}}
% variantes ASCII d'ordinaux écrites en macro
\providecommand{\iere}{\textsuperscript{re}}\providecommand{\ieme}{\textsuperscript{e}}\providecommand{\ere}{\textsuperscript{re}}\providecommand{\eme}{\textsuperscript{e}}\providecommand{\er}{\textsuperscript{er}}\providecommand{\ie}{\textsuperscript{e}}
% Ordinaux français écrits en macro par les modèles : 1\ière, 3\ième
\newcommand{\ière}{\textsuperscript{re}}\newcommand{\ième}{\textsuperscript{e}}
% \exemple (inventée par les modèles) : étiquette « Exemple : »
\providecommand{\exemple}{\textbf{Exemple~:}\ }
% \square utilisable aussi hors mode math (cases à cocher)
\AtBeginDocument{\let\squareMath\square\renewcommand{\square}{\ensuremath{\squareMath}}}
% Mot ou phrase en allemand à l'intérieur d'un paragraphe français
\newcommand{\all}[1]{\ifmmode\text{\textit{#1}}\else\begingroup\selectlanguage{ngerman}\itshape #1\endgroup\fi}
\let\esp\all % alias toléré
\newtcolorbox{exemplebox}[1][]{popbase,colframe=poporange,before upper={\popboxhead{Exemple}{poporange}{#1}}}
\newtcolorbox{definition}[1][]{popbase,colframe=popblue,before upper={\popboxhead{Définition}{popblue}{#1}}}
\newtcolorbox{propriete}[1][]{popbase,colframe=poppurple,before upper={\popboxhead{Propriété}{poppurple}{#1}}}
\newtcolorbox{methode}[1][]{popbase,colframe=popgold,before upper={\popboxhead{Méthode}{popgold}{#1}}}
\newtcolorbox{attention}[1][]{popbase,colframe=poppink,before upper={\popboxhead{Attention}{poppink}{#1}}}
\newtcolorbox{experience}[1][]{popbase,colframe=popblue,colback=popblueL,before upper={\popboxhead{Expérience}{popblue}{#1}}}
% « Le savais-tu ? » : carte violette pleine (contexte, culture, Cameroun)
\newtcolorbox{savaistu}[1][]{popbase,colback=poppurple,colframe=popink,
  fontupper=\popsemi\fontsize{10.4}{14.2}\selectfont\color{white},
  before upper={\poppill{Le savais-tu\,?}{white}{poppurple}\ifx\relax#1\relax\else\hspace{6pt}{\popxbold\color{white}#1}\fi\par\vspace{7pt}}}
% Exercice résolu : énoncé en haut, \tcblower puis la solution sur fond crème
\newcounter{popexo}\newcounter{popexr}
\newtcolorbox{exoresolu}[1][]{popbase,colframe=poporange,colbacklower=popcream,
  segmentation style={popink,line width=1.2pt,dashed},
  before upper={\stepcounter{popexr}\popboxhead{Exercice résolu \thepopexr}{poporange}{#1}},
  before lower={\poppill{Solution}{popink}{popcream}\par\vspace{6pt}}}
% Exercice (non résolu en place) — la correction est dans la partie Corrigés
\newtcolorbox{exercice}[1][]{popbase,colframe=popink,
  before upper={\stepcounter{popexo}\popboxhead{Exercice \thepopexo}{popink}{#1}}}
% Corrigé
\newtcolorbox{corrige}[1][]{popbase,colframe=popgreen,colback=popgreenL,
  before upper={\popboxhead{Corrigé}{popgreen}{#1}}}
% Objectifs / prérequis / bilan
\newtcolorbox{objectifs}{popbase,colframe=popink,colback=poporangeL,
  before upper={\popboxhead{Objectifs de la leçon}{popink}{}}}
\newtcolorbox{prerequis}{popbase,colframe=popink,colback=popblueL,
  before upper={\popboxhead{Prérequis}{popblue}{}}}
\newtcolorbox{bilan}{popbase,colframe=popink,colback=white,boxrule=2.8pt,
  before upper={\popboxhead{Fiche bilan}{poporange}{L'essentiel à connaître pour l'examen}}}
% Figure encadrée avec légende
\newtcolorbox{popfigure}[1][]{enhanced,breakable=false,colback=white,colframe=popline,boxrule=1.4pt,arc=8pt,
  left=10pt,right=10pt,top=10pt,bottom=8pt,before skip=10pt,after skip=12pt,halign=center,
  lower separated=false,halign lower=center,
  fontlower=\popsemi\fontsize{9}{11.5}\selectfont\color{popmuted},
  #1}
\newcommand{\legende}[1]{\par\vspace{6pt}{\popsemi\fontsize{9}{11.5}\selectfont\color{popmuted}#1}}

% ============================================================
% Signature OnBuch+
% ============================================================
\newcommand{\poplogoinline}[1]{{\poplogo\fontsize{#1}{#1}\selectfont\color{popink}OnBuch\color{poporange}+}}
\newcommand{\signatureonbuch}{%
  \begin{tcolorbox}[enhanced,colback=popink,colframe=popink,boxrule=2.2pt,arc=10pt,
      drop shadow=poporange,left=14pt,right=14pt,top=10pt,bottom=10pt,before skip=0pt,after skip=0pt]
    \tikz[baseline=-0.7ex]\node[circle,fill=popgreen,draw=popcream,line width=1.3pt,minimum size=22pt,inner sep=0]{\popcheck{white}};%
    \hspace{9pt}\begin{minipage}[c]{0.62\linewidth}
      {\popxbold\fontsize{12}{13}\selectfont\color{popcream}Signé\ \textbullet\ {\poplogo OnBuch\color{poporange}+}}\par\vspace{2pt}
      {\raggedright\popsemi\fontsize{8}{10.5}\selectfont\color{popline}\MakeUppercase{Équipe pédagogique OnBuch+}\\\MakeUppercase{Conforme au programme officiel MINESEC}\par}
    \end{minipage}\hfill
    {\poplogo\fontsize{22}{22}\selectfont\color{popcream}OnBuch\color{poporange}+}
  \end{tcolorbox}%
}

% ============================================================
% Page de garde
% #1 titre · #2 matière · #3 niveau · #4 module · #5 n° leçon · #6 durée ·
% #7 description / objectifs (texte ou itemize)
% ============================================================
\newcommand{\pagegarde}[7]{%
  \thispagestyle{empty}
  \newgeometry{a4paper,margin=1.7cm,top=1.6cm,bottom=1.4cm}%
  \begin{tikzpicture}[remember picture,overlay]
    \fill[poporange!18] ([xshift=-3.2cm,yshift=-3.6cm]current page.north east) circle (4.6cm);
    \fill[poppurple!14] ([xshift=2.4cm,yshift=7.6cm]current page.south west) circle (3.8cm);
  \end{tikzpicture}%
  {\poplogo\fontsize{30}{30}\selectfont\color{popink}OnBuch\color{poporange}+}\hfill
  \raisebox{0.6ex}{\poppill{Cours signé \textbullet\ Premium}{popink}{popcream}}\par
  \vspace{1pt}\popsquiggle{poppurple}\par
  \vspace{8pt}
  \begin{tcolorbox}[enhanced,colback=poporange,colframe=popink,boxrule=2.8pt,arc=13pt,
      drop shadow=popink,left=18pt,right=18pt,top=12pt,bottom=13pt,after skip=10pt]
    \poppill{#2 \textbullet\ #3}{popink}{popcream}\hspace{5pt}\poppill{#4}{white}{popink}\par\vspace{8pt}
    {\popdisplay\fontsize{14}{14}\selectfont\color{popink}LEÇON #5}\par\vspace{4pt}
    {\raggedright\hyphenpenalty=10000\exhyphenpenalty=10000\popdisplay\fontsize{25}{27}\selectfont\color{popcream}\MakeUppercase{#1}\par}
    \vspace{8pt}
    {\popxbold\fontsize{9.5}{9.5}\selectfont\color{popink}\MakeUppercase{Durée conseillée : #6}}
  \end{tcolorbox}
  \begin{tcolorbox}[enhanced,colback=white,colframe=popink,boxrule=2.2pt,arc=10pt,
      drop shadow=popink,left=16pt,right=16pt,top=10pt,bottom=10pt,after skip=8pt]
    \poppill{Dans cette leçon}{poppurple}{white}\par\vspace{5pt}
    \popsemi\fontsize{9.3}{12.6}\selectfont\color{popink}#7
  \end{tcolorbox}
  \noindent\begin{minipage}[t]{0.487\linewidth}
    \begin{tcolorbox}[enhanced,colback=popgreen,colframe=popink,boxrule=2.2pt,arc=10pt,
        drop shadow=popink,left=11pt,right=11pt,top=10pt,bottom=10pt,height=3.1cm,valign=center]
      \tcbox[colback=white,colframe=popink,boxrule=1.3pt,arc=5pt,boxsep=0pt,nobeforeafter,
        left=3pt,right=3pt,top=3pt,bottom=3pt,valign=center]{\qrcode[height=1.9cm]{https://play.google.com/store/apps/details?id=com.luvvix.onbuch&hl=fr}}%
      \hspace{8pt}\begin{minipage}[c]{0.5\linewidth}
        {\popdisplay\fontsize{11}{12.5}\selectfont\color{white}\MakeUppercase{Télécharge OnBuch}}\par\vspace{4pt}
        {\raggedright\popsemi\fontsize{8.4}{11}\selectfont\color{white}Gratuit \textbullet\ Google Play\\Tuteur IA, annales, concours, résultats\par}
      \end{minipage}
    \end{tcolorbox}
  \end{minipage}\hfill
  \begin{minipage}[t]{0.487\linewidth}
    \begin{tcolorbox}[enhanced,colback=poppurple,colframe=popink,boxrule=2.2pt,arc=10pt,
        drop shadow=popink,left=11pt,right=11pt,top=10pt,bottom=10pt,height=3.1cm,valign=center]
      \tikz[baseline=-0.7ex]\node[circle,fill=white,draw=popink,line width=1.3pt,minimum size=1.6cm,inner sep=0]{\popdisplay\fontsize{24}{24}\selectfont\color{poppurple}+};%
      \hspace{8pt}\begin{minipage}[c]{0.6\linewidth}
        {\popdisplay\fontsize{11}{12.5}\selectfont\color{white}\MakeUppercase{Passe à OnBuch+}}\par\vspace{4pt}
        {\raggedright\popsemi\fontsize{8.4}{11}\selectfont\color{white}Tous les cours signés, Tuteur IA illimité, fiches PDF — onglet OnBuch+ de l'app\par}
      \end{minipage}
    \end{tcolorbox}
  \end{minipage}
  \vspace{8pt}
  \popcontenu
  \vfill
  \signatureonbuch
  \restoregeometry
  \newpage
}

% Rangée « Ce cours contient » (page de garde)
\newcommand{\poptile}[3]{% #1 couleur #2 gros texte #3 légende
  \begin{tcolorbox}[enhanced,colback=#1,colframe=popink,boxrule=2pt,arc=9pt,shadow={2.5pt}{-2.5pt}{0pt}{popink},
      left=6pt,right=6pt,top=9pt,bottom=9pt,halign=center,valign=center,height=1.75cm,width=\linewidth,nobeforeafter]
    {\popdisplay\fontsize{12.5}{13}\selectfont\color{white}\MakeUppercase{#2}}\par\vspace{3pt}
    {\centering\popsemi\fontsize{7.8}{9.5}\selectfont\color{white}#3\par}
  \end{tcolorbox}}
\newcommand{\popcontenu}{%
  \par\noindent{\popxboldls\fontsize{7.5}{7.5}\selectfont\color{popmuted}CE COURS SIGNÉ CONTIENT}\par\vspace{7pt}
  \noindent\begin{minipage}[t]{0.235\linewidth}\poptile{popblue}{Cours}{complet, illustré, pas à pas}\end{minipage}\hfill
  \begin{minipage}[t]{0.235\linewidth}\poptile{poporange}{Méthodes}{et exercices résolus}\end{minipage}\hfill
  \begin{minipage}[t]{0.235\linewidth}\poptile{poppink}{Exercices}{type examen + corrigés}\end{minipage}\hfill
  \begin{minipage}[t]{0.235\linewidth}\poptile{popgreen}{Fiche bilan}{l'essentiel à retenir}\end{minipage}\par}

% Sommaire
\newtcolorbox{sommaire}{enhanced,colback=white,colframe=popink,boxrule=2.2pt,arc=10pt,
  drop shadow=popink,left=18pt,right=18pt,top=13pt,bottom=13pt,before skip=6pt,after skip=20pt,
  before upper={\poppill{Sommaire}{poppurple}{white}\par\vspace{9pt}\popsemi\fontsize{10.6}{14.5}\selectfont\color{popink}}}

% Signature de fin de cours — poussée tout en bas de la dernière page
\newcommand{\findecours}{%
  \par\vfill\nopagebreak
  \begin{center}\popsquiggle{poporange}\end{center}\vspace{4pt}
  \signatureonbuch
}
\AtBeginDocument{\def\llbracket{\mathopen{[\mkern-3.5mu[}}\def\rrbracket{\mathclose{]\mkern-3.5mu]}}} % intervalles d'entiers (glyphe ⟦ absent de la police)
% Glyphes absents de Fira Math : redéfinis à partir de glyphes disponibles
\AtBeginDocument{%
  \def\setminus{\mathbin{\backslash}}\let\smallsetminus\setminus
  \def\vdots{\mathord{\vbox{\baselineskip3.2pt\lineskiplimit0pt\kern4pt\hbox{.}\hbox{.}\hbox{.}}}}%
  \def\ddots{\mathinner{\mkern1mu\raise6.5pt\hbox{.}\mkern2mu\raise3.6pt\hbox{.}\mkern2mu\raise0.7pt\hbox{.}\mkern1mu}}%
  \def\triangle{\mathord{\scalebox{0.9}{$\symup{\Delta}$}}}%
  \def\nabla{\mathord{\raisebox{\depth}{\scalebox{1}[-1]{$\symup{\Delta}$}}}}%
  \def\checkmark{\popcheck{popdark}}%
  \def\star{\mathbin{\ast}}\def\bigstar{\mathord{\ast}}%
  \def\diamond{\mathbin{\scalebox{0.6}{$\Diamond$}}}%
  \def\bigcup{\mathop{\mathchoice{\raisebox{-0.3em}{\scalebox{1.7}{$\cup$}}}{\scalebox{1.25}{$\cup$}}{\cup}{\cup}}\displaylimits}%
  \def\bigcap{\mathop{\mathchoice{\raisebox{-0.3em}{\scalebox{1.7}{$\cap$}}}{\scalebox{1.25}{$\cap$}}{\cap}{\cap}}\displaylimits}%
  \def\lesseqqgtr{\mathrel{\lessgtr}}\def\bigoplus{\mathop{\oplus}\displaylimits}%
}
\providecommand{\ssi}{si et seulement si} % abréviation fréquente
\AtBeginDocument{\providecommand{\wideparen}{\overparen}} % arc de cercle

\begin{document}

\pagegarde{Les types de familles}{Allemand}{1ère A}{Chapitre 1 : Vie familiale et sociale}{ALA01}{2 h}{Leçon de type « texte » consacrée à la lecture et au commentaire d'un texte sur les types de familles dans les pays germanophones. Elle permet d'acquérir le vocabulaire de la famille, de réviser l'ordre des mots, les articles, les possessifs et la conjugaison au présent, puis de produire un court texte sur sa propre famille.
\vspace{4pt}
\begin{multicols}{2}\small
\begin{itemize}
\item À la fin de cette leçon, je sais comprendre globalement puis en détail un texte allemand sur les types de familles.
\item À la fin de cette leçon, je sais employer le vocabulaire thématique de la famille (die Familie, die Eltern, die Großeltern, die Geschwister, verheiratet, allein erziehend) avec les articles et les pluriels corrects.
\item À la fin de cette leçon, je sais construire une phrase simple allemande en respectant la place du verbe en 2e position.
\item À la fin de cette leçon, je sais utiliser les articles définis, indéfinis et les adjectifs possessifs au nominatif et à l'accusatif.
\item À la fin de cette leçon, je sais conjuguer au présent les verbes faibles, les verbes forts, sein et haben.
\item À la fin de cette leçon, je sais répondre par des phrases complètes à des questions de compréhension sur un texte.
\end{itemize}\end{multicols}}

\begin{sommaire}
\begin{multicols}{2}
\begin{enumerate}
\item Texte support : Familien heute
\item Compréhension du texte
\item Lexique thématique : die Familie
\item Grammaire 1 : l'ordre des mots et les articles/possessifs
\item Grammaire 2 : conjugaison au présent
\item Commentaire et production guidée
\item Activité d'intégration
\item Méthodes et exercices résolus
\item Exercices
\item Corrigés des exercices
\item Fiche bilan
\end{enumerate}
\end{multicols}
\end{sommaire}

\begin{objectifs}
\begin{itemize}
\item À la fin de cette leçon, je sais comprendre globalement puis en détail un texte allemand sur les types de familles.
\item À la fin de cette leçon, je sais employer le vocabulaire thématique de la famille (die Familie, die Eltern, die Großeltern, die Geschwister, verheiratet, allein erziehend) avec les articles et les pluriels corrects.
\item À la fin de cette leçon, je sais construire une phrase simple allemande en respectant la place du verbe en 2e position.
\item À la fin de cette leçon, je sais utiliser les articles définis, indéfinis et les adjectifs possessifs au nominatif et à l'accusatif.
\item À la fin de cette leçon, je sais conjuguer au présent les verbes faibles, les verbes forts, sein et haben.
\item À la fin de cette leçon, je sais répondre par des phrases complètes à des questions de compréhension sur un texte.
\item À la fin de cette leçon, je sais rédiger un court paragraphe (5 à 8 phrases) pour présenter ma famille.
\item À la fin de cette leçon, je sais comparer les réalités familiales du Cameroun et des pays germanophones.
\end{itemize}
\end{objectifs}

\begin{prerequis}
\begin{itemize}
\item L'alphabet, les salutations et la présentation personnelle (Seconde)
\item Les articles définis der/die/das et indéfinis ein/eine au nominatif (Seconde)
\item La conjugaison de sein et haben au présent (Seconde)
\item La phrase affirmative simple : sujet + verbe conjugué (Seconde)
\item Les chiffres et les nombres pour parler des membres de la famille (Seconde)
\end{itemize}
\end{prerequis}

\newpage

\coursec{Texte support : Familien heute}

\courssub{Situation d'accroche et problématique}

Imaginez la scène : à Yaoundé, chez les Mbarga, on vit sous le même toit avec le père, la mère, les enfants, mais aussi l'oncle, la tante et les cousins venus étudier en ville. Au Cameroun, on connaît bien la famille nucléaire et la famille élargie. Mais comment vit-on en famille en Allemagne, en Autriche ou en Suisse ? Là-bas aussi, les modèles familiaux changent : familles monoparentales, familles recomposées, couples sans enfants.

\textbf{Problématique :} Quels types de familles existe-t-il dans les pays germanophones, et comment en parler en allemand ?

Dans cette leçon, nous lisons un texte sur les types de familles et nous apprenons à parler de notre propre famille.

\courssub{Lecture globale du texte}

\begin{methode}[Bien lire un texte allemand en 5 étapes]
\begin{enumerate}
\item Lire le \cle{titre} et anticiper : de quoi va parler le texte ?
\item Faire une première lecture \cle{globale}, sans dictionnaire, pour saisir l'idée générale.
\item Repérer les \cle{mots transparents} (ressemblant au français) : \all{Familie}, \all{Mutter}, \all{Partner}, \all{klassisch}.
\item Faire une deuxième lecture \cle{détaillée} avec le glossaire : souligner le vocabulaire familial et les verbes conjugués.
\item Résumer le texte en une phrase, en français puis en allemand.
\end{enumerate}
\end{methode}

\begin{textebox}[Familien heute in Deutschland]
In Deutschland gibt es heute viele verschiedene Familienformen. Die Kernfamilie – Vater, Mutter und Kinder – ist klassisch. Aber viele Kinder leben bei einer allein erziehenden Mutter oder einem allein erziehenden Vater. Es gibt auch Patchworkfamilien: Die Eltern sind geschieden und haben neue Partner. Die Kinder haben dann Stiefbrüder und Stiefschwestern. Manche Paare sind verheiratet und haben keine Kinder; sie sind ein kinderloses Ehepaar. Die Großeltern wohnen oft nicht in der Familie, aber sie besuchen ihre Enkel gern. Auch in Österreich und in der Schweiz sieht man diese Familienformen. Die Familie ist heute bunt und vielfältig – und das ist gut so!
\end{textebox}

\textbf{Glossaire :}

\begin{center}
\begin{tabularx}{\linewidth}{|X|X|}\hline
\rowcolor{popblueL} \textbf{Deutsch} & \textbf{Français} \\ \hline
die Familie (-n) & la famille \\ \hline
die Eltern (Pl.) & les parents \\ \hline
die Großeltern (Pl.) & les grands-parents \\ \hline
die Geschwister (Pl.) & les frères et sœurs \\ \hline
das Kind (-er) & l'enfant \\ \hline
die Tochter (¨-) & la fille \\ \hline
der Sohn (¨-e) & le fils \\ \hline
der Enkel (-) / die Enkelin (-nen) & le petit-fils / la petite-fille \\ \hline
verheiratet & marié(e) \\ \hline
geschieden & divorcé(e) \\ \hline
allein erziehend & qui élève seul(e) \\ \hline
das Ehepaar (-e) & le couple (marié) \\ \hline
der Stiefbruder (¨-) & le demi-frère (par remariage) \\ \hline
leben & vivre \\ \hline
zusammenleben & vivre ensemble \\ \hline
besuchen & rendre visite à \\ \hline
vielfältig & divers, varié \\ \hline
\end{tabularx}
\end{center}

\textbf{Première compréhension globale :} Le texte est un petit \cle{article descriptif} (type : texte informatif). Thème : les formes de familles aujourd'hui en Allemagne. Personnages évoqués : les parents, les enfants, les grands-parents, les couples sans enfants.

\courssub{Les types de familles dans le texte}

\begin{definition}[Les cinq types de familles]
\begin{itemize}
\item \cle{die Kernfamilie} : la famille nucléaire — père, mère et enfants.
\item \cle{die Großfamilie} : la famille élargie — avec grands-parents, oncles, tantes, cousins (très courante au Cameroun).
\item \cle{die allein erziehende Mutter / der allein erziehende Vater} : la mère / le père qui élève seul(e) son enfant (famille monoparentale).
\item \cle{die Patchworkfamilie} : la famille recomposée après un divorce ou une séparation, avec de nouveaux partenaires et leurs enfants.
\item \cle{das kinderlose Ehepaar} : le couple marié sans enfants.
\end{itemize}
\end{definition}

\begin{popfigure}
\begin{center}
\begin{tikzpicture}[mindmap, grow cyclic, every node/.style={concept, align=center, font=\small}, concept color=popblue!60, level 1 concept/.append style={sibling angle=72, level distance=3.4cm}]
\node[concept color=popblue!70] {\all{die Familie}}
  child[concept color=popgreen!60] { node[align=center]{\all{die}\\\all{Kernfamilie}} }
  child[concept color=poporange!60] { node[align=center]{\all{die}\\\all{Großfamilie}} }
  child[concept color=poppink!60] { node[align=center]{\all{allein}\\\all{erziehend}} }
  child[concept color=poppurple!60] { node[align=center]{\all{die}\\\all{Patchworkfamilie}} }
  child[concept color=popgold!70] { node[align=center]{\all{das kinderlose}\\\all{Ehepaar}} };
\end{tikzpicture}
\end{center}
\legende{Carte mentale : les cinq types de familles}
\end{popfigure}

\courssub{Lecture détaillée : observation du vocabulaire et des verbes}

Lisons maintenant phrase par phrase. Soulignons les mots de la famille et encadrons les verbes conjugués :

\begin{itemize}
\item \all{In Deutschland \textbf{gibt} es heute viele verschiedene Familienformen.} « Il existe aujourd'hui en Allemagne beaucoup de formes de familles différentes. » — \all{es gibt} = « il y a ».
\item \all{Die Kernfamilie \textbf{ist} klassisch.} « La famille nucléaire est classique. » — verbe \all{sein} (être).
\item \all{Viele Kinder \textbf{leben} bei einer allein erziehenden Mutter.} « Beaucoup d'enfants vivent chez une mère qui élève seule. » — verbe \all{leben}.
\item \all{Die Eltern \textbf{sind} geschieden und \textbf{haben} neue Partner.} « Les parents sont divorcés et ont de nouveaux partenaires. » — verbes \all{sein} et \all{haben}.
\item \all{Die Großeltern \textbf{wohnen} oft nicht in der Familie, aber sie \textbf{besuchen} ihre Enkel gern.} « Les grands-parents n'habitent souvent pas dans la famille, mais ils rendent volontiers visite à leurs petits-enfants. »
\end{itemize}

\textbf{Repérage de 5 phrases simples (ordre des mots) :} observez que dans chaque phrase, le verbe conjugué est à la \cle{deuxième place} :

\begin{center}
\begin{tabularx}{\linewidth}{|X|X|}\hline
\rowcolor{popblueL} \textbf{Deutsch} & \textbf{Français} \\ \hline
Die Kernfamilie ist klassisch. & La famille nucléaire est classique. \\ \hline
Viele Kinder leben bei der Mutter. & Beaucoup d'enfants vivent chez la mère. \\ \hline
Die Eltern haben neue Partner. & Les parents ont de nouveaux partenaires. \\ \hline
Die Großeltern besuchen ihre Enkel gern. & Les grands-parents visitent volontiers leurs petits-enfants. \\ \hline
Manche Paare haben keine Kinder. & Certains couples n'ont pas d'enfants. \\ \hline
\end{tabularx}
\end{center}

\courssub{Écoute et relecture}

\begin{experience}[Compréhension orale et lecture à voix haute]
\begin{enumerate}
\item \textbf{Écoute :} le professeur lit le texte une première fois, livre fermé. Notez trois mots que vous reconnaissez.
\item \textbf{Écoute active :} deuxième écoute, livre ouvert. Levez la main quand vous entendez un type de famille (\all{Kernfamilie}, \all{Patchworkfamilie}...).
\item \textbf{Relecture individuelle :} chaque élève lit une phrase à voix haute. Attention à la prononciation : \all{Großeltern} se prononce « gros-èl-tèrne », \all{geschieden} « gue-chi-dène », \all{Ehepaar} « é-o-par ».
\end{enumerate}
\end{experience}

\begin{exoresolu}[Questions de compréhension]
Énoncé : Répondez en allemand par une phrase complète. 1) \all{Was ist die Kernfamilie?} 2) \all{Wo wohnen die Großeltern oft?} 3) \all{Was ist eine Patchworkfamilie?}
\tcblower
Solution détaillée : on reprend les mots de la question et on place le verbe en deuxième position.
\begin{enumerate}
\item \all{Die Kernfamilie ist Vater, Mutter und Kinder.} « La famille nucléaire, c'est le père, la mère et les enfants. »
\item \all{Die Großeltern wohnen oft nicht in der Familie.} « Les grands-parents n'habitent souvent pas dans la famille. »
\item \all{In einer Patchworkfamilie sind die Eltern geschieden und haben neue Partner.} « Dans une famille recomposée, les parents sont divorcés et ont de nouveaux partenaires. »
\end{enumerate}
\end{exoresolu}

\begin{attention}[Pièges pour les francophones]
\begin{itemize}
\item \all{die Eltern} est \textbf{toujours pluriel} et signifie « les parents » (père et mère). Ne pas confondre avec \all{die Verwandten} : « les parents » au sens de la parentèle (oncles, cousins...).
\item \all{das Kind} est neutre : on dit \all{das Kind}, jamais « der Kind ».
\item Majuscule à tous les noms : \all{die Mutter}, \all{der Sohn}, \all{die Familie}.
\item \all{verheiratet} se prononce « fè-r-haï-ra-tète » et s'emploie avec \all{sein} : \all{Sie sind verheiratet.} « Ils sont mariés. »
\end{itemize}
\end{attention}

\begin{savaistu}[Familles en Allemagne]
En Allemagne, environ un enfant sur cinq vit avec un parent seul. Le mot \all{Patchworkfamilie} vient de la couture en patchwork : on assemble des morceaux de tissu différents pour créer quelque chose de nouveau — comme une famille recomposée ! Et au Cameroun ? La famille élargie (\all{die Großfamilie}) reste le modèle le plus répandu : oncles, tantes et cousins font partie de la vie quotidienne.
\end{savaistu}

\begin{exercice}[À vous !]
1) Soulignez dans le texte tous les mots de la famille. 2) Relevez cinq verbes conjugués et donnez leur infinitif. 3) Résumez le texte en une phrase allemande en commençant par \all{In Deutschland gibt es ...}
\end{exercice}

\begin{corrige}[Exercice « À vous ! »]
1) \all{Familienformen, Vater, Mutter, Kinder, Eltern, Partner, Stiefbrüder, Stiefschwestern, Ehepaar, Großeltern, Enkel}. 2) \all{gibt} (geben), \all{ist} (sein), \all{leben} (leben), \all{haben} (haben), \all{wohnen} (wohnen), \all{besuchen} (besuchen). 3) Exemple : \all{In Deutschland gibt es viele verschiedene Familienformen.} « En Allemagne, il existe beaucoup de formes de familles différentes. »
\end{corrige}

\coursec{Compréhension du texte}

Dans la section précédente, nous avons lu et écouté le texte support \emph{Familien heute}, qui présente les différentes formes de familles en Allemagne : la famille nucléaire, la famille monoparentale, la famille recomposée et la famille élargie avec les grands-parents. Nous avons dégagé une première idée générale du texte. Il s'agit maintenant de vérifier et d'approfondir cette compréhension, étape par étape : d'abord la compréhension globale, puis la compréhension détaillée, enfin le résumé. C'est exactement la démarche exigée à l'examen.

\courssub{La méthode de compréhension : du global au détaillé}

Comprendre un texte allemand ne signifie pas traduire chaque mot. Un bon lecteur procède en deux lectures : une \cle{lecture globale} pour saisir le thème général, puis une \cle{lecture détaillée} pour repérer les informations précises demandées par les questions.

\begin{popfigure}
\begin{tikzpicture}[
  node distance=0.55cm,
  box/.style={draw=popblue, fill=popblueL, rounded corners=3pt, align=center, text width=4.2cm, minimum height=0.9cm, font=\small},
  arr/.style={-{Stealth[length=3mm]}, thick, poporange}
]
\node[box, align=center] (a){\textbf{1. Lecture globale}\\ lire tout le texte une fois};
\node[box, below=of a, align=center] (b){\textbf{2. Questions globales}\\ thème, personnages, lieu};
\node[box, below=of b, align=center] (c){\textbf{3. Lecture détaillée}\\ souligner les mots-clés};
\node[box, below=of c, align=center] (d){\textbf{4. Questions détaillées}\\ répondre en phrases complètes};
\node[box, below=of d, align=center] (e){\textbf{5. Résumé}\\ 2 à 3 phrases avec \all{und}, \all{aber}, \all{auch}};
\draw[arr] (a) -- (b);
\draw[arr] (b) -- (c);
\draw[arr] (c) -- (d);
\draw[arr] (d) -- (e);
\end{tikzpicture}
\legende{Organigramme de la méthode de compréhension d'un texte}
\end{popfigure}

\begin{methode}[Répondre à une question de compréhension]
Pour répondre correctement à une question de compréhension en allemand, suivez toujours ces quatre étapes :
\begin{enumerate}
\item \textbf{Lire la question} et repérer le mot interrogatif (\all{wer}, \all{was}, \all{wo}...) : il annonce le type d'information attendu.
\item \textbf{Retrouver dans le texte} la phrase qui contient la réponse et la souligner.
\item \textbf{Reprendre les mots de la question} pour construire la réponse : la question \all{Wo wohnen die Großeltern ?} donne le début de la réponse \all{Die Großeltern wohnen...}.
\item \textbf{Rédiger une phrase complète}, jamais un seul mot. On écrit \all{Die Großeltern wohnen oft nicht in der Familie.} et non simplement \all{Nicht in der Familie.}
\end{enumerate}
\end{methode}

\begin{attention}[Majuscule et place du verbe]
Deux règles absolues pour toute réponse rédigée en allemand :
\begin{itemize}
\item Toute phrase commence par une \cle{majuscule}, et tous les noms communs prennent aussi une majuscule : \all{die Mutter}, \all{das Kind}, \all{die Großeltern}.
\item Le \cle{verbe conjugué} est toujours en \cle{2e position}, même si la phrase commence par un complément. On dit \all{In Yaoundé wohnen viele Großeltern bei der Familie.} (« À Yaoundé, beaucoup de grands-parents habitent avec la famille. ») et jamais \all{In Yaoundé viele Großeltern wohnen...}. Le verbe ne recule jamais à la troisième place.
\end{itemize}
\end{attention}

\courssub{Les mots interrogatifs : la clé des questions}

Avant de répondre, il faut comprendre la question. Observons les questions posées sur notre texte : \all{Wer besucht die Enkel gern ?} (Qui visite volontiers les petits-enfants ?), \all{Wo wohnen die Großeltern oft ?} (Où habitent souvent les grands-parents ?), \all{Was ist eine Patchworkfamilie ?} (Qu'est-ce qu'une famille recomposée ?). Chaque question commence par un \cle{mot interrogatif}, suivi immédiatement du verbe conjugué.

\begin{propriete}[Les mots interrogatifs utiles]
Les mots interrogatifs allemands se placent en \cle{1re position} et le verbe conjugué en \cle{2e position} : \all{Wo wohnen die Großeltern ?}
\end{propriete}

\begin{center}
\begin{tabularx}{\linewidth}{|l|X|X|}
\hline
\rowcolor{popblueL} \textbf{Deutsch} & \textbf{Français} & \textbf{Exemple} \\
\hline
\all{wer} & qui (personne) & \all{Wer besucht die Enkel ?} — Qui visite les petits-enfants ? \\
\hline
\all{was} & que, quoi (chose) & \all{Was ist eine Patchworkfamilie ?} — Qu'est-ce qu'une famille recomposée ? \\
\hline
\all{wo} & où (lieu) & \all{Wo wohnen die Großeltern ?} — Où habitent les grands-parents ? \\
\hline
\all{wie} & comment & \all{Wie leben viele Kinder ?} — Comment vivent beaucoup d'enfants ? \\
\hline
\all{warum} & pourquoi & \all{Warum ist die Mutter allein ?} — Pourquoi la mère est-elle seule ? \\
\hline
\all{wie viele} & combien de & \all{Wie viele Kinder hat die Familie ?} — Combien d'enfants la famille a-t-elle ? \\
\hline
\end{tabularx}
\end{center}

\begin{attention}[Faux amis des mots interrogatifs]
Ne confondez pas \all{wer} (qui) avec \all{wo} (où) : les francophones les inversent souvent. De même, \all{wie viele} signifie « combien de » et non « comment ». Et attention : \all{was} interroge sur une chose, jamais sur une personne — pour une personne, c'est toujours \all{wer}.
\end{attention}

\courssub{Compréhension globale : Worum geht es im Text ?}

Les questions globales portent sur le texte entier. Les deux questions classiques sont : \all{Worum geht es im Text ?} (« De quoi parle le texte ? ») et \all{Welche Familienformen gibt es ?} (« Quelles formes de familles existe-t-il ? »).

Pour répondre à \all{Worum geht es im Text ?}, on utilise la formule : \all{Der Text handelt von...} (« Le texte parle de... ») ou \all{Im Text geht es um...} (« Il s'agit dans le texte de... »). La préposition \all{von} et \all{um} gouvernent le datif (ici \all{den Familien}, \all{verschiedenen Familienformen}).

\begin{exemplebox}[Réponses globales rédigées]
\all{Der Text handelt von den Familien in Deutschland heute.} « Le texte parle des familles en Allemagne aujourd'hui. »

\all{Im Text geht es um verschiedene Familienformen.} « Dans le texte, il s'agit de différentes formes de familles. »

\all{Es gibt die Kernfamilie, die allein erziehende Mutter und die Patchworkfamilie.} « Il y a la famille nucléaire, la mère qui élève seule et la famille recomposée. »
\end{exemplebox}

\courssub{Compréhension détaillée : les questions précises}

Les questions détaillées portent sur une information précise du texte. Reprenons les trois questions du programme et construisons les réponses en reprenant les mots de la question.

\textbf{Question 1 :} \all{Wer besucht die Enkel gern ?} (Qui visite volontiers les petits-enfants ?)

Réponse complète : \all{Die Großeltern besuchen die Enkel gern.} « Les grands-parents visitent volontiers les petits-enfants. » On reprend \all{besuchen die Enkel gern} et on remplace \all{wer} par le sujet trouvé dans le texte : \all{die Großeltern}.

\textbf{Question 2 :} \all{Wo wohnen die Großeltern oft ?} (Où habitent souvent les grands-parents ?)

Réponse complète : \all{Die Großeltern wohnen oft nicht in der Familie, sondern in einem eigenen Haus.} « Les grands-parents n'habitent souvent pas dans la famille, mais dans leur propre maison. »

\textbf{Question 3 :} \all{Was ist eine Patchworkfamilie ?} (Qu'est-ce qu'une famille recomposée ?)

Réponse complète : \all{Eine Patchworkfamilie ist eine Familie, in der die Eltern Kinder aus früheren Ehen zusammen erziehen.} « Une famille recomposée est une famille dans laquelle les parents élèvent ensemble des enfants de mariages précédents. »

\begin{exemplebox}[Phrases du texte à connaître]
\all{Viele Kinder leben bei einer allein erziehenden Mutter.} « Beaucoup d'enfants vivent avec une mère qui élève seule. »

\all{Die Eltern sind geschieden.} « Les parents sont divorcés. »

\all{Die Geschwister spielen zusammen.} « Les frères et sœurs jouent ensemble. »
\end{exemplebox}

\courssub{Vrai ou faux : Richtig oder falsch ? Begründen Sie !}

L'exercice vrai/faux exige une \cle{justification} par une citation du texte. La consigne \all{Begründen Sie !} signifie « Justifiez ! » : il faut recopier la phrase du texte qui prouve votre réponse, introduite par \all{Im Text steht: „...“} (« Dans le texte, il est écrit : ... »).

\begin{exoresolu}[Richtig oder falsch ?]
Énoncé. Répondez par \all{richtig} ou \all{falsch} et justifiez par une citation du texte \emph{Familien heute} :
\begin{enumerate}
\item \all{Alle Kinder leben bei Vater und Mutter.}
\item \all{Die Großeltern besuchen die Enkel gern.}
\item \all{In einer Patchworkfamilie leben nur die eigenen Kinder.}
\end{enumerate}
\tcblower
Solution détaillée.
\begin{enumerate}
\item \all{Falsch. Im Text steht: „Viele Kinder leben bei einer allein erziehenden Mutter.“} — Faux. Le texte dit que beaucoup d'enfants vivent avec une mère seule, donc pas tous avec le père et la mère.
\item \all{Richtig. Im Text steht: „Die Großeltern besuchen die Enkel gern.“} — Vrai. C'est exactement la phrase du texte.
\item \all{Falsch. Im Text steht: „In der Patchworkfamilie leben auch Kinder aus früheren Ehen.“} — Faux. Dans la famille recomposée vivent aussi des enfants de mariages précédents, pas seulement les enfants communs.
\end{enumerate}
\end{exoresolu}

\begin{attention}[Ne jamais répondre seulement « richtig » ou « falsch »]
Sans justification, la réponse vaut zéro point. La citation doit être recopiée \emph{exactement} comme dans le texte, entre guillemets allemands „...“. Si la phrase est \all{falsch}, la citation doit montrer la contradiction.
\end{attention}

\courssub{Trouver les équivalents dans le texte : Finden Sie im Text}

La consigne \all{Finden Sie im Text} (« Trouvez dans le texte ») demande de repérer l'équivalent allemand d'un mot français. Stratégie : penser au sens, repérer la zone du texte qui parle de cette idée, puis recopier le mot allemand \emph{avec son article} s'il s'agit d'un nom.

\begin{center}
\begin{tabularx}{\linewidth}{|X|X|X|}
\hline
\rowcolor{popblueL} \textbf{Français} & \textbf{Deutsch (dans le texte)} & \textbf{Remarque} \\
\hline
divorcé & \all{geschieden} & adjectif : \all{Die Eltern sind geschieden.} \\
\hline
les frères et sœurs & \all{die Geschwister} & toujours au pluriel, pas de singulier courant \all{*das Geschwister} \\
\hline
vivre ensemble & \all{zusammenleben} & verbe à particule séparable : infinitif \all{zusammenleben}, conjugué \all{sie leben zusammen} \\
\hline
la mère qui élève seule & \all{die allein erziehende Mutter} & expression figée du thème (participe présent adjectivé) \\
\hline
les petits-enfants & \all{die Enkel} & pluriel identique au singulier (\all{der Enkel}) \\
\hline
\end{tabularx}
\end{center}

\courssub{Entraînement avec un nouveau texte}

Appliquons maintenant toute la méthode à un texte nouveau et court, dans l'esprit du texte support, ancré dans le contexte camerounais.

\begin{textebox}[Familie Ngo Bell in Yaoundé]
Familie Ngo Bell wohnt in Yaoundé. Vater Jean und Mutter Marthe sind verheiratet und haben drei Kinder: zwei Söhne und eine Tochter. Die Familie ist eine Kernfamilie, aber sie ist nicht allein. Die Großmutter wohnt auch im Haus. Sie hilft der Mutter jeden Tag: Sie kocht, und sie erzählt den Enkeln Geschichten. Die Kinder lieben die Großmutter sehr.

Der älteste Sohn heißt Paul. Er ist siebzehn Jahre alt und besucht das Gymnasium. Seine Geschwister sind jünger. Am Wochenende besucht die Familie oft den Onkel in Douala. Dort lebt auch eine Tante mit ihren zwei Kindern. Die Kinder spielen zusammen Fußball.

Viele Familien in Kamerun sind groß. Onkel, Tanten und Cousins sind wichtig. Aber es gibt auch allein erziehende Mütter und Väter. Familie Ngo Bell hilft zum Beispiel einer Nachbarin: Sie ist geschieden und erzieht ihre zwei Kinder allein. „Familie ist sehr wichtig“, sagt Vater Jean, „und Familie bedeutet auch helfen.“
\end{textebox}

\textbf{Glossaire :} \all{verheiratet} (marié) ; \all{die Kernfamilie, -n} (la famille nucléaire) ; \all{der Enkel, - (Pl. die Enkel)} (le petit-fils / les petits-enfants) ; \all{der Onkel, -} (l'oncle) ; \all{die Tante, -n} (la tante) ; \all{der Cousin, -s} (le cousin) ; \all{die Nachbarin, -nen} (la voisine) ; \all{helfen} (aider, + datif) ; \all{bedeuten} (signifier).

\begin{savaistu}[La famille au Cameroun et en Allemagne]
Au Cameroun, la \cle{famille élargie} (\all{die Großfamilie} / \all{der erweiterte Familienkreis}) joue un rôle central : oncles, tantes, cousins, grands-parents vivent souvent ensemble ou à proximité et partagent l'éducation des enfants (\all{gemeinsame Erziehung}). En Allemagne, la \cle{famille nucléaire} (\all{die Kernfamilie}) est le modèle dominant, mais les familles monoparentales (\all{alleinerziehende Eltern}) et recomposées (\all{Patchworkfamilien}) sont en forte augmentation. Le texte \emph{Familie Ngo Bell} montre une situation mixte : une famille nucléaire qui intègre la grand-mère et maintient des liens forts avec la famille élargie à Douala — une réalité très fréquente au Cameroun.
\end{savaistu}

\begin{exercice}[Questions sur le texte « Familie Ngo Bell »]
Répondez par des phrases complètes en allemand :
\begin{enumerate}
\item \all{Wo wohnt Familie Ngo Bell ?}
\item \all{Wie viele Kinder haben Vater Jean und Mutter Marthe ?}
\item \all{Wer wohnt auch im Haus ?}
\item \all{Was macht die Familie am Wochenende ?}
\item \all{Richtig oder falsch ? Begründen Sie ! Die Nachbarin ist verheiratet.}
\end{enumerate}
\end{exercice}

\begin{corrige}[Exercice : Familie Ngo Bell]
\begin{enumerate}
\item \all{Familie Ngo Bell wohnt in Yaoundé.} — La famille Ngo Bell habite à Yaoundé.
\item \all{Vater Jean und Mutter Marthe haben drei Kinder.} — Ils ont trois enfants.
\item \all{Die Großmutter wohnt auch im Haus.} — La grand-mère habite aussi dans la maison.
\item \all{Am Wochenende besucht die Familie oft den Onkel in Douala.} — Le week-end, la famille visite souvent l'oncle à Douala. (Noter le verbe \all{besucht} en 2e position après le complément de temps \all{Am Wochenende}.)
\item \all{Falsch. Im Text steht: „Sie ist geschieden und erzieht ihre zwei Kinder allein.“} — Faux : la voisine est divorcée, pas mariée.
\end{enumerate}
\end{corrige}

\courssub{Le résumé oral : 2 à 3 phrases avec und, aber, auch}

Dernière étape de la méthode : résumer le texte oralement en deux ou trois phrases. On utilise des \cle{connecteurs simples} qui ne changent pas l'ordre des mots (ce ne sont pas des conjonctions de subordination) : \all{und} (et), \all{aber} (mais), \all{auch} (aussi). Après \all{und} et \all{aber}, le sujet vient normalement, puis le verbe en 2e position.

\begin{exemplebox}[Modèle de résumé oral]
\all{Der Text handelt von den Familien heute, und es gibt viele Familienformen. Viele Kinder leben bei einer allein erziehenden Mutter, aber die Großeltern sind auch wichtig. Die Patchworkfamilie ist auch eine moderne Familienform.}

« Le texte parle des familles aujourd'hui, et il existe beaucoup de formes de familles. Beaucoup d'enfants vivent avec une mère seule, mais les grands-parents sont aussi importants. La famille recomposée est aussi une forme de famille moderne. »
\end{exemplebox}

\begin{experience}[Résumé oral en binôme]
Par groupes de deux : l'élève A résume le texte \emph{Familien heute} en trois phrases avec \all{und}, \all{aber}, \all{auch} ; l'élève B écoute et vérifie que chaque phrase commence par une majuscule, que le verbe est en 2e position et que les trois idées principales sont présentes. Puis on échange les rôles avec le texte \emph{Familie Ngo Bell}.
\end{experience}

\begin{aretenir}[L'essentiel de la compréhension de texte]
\begin{itemize}
\item Deux lectures : globale (le thème) puis détaillée (les informations précises).
\item Le mot interrogatif annonce la réponse : \all{wer} → une personne, \all{wo} → un lieu, \all{was} → une chose, \all{warum} → une raison, \all{wie viele} → un nombre.
\item Toujours répondre par une phrase complète : majuscule initiale, verbe conjugué en 2e position.
\item Vrai/faux : toujours justifier par une citation exacte introduite par \all{Im Text steht: „...“}.
\item Résumé : 2 à 3 phrases reliées par \all{und}, \all{aber}, \all{auch}.
\end{itemize}
\end{aretenir}

\coursec{Lexique thématique : die Familie}

Après avoir travaillé la compréhension globale et détaillée du texte \emph{Familien heute}, nous allons maintenant fixer et approfondir le vocabulaire essentiel du thème. Cette section vous permettra de nommer précisément les membres de la famille, de décrire les situations matrimoniales, de classer les types de familles et d'utiliser les verbes et expressions courants. Comme toujours, chaque nom sera appris avec \textbf{son article} et \textbf{son pluriel} — condition indispensable pour éviter les erreurs d'accord par la suite.

\courssub{Champ 1 — Les membres de la famille (Die Familienmitglieder)}

Observons d'abord une phrase extraite de notre texte support :
\all{Meine Großeltern besuchen uns oft und bringen den Enkeln Geschenke mit.}
« Mes grands-parents nous rendent souvent visite et apportent des cadeaux aux petits-enfants. »

On y relève plusieurs noms de parenté au pluriel (\all{Großeltern}, \all{Enkeln}) et au cas datif pluriel (\all{Enkeln}). Le vocabulaire de la parenté suit des règles de formation régulières (souvent Umlaut au pluriel pour les monosyllabes) mais comporte des irrégularités à mémoriser par cœur.

\begin{definition}[Le mot collectif \cle{Geschwister}]
\all{die Geschwister} (toujours au pluriel) = « les frères et sœurs » (ensemble).
On ne dit jamais « \emph{ein Geschwister} » pour dire « un frère » ou « une sœur ». Pour l'indéfini, on utilise \all{ein Bruder} ou \all{eine Schwester}. En français, il n'existe pas de mot unique équivalent : on est obligé de dire « frères et sœurs » ou « fratrie » (terme plus formel).
\end{definition}

\begin{propriete}[Les principaux noms de parenté — Article, Pluriel, Umlaut]
Le tableau ci-dessous regroupe les termes du programme. Les flèches indiquent la modification du radical (Umlaut : a/au/o/u → ä/äu/ö/ü) ou l'ajout de suffixe.
\begin{center}
\begin{tabularx}{\linewidth}{|X|X|X|X|}
\hline
\rowcolor{popblueL} \textbf{Singulier (Article + Nom)} & \textbf{Pluriel} & \textbf{Français} & \textbf{Remarque} \\
\hline
der Vater & die \textbf{Väter} & le père & Umlaut a→ä \\
die Mutter & die \textbf{Mütter} & la mère & Umlaut u→ü \\
die Eltern (Pl.) & — & les parents & \emph{Plurale tantum} \\
der Sohn & die \textbf{Söhne} & le fils & Umlaut o→ö + -e \\
die Tochter & die \textbf{Töchter} & la fille & Umlaut o→ö + -er \\
das Kind & die \textbf{Kinder} & l'enfant & Pluriel en -er (irrégulier) \\
der Bruder & die \textbf{Brüder} & le frère & Umlaut u→ü \\
die Schwester & die \textbf{Schwestern} & la sœur & Pluriel régulier en -n \\
die Geschwister (Pl.) & — & les frères et sœurs & Collectif, pas de singulier \\
der Großvater & die \textbf{Großväter} & le grand-père & Composé de Vater \\
die Großmutter & die \textbf{Großmütter} & la grand-mère & Composé de Mutter \\
die Großeltern (Pl.) & — & les grands-parents & Plurale tantum \\
der Großonkel & die \textbf{Großonkel} & le grand-oncle & Composé de Onkel \\
die Großtante & die \textbf{Großtanten} & la grand-tante & Composé de Tante \\
der Enkel & die \textbf{Enkel} & le petit-fils & Pluriel sans changement \\
die Enkelin & die \textbf{Enkelinnen} & la petite-fille & -in (féminin) + -nen \\
der Onkel & die \textbf{Onkel} & l'oncle & Pluriel sans changement (emprunt fr.) \\
die Tante & die \textbf{Tanten} & la tante & Pluriel en -n (emprunt fr.) \\
der Cousin & die \textbf{Cousins} & le cousin & Pluriel en -s (emprunt fr.) \\
die Cousine & die \textbf{Cousinen} & la cousine & Pluriel en -nen \\
\hline
\end{tabularx}
\end{center}
\end{propriete}

\begin{exemplebox}[Phrases modèles avec les membres de la famille]
\all{Mein Vater heißt Paul und meine Mutter heißt Anna.} — « Mon père s'appelle Paul et ma mère s'appelle Anna. » \\
\all{Wir haben zwei Söhne und eine Tochter.} — « Nous avons deux fils et une fille. » \\
\all{Meine Geschwister sind älter als ich.} — « Mes frères et sœurs sont plus âgés que moi. » \\
\all{Die Großeltern kommen am Wochenende zu Besuch.} — « Les grands-parents viennent en visite le week-end. » \\
\all{Der Onkel meiner Mutter ist mein Großonkel.} — « L'oncle de ma mère est mon grand-oncle. »
\end{exemplebox}

\begin{savaistu}[Petit-fils / Petite-fille : la distinction de genre]
En français, « petit-fils » et « petite-fille » se distinguent par le nom (\emph{fils} vs \emph{fille}). En allemand, la distinction se fait par le suffixe féminin \textbf{-in} ajouté au masculin : \all{der Enkel} (petit-fils) → \all{die Enkelin} (petite-fille). Au pluriel : \all{die Enkel} (petits-fils / petits-enfants mixtes) vs \all{die Enkelinnen} (petites-filles). Cela vaut pour presque tous les noms de personnes : \all{der Lehrer} / \all{die Lehrerin}, \all{der Freund} / \all{die Freundin}.
\end{savaistu}

\courssub{Champ 2 — L'état civil et les situations familiales (Stand und Situation)}

Pour décrire la structure d'une famille, il faut maîtriser les adjectifs d'état civil et les noms désignant le conjoint ou le partenaire.

\begin{propriete}[Adjectifs d'état civil (invariables à l'attribut, accordés à l'épithète)]
\begin{center}
\begin{tabularx}{\linewidth}{|X|X|X|}
\hline
\rowcolor{poporangeL} \textbf{Allemand} & \textbf{Français} & \textbf{Exemple (attribut)} \\
\hline
verheiratet & marié(e) & \all{Er ist verheiratet.} (Il est marié.) \\
ledig & célibataire & \all{Sie ist ledig.} (Elle est célibataire.) \\
geschieden & divorcé(e) & \all{Mein Bruder ist geschieden.} \\
verwitwet & veuf / veuve & \all{Die Großmutter ist verwitwet.} \\
allein erziehend & qui élève seul(e) (parent solo) & \all{Meine Schwester ist allein erziehend.} \\
\hline
\end{tabularx}
\end{center}
\emph{Note :} À l'attribut (après \all{sein}), l'adjectif ne prend \textbf{pas} de terminaison. En épithète (devant le nom), il s'accorde : \all{ein verheirateter Mann}, \all{eine verheiratete Frau}.
\end{propriete}

\begin{definition}[Conjoint et partenaire]
\begin{itemize}
\item \all{der Ehemann} (die \textbf{Ehemänner}) — l'époux, le mari.
\item \all{die Ehefrau} (die \textbf{Ehefrauen}) — l'épouse, la femme.
\item \all{der Partner} (die \textbf{Partner}) — le partenaire (non marié, masculin).
\item \all{die Partnerin} (die \textbf{Partnerinnen}) — la partenaire (non mariée, féminin).
\end{itemize}
\emph{Usage :} \all{Ehemann/Ehefrau} implique un mariage civil ou religieux. \all{Partner/Partnerin} s'utilise pour les couples non mariés (Lebensgemeinschaft) ou de façon neutre.
\end{definition}

\begin{exemplebox}[Décrire une situation familiale]
\all{Mein Onkel ist geschieden und lebt jetzt mit seiner neuen Partnerin zusammen.} — « Mon oncle est divorcé et vit maintenant avec sa nouvelle compagne. » \\
\all{Sie ist verwitwet und allein erziehend; sie hat drei Kinder.} — « Elle est veuve et élève seule ; elle a trois enfants. »
\end{exemplebox}

\courssub{Champ 3 — Les types de familles (Familienformen)}

La société germanophone (comme la société camerounaise) connaît une grande diversité de structures familiales. Voici les termes officiels du programme.

\begin{propriete}[Types de familles — Noms composés, genres et pluriels]
\begin{center}
\begin{tabularx}{\linewidth}{|X|X|X|}
\hline
\rowcolor{popgreenL} \textbf{Allemand (Singulier)} & \textbf{Pluriel} & \textbf{Français} \\
\hline
die Kernfamilie & die Kernfamilien & la famille nucléaire (parents + enfants) \\
die Großfamilie & die Großfamilien & la famille élargie (plusieurs générations) \\
die Patchworkfamilie & die Patchworkfamilien & la famille recomposée \\
die Ein-Eltern-Familie & die Ein-Eltern-Familien & la famille monoparentale \\
das kinderlose Paar & die kinderlosen Paare & le couple sans enfants \\
\hline
\end{tabularx}
\end{center}
\emph{Analyse morphologique :}
\begin{itemize}
\item \all{Kernfamilie}, \all{Großfamilie}, \all{Patchworkfamilie}, \all{Ein-Eltern-Familie} : composés en \textbf{-familie} (féminin, pluriel en \textbf{-n}).
\item \all{das kinderlose Paar} : forme déclinée avec article défini de l'adjectif \all{kinderlos} (sans enfant) + \all{das Paar} (le couple, pluriel \all{die Paare}).
\end{itemize}
\end{propriete}

\begin{textebox}[Texte d'entraînement : Vielfalt der Familienformen]
In Deutschland gibt es viele verschiedene Familienformen. Die klassische \textbf{Kernfamilie} mit Vater, Mutter und Kindern ist immer noch häufig, aber nicht mehr die einzige Norm. Immer mehr Kinder wachsen in \textbf{Ein-Eltern-Familien} auf, weil sich die Eltern \textbf{getrennt} haben oder weil ein Elternteil \textbf{verwitwet} ist. Auch \textbf{Patchworkfamilien} sind normal geworden: Wenn Mutter oder Vater einen neuen \textbf{Partner} hat, der eigene Kinder mitbringt, entstehen neue Geschwisterbeziehungen zwischen \textbf{Stiefgeschwistern}. In \textbf{Großfamilien}, wie man sie oft in ländlichen Gegenden oder in migrantischen Communities findet, leben manchmal drei Generationen unter einem Dach: \textbf{Großeltern}, Eltern und Enkel unterstützen sich gegenseitig. Schließlich gibt es \textbf{kinderlose Paare}, die bewusst oder ungewollt keine Kinder haben. Alle diese Formen sind rechtlich anerkannt und sozial akzeptiert.

\medskip
\noindent\textbf{Glossaire :}
\begin{itemize}
\item \all{die Norm} (die Normen) — la norme
\item \all{aufwachsen} (wächst auf, ist aufgewachsen) — grandir
\item \all{sich trennen} (trennt sich, hat sich getrennt) — se séparer
\item \all{mitbringen} (bringt mit, hat mitgebracht) — apporter, amener
\item \all{entstehen} (entsteht, ist entstanden) — naître, se former
\item \all{das Stiefgeschwister} (die Stiefgeschwister) — le demi-frère / la demi-sœur (par alliance)
\item \all{unter einem Dach} — sous un même toit
\item \all{sich gegenseitig unterstützen} — se soutenir mutuellement
\item \all{bewusst / ungewollt} — consciemment / involontairement
\item \all{rechtlich anerkannt} — légalement reconnu
\item \all{sozial akzeptiert} — socialement accepté
\end{itemize}

\medskip
\noindent\textbf{Questions de compréhension :}
\begin{enumerate}
\item Nenne drei Familienformen, die im Text erwähnt werden.
\item Was bedeutet „Patchworkfamilie“? Erkläre mit eigenen Worten.
\item Warum leben in Großfamilien oft drei Generationen zusammen?
\end{enumerate}
\end{textebox}

\courssub{Champ 4 — Verbes et expressions utiles (Verben und Redewendungen)}

Ces verbes permettent de décrire la vie familiale, les événements (mariage, naissance, séparation) et les relations affectives.

\begin{propriete}[Verbes thématiques — Présent, Participe passé, Auxiliaire]
\begin{center}
\begin{tabularx}{\linewidth}{|X|X|X|X|X|}
\hline
\rowcolor{poppurpleL} \textbf{Infinitif} & \textbf{Présent (ich/er/sie/es)} & \textbf{Perfekt} & \textbf{Aux.} & \textbf{Français} \\
\hline
heiraten & ich heirate / er heiratet & hat geheiratet & haben & se marier (avec qqn) \\
leben & ich lebe / er lebt & hat gelebt & haben & vivre (habiter) \\
zusammenleben & ich lebe zusammen / er lebt zusammen & hat zusammengelebt & haben & vivre ensemble (en couple) \\
sich trennen & ich trenne mich / er trennt sich & hat sich getrennt & haben & se séparer (divorcer/quitter) \\
erziehen & ich erziehe / er erzieht & hat erzogen & haben & élever (éduquer) \\
besuchen & ich besuche / er besucht & hat besucht & haben & rendre visite à \\
lieben & ich liebe / er liebt & hat geliebt & haben & aimer \\
\hline
\end{tabularx}
\end{center}
\emph{Attention :} \all{heiraten} est transitif direct en allemand (\all{jemanden heiraten} — épouser qqn), pas de préposition. \all{besuchen} aussi : \all{jemanden besuchen}.
\end{propriete}

\begin{exemplebox}[Verbes en contexte]
\all{Meine Schwester heiratet nächsten Sommer.} — « Ma sœur se marie l'été prochain. » \\
\all{Wir leben seit zehn Jahren zusammen.} — « Nous vivons ensemble depuis dix ans. » \\
\all{Eltern erziehen ihre Kinder mit Liebe.} — « Les parents élèvent leurs enfants avec amour. » \\
\all{Die Großeltern besuchen uns jeden Sonntag.} — « Les grands-parents nous rendent visite chaque dimanche. »
\end{exemplebox}

\begin{attention}[Faux ami majeur : \cle{bekommen} $\neq$ devenir]
\all{bekommen} (reçoit, a reçu) signifie \textbf{« recevoir », « obtenir »}, jamais « devenir ».
\begin{itemize}
\item \all{Sie bekommt ein Kind.} = « Elle attend un enfant / Elle a un enfant » (réception/naissance), \textbf{pas} « Elle devient un enfant ».
\item \all{Ich habe ein Geschenk bekommen.} = « J'ai reçu un cadeau. »
\end{itemize}
Pour « devenir », on utilise \all{werden} (wird, ist geworden) :
\all{Er wird Arzt.} = « Il devient médecin. » \\
\all{Das Kind wird krank.} = « L'enfant tombe malade / devient malade. »
\end{attention}

\courssub{Familles de mots (Wortfamilien) — Dérivation et composition}

L'allemand construit son vocabulaire par dérivation (suffixes) et composition. Repérer la racine permet de deviner le sens de mots inconnus.

\begin{propriete}[Principales familles de mots du thème]
\begin{center}
\begin{tabularx}{\linewidth}{|X|X|X|}
\hline
\rowcolor{poppinkL} \textbf{Racine (Nom/Verbe)} & \textbf{Dérivés (Adjectifs / Noms abstraits)} & \textbf{Sens} \\
\hline
die Familie & \all{familiär} (adj.), \all{die Familiarität} & familial, familiarité \\
heiraten (v.) & \all{die Heirat} (n.f.), \all{verheiratet} (adj.) & le mariage, marié \\
erziehen (v.) & \all{die Erziehung} (n.f.), \all{erzieherisch} (adj.) & l'éducation, éducatif \\
das Kind & \all{die Kindheit} (n.f.), \all{kinderlos} (adj.), \all{kindlich} (adj.) & l'enfance, sans enfant, enfantin \\
die Ehe & \all{ehelich} (adj.), \all{die Eheleute} (Pl.) & conjugal, les époux \\
\hline
\end{tabularx}
\end{center}
\end{propriete}

\begin{methode}[Mémoriser le vocabulaire nominal en 3 temps]
Pour chaque nom nouveau, appliquez systématiquement ce rituel :
\begin{enumerate}
\item \textbf{Article + Singulier + Pluriel} : écrivez \all{der Sohn, die Söhne} ; \all{die Tochter, die Töchter} ; \all{das Kind, die Kinder}. Apprenez le triplet comme un bloc indissociable.
\item \textbf{Phrase personnelle} : écrivez une phrase vraie vous concernant. Ex. \all{Ich habe einen Bruder und zwei Schwestern.}
\item \textbf{Contexte grammatical} : déclinez le nom aux \textbf{quatre cas} (nominatif, accusatif, datif, génitif) au singulier et au pluriel avec l'article défini. Ex. \all{der Sohn, den Sohn, dem Sohn, des Sohnes} / \all{die Söhne, die Söhne, den Söhnen, der Söhne}.
\end{enumerate}
Cette méthode ancre le genre, le pluriel et la déclinaison simultanément.
\end{methode}

\courssub{Complément Bac — La belle-famille (Schwieger- / Stief-)}

Ces termes ne sont pas au programme de base mais tombent fréquemment aux épreuves du Baccalauréat (compréhension ou expression).

\begin{definition}[Beaux-parents et beaux-enfants — Deux séries distinctes]
\begin{itemize}
\item \textbf{Par alliance (mariage) : préfixe \cle{Schwieger-}}
\begin{itemize}
\item \all{der Schwiegervater} (die Schwiegerväter) — le beau-père (père du conjoint)
\item \all{die Schwiegermutter} (die Schwiegermütter) — la belle-mère (mère du conjoint)
\item \all{der Schwiegersohn} (die Schwiegersöhne) — le gendre
\item \all{die Schwiegertochter} (die Schwiegertöchter) — la bru
\end{itemize}
\item \textbf{Par recomposition (remariage) : préfixe \cle{Stief-}}
\begin{itemize}
\item \all{der Stiefvater} (die Stiefväter) — le beau-père (nouveau mari de la mère)
\item \all{die Stiefmutter} (die Stiefmütter) — la belle-mère (nouvelle femme du père)
\item \all{das Stiefkind} (die Stiefkinder) — le beau-fils / la belle-fille (enfant du conjoint)
\item \all{der Stiefbruder} / \all{die Stiefschwester} — le demi-frère / la demi-sœur (par alliance)
\end{itemize}
\end{itemize}
\emph{Astuce mnémotechnique :} \all{Schwieger-} sonne comme « suisse » (lien par \emph{échange} d'anneaux/mariage) ; \all{Stief-} vient de \all{stief} (ancien : « privé de parents ») → lien par \emph{deuil ou séparation} puis recomposition.
\end{definition}

\courssub{Visualisation : Arbre généalogique simplifié}

Le schéma ci-dessous résume les relations verticales (générations) et horizontales (fratrie, couple). Les flèches indiquent la relation « est parent de » (direction haut→bas).

\begin{popfigure}
\begin{tikzpicture}[
  node distance=1.2cm and 1.5cm,
  every node/.style={font=\small, align=center},
  male/.style={rectangle, draw=popblue, thick, fill=popblueL, rounded corners=3pt, minimum width=2.2cm},
  female/.style={rectangle, draw=poppink, thick, fill=poppinkL, rounded corners=3pt, minimum width=2.2cm},
  couple/.style={rectangle, draw=poppurple, thick, fill=poppurpleL, rounded corners=3pt, minimum width=3cm},
  arrow/.style={-Stealth, popdark, thick}
]
% Génération 1 : Großeltern
\node[male, align=center] (gv1){Großvater\\ (Vater des Vaters)};
\node[female, right=of gv1, align=center] (gm1){Großmutter\\ (Mutter des Vaters)};
\node[couple, below=of $(gv1)!0.5!(gm1)$, align=center] (gp1){Großeltern\\ (väterlicherseits)};

\node[male, right=3cm of gm1, align=center] (gv2){Großvater\\ (Vater der Mutter)};
\node[female, right=of gv2, align=center] (gm2){Großmutter\\ (Mutter der Mutter)};
\node[couple, below=of $(gv2)!0.5!(gm2)$, align=center] (gp2){Großeltern\\ (mütterlicherseits)};

% Génération 2 : Eltern
\node[male, below=1.5cm of gp1] (vater) {Vater};
\node[female, right=of vater] (mutter) {Mutter};
\node[couple, below=of $(vater)!0.5!(mutter)$, align=center] (eltern){Eltern\\ (Ehepaar)};

% Génération 3 : Kinder
\node[male, below=1.5cm of eltern, xshift=-2cm] (sohn) {Sohn / Bruder};
\node[female, right=of sohn] (tochter) {Tochter / Schwester};
\node[male, right=of tochter] (sohn2) {Sohn / Bruder};

% Flèches
\draw[arrow] (gp1) -- (vater);
\draw[arrow] (gp2) -- (mutter);
\draw[arrow] (eltern) -- (sohn);
\draw[arrow] (eltern) -- (tochter);
\draw[arrow] (eltern) -- (sohn2);

% Étiquettes générations
\node[above=0.2cm of gp1, font=\bfseries\small, text=popblue] {1. Generation: Großeltern};
\node[above=0.2cm of eltern, font=\bfseries\small, text=poporange] {2. Generation: Eltern};
\node[below=0.2cm of sohn, font=\bfseries\small, text=popgreen] {3. Generation: Kinder / Enkel};
\end{tikzpicture}
\legende{Arbre généalogique à trois générations : termes allemands clés. Bleu = masculin, Rose = féminin, Violet = couple.}
\end{popfigure}

\courssub{Exercice résolu : Construire un arbre généalogique lexical}

\begin{exoresolu}[Compléter le tableau de parenté]
\textbf{Énoncé :} Complétez le tableau ci-dessous en donnant le terme allemand correspondant (article + singulier + pluriel) et le lien de parenté précis.

\begin{center}
\begin{tabularx}{\linewidth}{|X|X|X|}
\hline
\rowcolor{popgoldL} \textbf{Français} & \textbf{Allemand (Art. + Sg. + Pl.)} & \textbf{Précision} \\
\hline
Le père de mon père & & \\
La mère de ma mère & & \\
Le fils de mon frère & & \\
La fille de ma sœur & & \\
Le frère de mon père & & \\
La sœur de ma mère & & \\
L'époux de ma tante & & \\
L'épouse de mon oncle & & \\
Le nouveau mari de ma mère (divorcée) & & \\
La nouvelle femme de mon père (veuf) & \\
\hline
\end{tabularx}
\end{center}

\tcblower

\textbf{Solution détaillée :}
\begin{center}
\begin{tabularx}{\linewidth}{|X|X|X|}
\hline
\rowcolor{popgoldL} \textbf{Français} & \textbf{Allemand (Art. + Sg. + Pl.)} & \textbf{Précision} \\
\hline
Le père de mon père & der Großvater, die Großväter & Großvater väterlicherseits \\
La mère de ma mère & die Großmutter, die Großmütter & Großmutter mütterlicherseits \\
Le fils de mon frère & der Neffe, die Neffen & (Neveu — vocabulaire étendu) \\
La fille de ma sœur & die Nichte, die Nichten & (Nièce — vocabulaire étendu) \\
Le frère de mon père & der Onkel, die Onkel & Onkel väterlicherseits \\
La sœur de ma mère & die Tante, die Tanten & Tante mütterlicherseits \\
L'époux de ma tante & der Onkel, die Onkel & Par alliance (Onkel) \\
L'épouse de mon oncle & die Tante, die Tanten & Par alliance (Tante) \\
Le nouveau mari de ma mère (divorcée) & der Stiefvater, die Stiefväter & Recomposition (Stief-) \\
La nouvelle femme de mon père (veuf) & die Stiefmutter, die Stiefmütter & Recomposition (Stief-) \\
\hline
\end{tabularx}
\end{center}
\emph{Remarque :} \all{Neffe} (neveu) et \all{Nichte} (nièce) ne sont pas au programme officiel de Première mais sont indispensables pour décrire une famille élargie. Apprenez-les par la même occasion.
\end{exoresolu}

\courssub{Exercices d'application}

\begin{exercice}[Vocabulaire : Articles, pluriels et phrases]
\begin{enumerate}
\item Donnez le pluriel avec article : \all{der Großvater}, \all{die Enkelin}, \all{das Stiefkind}, \all{die Patchworkfamilie}, \all{der Schwiegersohn}.
\item Traduisez en allemand :
      \begin{enumerate}
      \item Ma sœur est mariée et a deux enfants.
      \item Nos grands-parents nous rendent visite dimanche.
      \item C'est une famille monoparentale ; la mère élève seule.
      \item Mon oncle est divorcé ; il vit avec sa nouvelle partenaire.
      \item Nous sommes une famille nombreuse (Großfamilie) : trois générations sous un même toit.
      \end{enumerate}
\item Mettez l'adjectif à la forme correcte (attribut ou épithète) :
      \begin{enumerate}
      \item Mein Onkel ist \_\_\_\_\_\_\_\_ (verheiratet).
      \item Eine \_\_\_\_\_\_\_\_ (allein erziehend) Mutter hat viel zu tun.
      \item Das Paar ist \_\_\_\_\_\_\_\_ (kinderlos).
      \item Er kommt aus einer \_\_\_\_\_\_\_\_ (groß) Familie.
      \end{enumerate}
\end{enumerate}
\end{exercice}


\begin{corrige}[Exercice 1]
\begin{enumerate}
\item \all{die Großväter} ; \all{die Enkelinnen} ; \all{die Stiefkinder} ; \all{die Patchworkfamilien} ; \all{die Schwiegersöhne}.
\item \begin{enumerate}
      \item \all{Meine Schwester ist verheiratet und hat zwei Kinder.}
      \item \all{Unsere Großeltern besuchen uns am Sonntag.}
      \item \all{Das ist eine Ein-Eltern-Familie; die Mutter ist allein erziehend.}
      \item \all{Mein Onkel ist geschieden; er lebt mit seiner neuen Partnerin zusammen.}
      \item \all{Wir sind eine Großfamilie: drei Generationen leben unter einem Dach.}
      \end{enumerate}
\item \begin{enumerate}
      \item \all{verheiratet} (attribut, pas de terminaison)
      \item \all{allein erziehende} (épithète, nom fém. sing. après article indéfini → -e)
      \item \all{kinderlos} (attribut)
      \item \all{großen} (épithète, \all{aus} + datif → \all{einer großen Familie}, datif féminin singulier article indéfini → -en)
      \end{enumerate}
\end{enumerate}
\end{corrige}

\begin{exercice}[Expression guidée : Ma famille]
Rédigez un paragraphe de 60–80 mots en allemand pour présenter votre famille (réelle ou imaginaire). Utilisez au moins :
\begin{itemize}
\item 5 noms de parenté différents (avec articles corrects),
\item 2 adjectifs d'état civil,
\item 3 verbes du champ lexical (\all{leben}, \all{heiraten}, \all{erziehen}, \all{besuchen}, \all{lieben}…),
\item 1 type de famille (\all{Kernfamilie}, \all{Patchworkfamilie}…).
\end{itemize}
\emph{Aide au démarrage :} \all{Ich komme aus einer … Familie. Mein Vater … und meine Mutter … Sie sind … Ich habe … Geschwister. Wir … gern zusammen. Meine Großeltern …}
\end{exercice}


\begin{corrige}[Exercice 2 — Proposition de production]
\all{Ich komme aus einer Kernfamilie. Mein Vater heißt Thomas und meine Mutter heißt Sabine. Sie sind seit 25 Jahren verheiratet. Ich habe zwei Geschwister: einen älteren Bruder und eine jüngere Schwester. Mein Bruder heißt Markus und er ist ledig. Meine Schwester Anna ist verheiratet und hat ein Kind, meine Nichte Lisa. Wir leben in Yaoundé, aber meine Großeltern leben in einem Dorf in der Nähe von Dschang. Sie besuchen uns oft in den Ferien. Wir lieben uns sehr und unterstützen uns gegenseitig.}
\end{corrige}

\coursec{Grammaire 1 : l'ordre des mots et les articles/possessifs}

Dans la section précédente, nous avons appris le vocabulaire de la famille : \all{die Familie}, \all{die Eltern}, \all{die Großeltern}, \all{die Geschwister}, \all{verheiratet}, \all{allein erziehend}. Nous allons maintenant observer comment ces mots se comportent \emph{dans la phrase}. Car en allemand, savoir un mot ne suffit pas : il faut savoir \cle{où le placer} et \cle{sous quelle forme} l'employer. C'est l'objet de cette section : d'abord la place du verbe (la règle la plus importante de toute la grammaire allemande), puis les articles et les adjectifs possessifs.

\courssub{Observation : que remarque-t-on dans ces phrases ?}

Lisons d'abord un court texte qui reprend notre thème et que nous allons observer attentivement.

\begin{textebox}[Familien in Kamerun — Text zur Beobachtung]
Meine Familie lebt in Yaoundé. Wir sind sechs Personen : mein Vater, meine Mutter, meine zwei Brüder, meine Schwester und ich. Mein Vater arbeitet in einem Büro, und meine Mutter verkauft Gemüse auf dem Markt. Am Wochenende besuchen wir oft unsere Großeltern. Sie wohnen in einem Dorf bei Bafoussam. Mein Großvater hat einen großen Garten, und meine Großmutter kocht sehr gut.

Mein Onkel Paul ist verheiratet. Er hat drei Kinder. Seine Frau heißt Marie, und ihre Kinder gehen schon zur Schule. Meine Tante Chantal ist allein erziehend. Sie hat eine Tochter. Ihre Tochter heißt Sandrine und besucht meine Schule.

Heute besucht uns meine Tante. Sie bringt ihre Tochter mit. Am Abend essen wir zusammen, und wir sprechen über die Schule und den Fußball. Familien sind in Kamerun sehr wichtig : die Eltern helfen den Kindern, und die Kinder helfen später den Eltern.
\end{textebox}

\textbf{Glossaire :}
\begin{itemize}
\item \all{das Gemüse} (–) : les légumes
\item \all{der Markt}, \all{die Märkte} : le marché
\item \all{das Dorf}, \all{die Dörfer} : le village
\item \all{der Garten}, \all{die Gärten} : le jardin
\item \all{der Onkel} (–) : l'oncle
\item \all{die Tante} (–n) : la tante
\item \all{verheiratet} : marié(e)
\item \all{allein erziehend} : parent isolé
\item \all{mitbringen} (trennbar) : apporter (avec soi)
\item \all{helfen} (hilft, half, hat geholfen + Dat.) : aider
\item \all{später} : plus tard
\end{itemize}

Observons maintenant quatre phrases tirées du texte :

\begin{exemplebox}[Observation des positions]
\all{Meine Familie lebt in Yaoundé.} « Ma famille vit à Yaoundé. »

\all{Heute besucht uns meine Tante.} « Aujourd'hui, ma tante nous rend visite. »

\all{Mein Vater arbeitet in einem Büro.} « Mon père travaille dans un bureau. »

\all{Am Wochenende besuchen wir oft unsere Großeltern.} « Le week-end, nous rendons souvent visite à nos grands-parents. »
\end{exemplebox}

Dans les phrases 1 et 3, le sujet est en tête et le verbe le suit. Dans les phrases 2 et 4, un complément (\all{Heute}, \all{Am Wochenende}) est en tête : le verbe reste à la même place, et c'est le \emph{sujet} qui recule après le verbe. Le verbe conjugué ne bouge jamais de sa deuxième position.

\courssub{La règle d'or : le verbe conjugué en 2e position}

\begin{propriete}[Règle d'or de la phrase déclarative allemande]
Dans la phrase déclarative allemande, le \cle{verbe conjugué} occupe \textbf{toujours la 2e position}.

\begin{itemize}
\item \textbf{Position 1} : le sujet \emph{ou} un complément (de temps, de lieu, etc.) — un seul élément (groupe syntaxique) ;
\item \textbf{Position 2} : le verbe conjugué ;
\item \textbf{Position 3} : le reste de la phrase (sujet s'il n'est pas en tête, compléments, etc.).
\end{itemize}

Quand un complément ouvre la phrase, le sujet passe \emph{après} le verbe : c'est l'\cle{inversion sujet-verbe}.

\all{Heute besucht uns meine Tante.} « Aujourd'hui, ma tante nous rend visite. »
\end{propriete}

\begin{attention}[Attention : la « position » n'est pas un mot]
Un groupe nominal comme \all{Meine Familie} (deux mots) ou une expression prépositionnelle comme \all{Am Wochenende} (deux mots) forme \textbf{une seule position}. Ne compte pas les mots, compte les \cle{groupes syntaxiques}.
\end{attention}

\begin{popfigure}
\begin{tikzpicture}[node distance=0.5cm, >=Stealth]
\node[draw=popblue, fill=popblueL, rounded corners, minimum width=4.5cm, minimum height=1.8cm, align=center] (p1) {\textbf{Position 1}\\ Sujet \emph{ou} Complément};
\node[draw=poporange, fill=poporangeL, rounded corners, minimum width=3.5cm, minimum height=1.8cm, align=center, right=of p1] (p2) {\textbf{Position 2}\\ \textbf{VERBE conjugué}};
\node[draw=popgreen, fill=popgreenL, rounded corners, minimum width=4.5cm, minimum height=1.8cm, align=center, right=of p2] (p3) {\textbf{Position 3}\\ Reste de la phrase (sujet, compléments...)};
\draw[->, thick, popdark] (p1) -- (p2);
\draw[->, thick, popdark] (p2) -- (p3);
\node[draw=popdark, fill=popcream, rounded corners, minimum width=4.5cm, minimum height=1cm, align=center, below=0.8cm of p1] (e1) {\all{Meine Familie}};
\node[draw=popdark, fill=popcream, rounded corners, minimum width=3.5cm, minimum height=1cm, align=center, below=0.8cm of p2] (e2) {\all{lebt}};
\node[draw=popdark, fill=popcream, rounded corners, minimum width=4.5cm, minimum height=1cm, align=center, below=0.8cm of p3] (e3) {\all{in Douala.}};
\draw[->, popmuted] (p1) -- (e1);
\draw[->, popmuted] (p2) -- (e2);
\draw[->, popmuted] (p3) -- (e3);
\end{tikzpicture}
\legende{Structure de la phrase simple : \all{Meine Familie lebt in Douala.}}
\end{popfigure}

Vérifions l'inversion avec le même exemple transformé :

\begin{center}
\begin{tabularx}{\linewidth}{|X|X|X|}
\hline
\rowcolor{popblueL} Français & Deutsch & Remarque \\
\hline
Ma famille vit à Douala. & \all{Meine Familie lebt in Douala.} & Sujet en position 1 \\
\hline
À Douala, ma famille vit. & \all{In Douala lebt meine Familie.} & Complément en tête : inversion \\
\hline
Aujourd'hui, ma tante nous rend visite. & \all{Heute besucht uns meine Tante.} & Le français garde le sujet avant le verbe, pas l'allemand \\
\hline
Le week-end, nous visitons nos grands-parents. & \all{Am Wochenende besuchen wir unsere Großeltern.} & \all{wir} passe après \all{besuchen} \\
\hline
\end{tabularx}
\end{center}

\begin{attention}[Erreur typique des francophones]
En français, on dit « À Yaoundé, la famille vit » : le sujet reste avant le verbe. En allemand, c'est \textbf{impossible}. N'écris jamais :

\emph{Faux :} \all{In Yaoundé die Familie lebt.} \quad \all{Heute meine Tante besucht uns.} \quad \all{Meine Familie sie lebt in Douala.} (double sujet)

\emph{Correct :} \all{In Yaoundé lebt die Familie.} \quad \all{Heute besucht uns meine Tante.}

\cle{Règle pratique} : \textbf{un seul élément avant le verbe}, jamais deux. Dès que tu places un complément en tête, fais immédiatement suivre le verbe, puis le sujet.
\end{attention}

\courssub{La phrase interrogative}

\textbf{1. Question sans mot interrogatif} (réponse oui/non) : le verbe conjugué passe en \cle{1re position} (comme en français « As-tu... ? »).

\begin{exemplebox}[Questions fermées]
\all{Hast du Geschwister?} « As-tu des frères et sœurs ? »

\all{Besuchst du deine Großeltern?} « Rends-tu visite à tes grands-parents ? »

\all{Wohnen deine Eltern in Douala?} « Tes parents habitent-ils à Douala ? »
\end{exemplebox}

\textbf{2. Question avec mot interrogatif} : mot interrogatif (position 1) + verbe (position 2) + sujet. La règle de la 2e position est donc respectée !

\begin{exemplebox}[Questions ouvertes]
\all{Wo wohnen deine Großeltern?} « Où habitent tes grands-parents ? »

\all{Wie heißt deine Schwester?} « Comment s'appelle ta sœur ? »

\all{Wann besuchst du deine Familie?} « Quand rends-tu visite à ta famille ? »
\end{exemplebox}

\begin{aretenir}[Les mots interrogatifs utiles]
\all{wer} (qui), \all{was} (quoi, que), \all{wo} (où), \all{woher} (d'où), \all{wohin} (vers où), \all{wann} (quand), \all{wie} (comment), \all{warum} (pourquoi), \all{wie viele} (combien de).
\end{aretenir}

\courssub{Les articles : der, die, das / ein, eine, ein}

Tu connais déjà les articles au \cle{nominatif} (cas du sujet). Mais quand le nom est \cle{complément d'objet direct} (COD), il passe à l'\cle{accusatif}. Bonne nouvelle : \textbf{seul le masculin change de forme}.

\begin{propriete}[Nominatif et accusatif des articles]
\begin{itemize}
\item Le \textbf{nominatif} marque le sujet : \all{Der Vater arbeitet.} « Le père travaille. »
\item L'\textbf{accusatif} marque le COD : \all{Ich besuche den Vater.} « Je rends visite au père. »
\item À l'accusatif, \textbf{seul le masculin change} : \all{der} $\to$ \all{den}, \all{ein} $\to$ \all{einen}. Le féminin (\all{die}/\all{eine}), le neutre (\all{das}/\all{ein}) et le pluriel (\all{die}/\all{–}) restent identiques au nominatif.
\end{itemize}
\end{propriete}

\begin{center}
\begin{tabular}{|l|l|l|l|}
\hline
\rowcolor{popblueL} & Masculin & Féminin & Neutre \\
\hline
Nominatif défini & der Vater & die Mutter & das Kind \\
\hline
Accusatif défini & \textbf{den} Vater & die Mutter & das Kind \\
\hline
Nominatif indéfini & ein Bruder & eine Schwester & ein Kind \\
\hline
Accusatif indéfini & \textbf{einen} Bruder & eine Schwester & ein Kind \\
\hline
\end{tabular}
\end{center}

\begin{exemplebox}[Articles en action]
\all{Ich habe einen Bruder.} « J'ai un frère. » (accusatif masculin : \all{einen})

\all{Mein Bruder wohnt in Bafoussam.} « Mon frère habite à Bafoussam. » (nominatif : \all{Mein Bruder} est sujet)

\all{Sie besucht ihre Tante.} « Elle rend visite à sa tante. » (accusatif féminin : \all{ihre}, inchangé)

\all{Wir sehen das Kind.} « Nous voyons l'enfant. » (accusatif neutre : \all{das}, inchangé)
\end{exemplebox}

\courssub{Les adjectifs possessifs : mein, dein, sein...}

Le possessif s'accorde avec le \textbf{possesseur} (qui possède ?), puis se \textbf{décline comme} \all{ein} selon le genre, le nombre et le cas du nom possédé.

\begin{center}
\begin{tabular}{|l|l|l|}
\hline
\rowcolor{popblueL} Possesseur & Possessif (base) & Exemple (Nominatif) \\
\hline
ich (je) & mein & \all{mein Vater} (mon père) \\
\hline
du (tu) & dein & \all{deine Mutter} (ta mère) \\
\hline
er (il) & sein & \all{sein Bruder} (son frère) \\
\hline
sie (elle) & ihr & \all{ihre Schwester} (sa sœur) \\
\hline
wir (nous) & unser & \all{unser Haus} (notre maison) \\
\hline
ihr (vous) & euer & \all{eure Eltern} (vos parents) \\
\hline
sie / Sie (ils, elles, Vous poli) & ihr / Ihr & \all{ihre Kinder} (leurs / Vos enfants) \\
\hline
\end{tabular}
\end{center}

\begin{popfigure}
\begin{tikzpicture}
\matrix[matrix of nodes, nodes={draw=popline, minimum width=2.8cm, minimum height=0.85cm, align=center, anchor=center, font=\small}, row sep=-\pgflinewidth, column sep=-\pgflinewidth] (m) {
|[fill=popblueL]| & |[fill=popblueL]| Maskulin & |[fill=popblueL]| Feminin & |[fill=popblueL]| Neutrum & |[fill=popblueL]| Plural \\
|[fill=poporangeL]| ich & mein & meine & mein & meine \\
|[fill=poporangeL]| du & dein & deine & dein & deine \\
|[fill=poporangeL]| er & sein & seine & sein & seine \\
|[fill=poporangeL]| sie & ihr & ihre & ihr & ihre \\
|[fill=poporangeL]| wir & unser & unsere & unser & unsere \\
|[fill=poporangeL]| ihr & euer & eure & euer & eure \\
};
\end{tikzpicture}
\legende{Les possessifs au nominatif. À l'accusatif masculin : \all{meinen}, \all{deinen}, \all{seinen}, \all{ihren}, \all{unseren}, \all{euren}.}
\end{popfigure}

\begin{propriete}[Déclinaison du possessif]
Le possessif se décline exactement comme l'article indéfini \all{ein} :
\begin{itemize}
\item Nominatif : \all{mein Vater} (m.), \all{meine Mutter} (f.), \all{mein Kind} (n.), \all{meine Eltern} (pl.).
\item Accusatif : seul le masculin ajoute \textbf{-en} : \all{meinen Vater}, \all{deinen Bruder}, \all{seinen Onkel}, \all{ihren Mann}, \all{unseren Opa}, \all{euren Onkel}. Féminin, neutre, pluriel : inchangés.
\end{itemize}

\all{Ich besuche meinen Vater.} « Je rends visite à mon père. »
\all{Sie liebt ihre Mutter.} « Elle aime sa mère. » (féminin accusatif = nominatif)
\end{propriete}

\begin{attention}[sein ou ihr ? Le piège du possesseur]
En français, « son / sa / ses » s'accordent avec l'objet possédé. En allemand, le choix de la base (\all{sein} vs \all{ihr}) dépend \textbf{uniquement du possesseur} :

\begin{itemize}
\item Possesseur masculin (\all{er}) $\to$ base \all{sein} : \all{Paul liebt seine Mutter.} « Paul aime sa mère. »
\item Possesseur féminin (\all{sie}) $\to$ base \all{ihr} : \all{Marie liebt ihren Vater.} « Marie aime son père. »
\end{itemize}

Ne traduis donc jamais « sa mère » par \all{ihre Mutter} si le possesseur est un homme (c'est \all{seine Mutter}). Et n'oublie pas l'accusatif masculin : \all{ihren Vater}, \all{seinen Bruder}.
\end{attention}

\begin{methode}[Vérifier une phrase allemande en 3 étapes]
\begin{enumerate}
\item \textbf{Souligne le verbe conjugué.} Est-il en 2e position (phrase déclarative) ou en 1re position (question oui/non) ?
\item \textbf{Compte les éléments avant le verbe.} Un seul groupe syntaxique est autorisé. Si tu as mis un complément en tête (\all{Heute}, \all{In Douala}...), vérifie que le sujet est \emph{après} le verbe (inversion).
\item \textbf{Vérifie le possessif.} Qui possède ? (masculin $\to$ \all{sein}, féminin $\to$ \all{ihr}). Le nom possédé est-il COD (accusatif) ? Alors masculin $\to$ \all{-en} (\all{meinen}, \all{deinen}, \all{seinen}, \all{ihren}, \all{unseren}, \all{euren}).
\end{enumerate}
\end{methode}

\begin{exoresolu}[Traduis et justifie]
Énoncé : traduis en allemand.
\begin{enumerate}
\item « Mon frère habite à Douala. »
\item « Aujourd'hui, nous visitons notre grand-mère. »
\item « As-tu des frères et sœurs ? »
\item « Paul aime son père. »
\end{enumerate}
\tcblower
Solution détaillée :

\begin{enumerate}
\item \all{Mein Bruder wohnt in Douala.} — Sujet \all{Mein Bruder} en position 1, verbe \all{wohnt} en position 2. \all{Bruder} masculin, nominatif $\to$ \all{mein}.
\item \all{Heute besuchen wir unsere Großmutter.} — Complément \all{Heute} en position 1 $\to$ verbe \all{besuchen} en position 2, sujet \all{wir} après (inversion). \all{Großmutter} féminin, accusatif = nominatif $\to$ \all{unsere}.
\item \all{Hast du Geschwister?} — Question fermée : verbe \all{hast} en 1re position.
\item \all{Paul liebt seinen Vater.} — Possesseur Paul (masculin) $\to$ base \all{sein}. \all{Vater} COD (accusatif masculin) $\to$ \all{seinen}.
\end{enumerate}
\end{exoresolu}

\begin{exercice}[À toi de jouer]
\begin{enumerate}
\item Remets les mots dans l'ordre correct :
      \begin{enumerate}
      \item \all{in Yaoundé / lebt / Familie / meine}
      \item \all{besucht / heute / mein / uns / Onkel}
      \end{enumerate}
\item Complète avec l'article ou le possessif correct (cas indiqué) :
      \begin{enumerate}
      \item \all{Ich habe \dots\ Bruder.} (ein, accusatif)
      \item \all{Sie besucht \dots\ Tante.} (ihr, accusatif)
      \item \all{Er liebt \dots\ Mutter.} (sein, accusatif)
      \item \all{Wir sehen \dots\ Kind.} (das, accusatif)
      \end{enumerate}
\item Traduis en allemand :
      \begin{enumerate}
      \item « Le week-end, ma sœur rend visite à ses grands-parents. »
      \item « Où habitent tes parents ? »
      \end{enumerate}
\end{enumerate}
\end{exercice}

\begin{corrige}[Exercice « À toi de jouer »]
\begin{enumerate}
\item \all{Meine Familie lebt in Yaoundé.} ; \all{Heute besucht uns mein Onkel.} (inversion après \all{Heute} : verbe 2e position, sujet après).
\item \all{Ich habe einen Bruder.} (acc. masc. \all{einen}) ; \all{Sie besucht ihre Tante.} (acc. fém. inchangé) ; \all{Er liebt seine Mutter.} (acc. fém. inchangé) ; \all{Wir sehen das Kind.} (acc. neutre inchangé).
\item \all{Am Wochenende besucht meine Schwester ihre Großeltern.} (complément en tête $\to$ inversion) ; \all{Wo wohnen deine Eltern?} (mot interrogatif + verbe + sujet).
\end{enumerate}
\end{corrige}

\begin{aretenir}[L'essentiel de la section]
\begin{itemize}
\item Phrase déclarative : le verbe conjugué est \textbf{toujours en 2e position} ; un seul élément (groupe) devant ; complément en tête $\to$ inversion du sujet.
\item Question oui/non : verbe en 1re position. Question avec \all{wo}, \all{wann}, \all{wie}... : mot interrogatif (pos. 1) + verbe (pos. 2).
\item Accusatif : seul le masculin change (\all{der} $\to$ \all{den}, \all{ein} $\to$ \all{einen}, \all{mein} $\to$ \all{meinen}, \all{dein} $\to$ \all{deinen}, \all{sein} $\to$ \all{seinen}, \all{ihr} $\to$ \all{ihren}, \all{unser} $\to$ \all{unseren}, \all{euer} $\to$ \all{euren}).
\item Possessifs : le choix de la base (\all{sein}/\all{ihr}) dépend du possesseur ; la terminaison dépend du nom possédé (genre, nombre, cas).
\end{itemize}
\end{aretenir}

\coursec{Grammaire 2 : conjugaison au présent}

Dans la section précédente, nous avons appris à placer les mots dans la phrase simple : le verbe conjugué occupe toujours la \cle{deuxième place}, et le sujet se décline avec les articles et les possessifs (\all{mein Vater}, \all{ihre Mutter}...). Mais une question reste ouverte : \textbf{quelle forme prend le verbe lui-même ?} En français, nous disons « je parle, tu parles, il parle » ; en allemand aussi, le verbe change selon la personne. C'est ce qu'on appelle la \cle{conjugaison}. Dans cette section, nous allons conjuguer au présent les verbes de notre texte \emph{Familien heute} : \all{leben}, \all{wohnen}, \all{besuchen}, \all{heiraten}, \all{haben} et \all{sein}.

\courssub{Observation : le verbe change selon la personne}

\begin{textebox}[Mini-texte d'observation : Unsere Familie]
\all{Ich wohne in Yaoundé und ich lebe mit meinen Eltern. Mein Bruder wohnt in Douala, er arbeitet dort. Meine Schwester spricht gut Deutsch, sie liest viele Bücher. Am Wochenende fahren wir zu unseren Großeltern. Wir sind eine große Familie und wir haben viele Cousins.}

\emph{Je vis à Yaoundé et je vis avec mes parents. Mon frère habite à Douala, il y travaille. Ma sœur parle bien allemand, elle lit beaucoup de livres. Le week-end, nous allons (en voiture) chez nos grands-parents. Nous sommes une grande famille et nous avons beaucoup de cousins.}
\end{textebox}

Observons les verbes soulignés par la logique des formes : \all{ich wohne} mais \all{er wohnt} ; \all{sie spricht} mais \all{ich spreche} ; \all{wir fahren} mais... que se passe-t-il à \all{du} et \all{er} ? Nous allons répondre point par point.

\courssub{Les verbes faibles (réguliers)}

\begin{definition}[Verbe faible]
Un \cle{verbe faible} est un verbe régulier : son \cle{radical} (la partie devant \all{-en} de l'infinitif) ne change jamais. On ajoute simplement les terminaisons personnelles. Exemple : \all{wohnen} → radical \all{wohn-}.
\end{definition}

\begin{propriete}[Terminaisons du présent des verbes faibles]
On enlève \all{-en} de l'infinitif et on ajoute au radical :

\begin{center}
\begin{tabular}{|l|l|l|}
\hline
\rowcolor{popblueL} Personne & Terminaison & Exemple : wohnen \\ \hline
ich & -e & ich wohne \\ \hline
du & -st & du wohnst \\ \hline
er/sie/es & -t & er wohnt \\ \hline
wir & -en & wir wohnen \\ \hline
ihr & -t & ihr wohnt \\ \hline
sie/Sie & -en & sie/Sie wohnen \\ \hline
\end{tabular}
\end{center}

Moyen mnémotechnique : \textbf{-e, -st, -t, -en, -t, -en}. Remarquez que \all{wir} et \all{sie/Sie} gardent la forme de l'infinitif.
\end{propriete}

\begin{exemplebox}[Exemples traduits]
\all{Ich lebe mit meinen Eltern in Yaoundé.} « Je vis avec mes parents à Yaoundé. »

\all{Du besuchst deine Großeltern am Sonntag.} « Tu rends visite à tes grands-parents le dimanche. »

\all{Mein Onkel heiratet im Juli.} « Mon oncle se marie en juillet. »

\all{Wir wohnen in einer kleinen Wohnung.} « Nous habitons dans un petit appartement. »
\end{exemplebox}

\textbf{Comparaison avec le français :} en français, « je parle / tu parles / il parle » s'entendent pareil mais s'écrivent différemment. En allemand, chaque forme s'entend \emph{et} s'écrit différemment : \all{ich wohne}, \all{du wohnst}, \all{er wohnt}. Il faut donc prononcer les terminaisons clairement.

\courssub{Verbes faibles à particularités orthographiques}

\begin{propriete}[Radicaux en -t, -d : insertion d'un -e-]
Quand le radical se termine en \all{-t} ou \all{-d} (\all{arbeiten}, \all{finden}, \all{warten}), on insère un \textbf{-e-} avant \all{-st} et \all{-t} pour faciliter la prononciation :

\all{arbeiten} → du arbeit\textbf{e}st, er arbeit\textbf{e}t, ihr arbeit\textbf{e}t (mais ich arbeite, wir arbeiten).
\end{propriete}

\begin{propriete}[Verbes en -ern et -eln]
Pour les verbes en \all{-ern} (\all{wandern} « randonner »), à la 1re personne le \all{-e-} du radical peut tomber : \all{ich wandre} (formes courantes aussi : \all{ich wandere}). Pour les verbes en \all{-eln} (\all{sammeln} « collectionner ») : \all{ich sammle}. Les autres personnes restent régulières : \all{du wanderst}, \all{er sammelt}.
\end{propriete}

\courssub{Les verbes forts : le changement de voyelle}

\begin{definition}[Verbe fort]
Un \cle{verbe fort} est un verbe irrégulier dont la \cle{voyelle du radical} change à certaines personnes. Les terminaisons restent celles des verbes faibles.
\end{definition}

\begin{propriete}[Règle d'or du changement de voyelle]
Le changement de voyelle ne concerne \textbf{que la 2e et la 3e personne du singulier} (\all{du} et \all{er/sie/es}). Trois types de changement :

\begin{center}
\begin{tabular}{|l|l|l|}
\hline
\rowcolor{popblueL} Changement & Infinitif & du / er, sie, es \\ \hline
e → i & sprechen, essen & sprichst / spricht ; isst / isst \\ \hline
e → ie & sehen, lesen & siehst / sieht ; liest / liest \\ \hline
a → ä & fahren, schlafen & fährst / fährt ; schläfst / schläft \\ \hline
\end{tabular}
\end{center}

Toutes les autres personnes (\all{ich}, \all{wir}, \all{ihr}, \all{sie/Sie}) gardent la voyelle de l'infinitif : \all{ich spreche}, \all{wir fahren}, \all{ihr lest}.
\end{propriete}

\begin{exemplebox}[Exemples traduits]
\all{Meine Schwester spricht gut Deutsch.} « Ma sœur parle bien allemand. »

\all{Er liest ein Buch über Familien.} « Il lit un livre sur les familles. »

\all{Du siehst deine Cousine am Wochenende.} « Tu vois ta cousine le week-end. »

\all{Mein Vater fährt nach Douala.} « Mon père va (en voiture) à Douala. »

\all{Das Baby schläft viel.} « Le bébé dort beaucoup. »

\all{Du isst mit deiner Familie.} « Tu manges avec ta famille. »
\end{exemplebox}

\begin{attention}[Pièges des francophones]
\begin{itemize}
\item Ne mets \textbf{jamais} la voyelle changée à \all{ich} : on dit \all{ich fahre}, pas \all{ich fähre} !
\item \all{essen} : à \all{du} et \all{er}, le \all{-s-} du radical fusionne avec la terminaison : \all{du isst} (et non \all{issst}).
\item Le tréma de \all{fährt} est obligatoire : \all{fahrt} sans tréma est une faute (cela signifie « trajet »).
\end{itemize}
\end{attention}

\begin{methode}[Conjuguer un verbe fort en 4 étapes]
\begin{enumerate}
\item \textbf{Identifier} la voyelle du radical : \all{sehen} → voyelle \all{e}.
\item \textbf{Vérifier} si elle change (liste à mémoriser) : \all{e → ie} pour \all{sehen}.
\item \textbf{Appliquer} le changement seulement à \all{du} et \all{er/sie/es} : \all{sieh-}.
\item \textbf{Ajouter} les terminaisons normales : \all{du siehst}, \all{er sieht} ; les autres : \all{ich sehe}, \all{wir sehen}, \all{ihr seht}, \all{sie sehen}.
\end{enumerate}
\end{methode}

\courssub{Les verbes sein et haben : les deux irréguliers essentiels}

\begin{propriete}[Conjugaison de sein et haben]
\begin{center}
\begin{tabular}{|l|l|l|}
\hline
\rowcolor{popblueL} Personne & sein (être) & haben (avoir) \\ \hline
ich & bin & habe \\ \hline
du & bist & hast \\ \hline
er/sie/es & ist & hat \\ \hline
wir & sind & haben \\ \hline
ihr & seid & habt \\ \hline
sie/Sie & sind & haben \\ \hline
\end{tabular}
\end{center}

Ces deux verbes sont les plus fréquents de la langue allemande : ils doivent être \textbf{mémorisés par cœur}.
\end{propriete}

\begin{exemplebox}[Emplois de sein et haben]
\all{sein} + état, profession, identité :

\all{Sie ist verheiratet.} « Elle est mariée. »

\all{Mein Vater ist Lehrer.} « Mon père est professeur. »

\all{Wir sind eine Großfamilie.} « Nous sommes une famille élargie. »

\all{haben} + possession, famille :

\all{Wir haben drei Kinder.} « Nous avons trois enfants. »

\all{Ich habe zwei Geschwister.} « J'ai deux frères et sœurs. »

\all{Hast du Geschwister?} « As-tu des frères et sœurs ? »
\end{exemplebox}

\begin{attention}[Erreurs fréquentes]
\begin{itemize}
\item On dit \all{du hast} et non \all{du habst} ; \all{er ist} et non \all{er seist}.
\item Ne pas confondre \all{sie sind} (« ils/elles sont ») et \all{Sie sind} (« vous êtes », formel) : la majuscule change le sens !
\item Ne traduis pas « avoir faim » par \all{haben Hunger}... si, justement ! Mais attention : « j'ai 15 ans » se dit \all{ich bin 15 Jahre alt} (avec \all{sein}), pas avec \all{haben}.
\end{itemize}
\end{attention}

\courssub{Tableau récapitulatif de conjugaison}

\begin{popfigure}
\begin{center}
\begin{tikzpicture}
\matrix[matrix of nodes, nodes={minimum height=0.65cm, align=center, anchor=center}, column sep=0pt, row sep=0pt,
 row 1/.style={nodes={fill=popblue, text=white, font=\bfseries, minimum width=2.1cm}},
 column 1/.style={nodes={fill=popblueL, font=\bfseries, minimum width=2.3cm}},
 every node/.style={draw=popline, minimum width=2.1cm}] {
Infinitiv & ich & du & er/sie/es & wir & ihr & sie/Sie \\
wohnen & wohne & wohnst & wohnt & wohnen & wohnt & wohnen \\
sprechen & spreche & \textcolor{poporange}{\textbf{spri}chst} & \textcolor{poporange}{\textbf{spri}cht} & sprechen & sprecht & sprechen \\
lesen & lese & \textcolor{poppurple}{\textbf{lie}st} & \textcolor{poppurple}{\textbf{lie}st} & lesen & lest & lesen \\
fahren & fahre & \textcolor{popgreen}{\textbf{fä}hrst} & \textcolor{popgreen}{\textbf{fä}hrt} & fahren & fahrt & fahren \\
sein & bin & bist & ist & sind & seid & sind \\
haben & habe & \textcolor{poporange}{\textbf{ha}st} & \textcolor{poporange}{\textbf{ha}t} & haben & habt & haben \\
};
\end{tikzpicture}
\end{center}
\legende{Conjugaison au présent : les formes modifiées (2e et 3e personnes du singulier) sont en couleur.}
\end{popfigure}

\begin{aretenir}[L'essentiel]
\begin{itemize}
\item Verbes faibles : radical + \textbf{-e, -st, -t, -en, -t, -en}.
\item Verbes forts : changement de voyelle \textbf{uniquement} à \all{du} et \all{er/sie/es} : e→i, e→ie, a→ä.
\item \all{sein} et \all{haben} : irréguliers, à apprendre par cœur.
\end{itemize}
\end{aretenir}

\courssub{Texte d'application : Familie Mballa}

\begin{textebox}[Familie Mballa aus Yaoundé]
\all{Familie Mballa wohnt in einem Viertel von Yaoundé. Der Vater, Herr Mballa, arbeitet als Lehrer in einem Gymnasium. Er spricht Deutsch, Französisch und Ewondo. Die Mutter ist Ärztin, sie arbeitet in einem Krankenhaus im Zentrum. Die Familie hat drei Kinder: zwei Söhne und eine Tochter. Der älteste Sohn heißt Junior, er ist siebzehn Jahre alt und liest gern Bücher über Fußball. Die Tochter Sandrine ist fünfzehn, sie spricht sehr gut Deutsch und sieht oft deutsche Filme. Der jüngste Sohn ist erst sechs Jahre alt, er schläft viel und isst gern Bananen. Am Wochenende fährt die ganze Familie zu den Großeltern aufs Land. Die Großeltern leben in einem Dorf in der Nähe von Mbalmayo. Sie haben einen großen Garten mit Mangobäumen und Kakao. Die Kinder besuchen ihre Großeltern gern, denn die Großmutter kocht sehr gut. „Wir sind eine Großfamilie“, sagt Herr Mballa, „und wir haben viel Glück.“ Am Sonntagabend fährt die Familie müde, aber glücklich nach Yaoundé zurück.}

\textbf{Glossaire :} das Viertel, - (le quartier) ; das Gymnasium, -nasien (le lycée) ; der Arzt / die Ärztin, -/-innen (le médecin) ; das Krankenhaus, -häuser (l'hôpital) ; der Sohn, -e (le fils) ; die Tochter, - (la fille) ; aufs Land fahren (aller à la campagne) ; das Dorf, -er (le village) ; der Garten, - (le jardin) ; der Mangobaum, -bäume (le manguier) ; der Kakao (le cacao) ; kochen (cuisiner) ; das Glück (le bonheur, la chance) ; müde (fatigué) ; zurückfahren (rentrer).
\end{textebox}

\textbf{Questions de compréhension :}
\begin{enumerate}
\item \all{Wo wohnt Familie Mballa?}
\item \all{Was ist Herr Mballa von Beruf?}
\item \all{Wie viele Kinder hat die Familie?}
\item \all{Wer spricht sehr gut Deutsch?}
\item \all{Wohin fährt die Familie am Wochenende?}
\end{enumerate}

\courssub{Exercices}

\begin{exoresolu}[Conjugue les verbes entre parenthèses]
Énoncé — Conjugue au présent : 1. Ich (wohnen) in Douala. 2. Du (sprechen) gut Deutsch. 3. Mein Vater (fahren) nach Bafoussam. 4. Wir (haben) zwei Kinder. 5. Ihr (lesen) ein Buch. 6. Sie (sein) verheiratet.
\tcblower
Solution détaillée : 1. \all{Ich wohne in Douala.} (verbe faible, ich + -e). 2. \all{Du sprichst gut Deutsch.} (verbe fort e→i, seulement à du/er). 3. \all{Mein Vater fährt nach Bafoussam.} (verbe fort a→ä, 3e personne). 4. \all{Wir haben zwei Kinder.} (haben : forme de l'infinitif à wir). 5. \all{Ihr lest ein Buch.} (pas de changement à ihr !). 6. \all{Sie ist verheiratet.} (sein, 3e personne : ist).
\end{exoresolu}

\begin{exercice}[À toi de jouer]
1. Conjugue : a) er (arbeiten) ; b) du (essen) ; c) wir (sein) ; d) ihr (schlafen) ; e) sie (haben) ; f) du (sehen).

2. Traduis en allemand : a) « Ma sœur lit un livre. » b) « Nous sommes une famille élargie. » c) « As-tu des frères et sœurs ? » d) « Mon grand-père habite à Mbalmayo. »
\end{exercice}

\begin{corrige}[Exercice « À toi de jouer »]
1. a) \all{er arbeitet} (radical en -t : insertion de -e-) ; b) \all{du isst} ; c) \all{wir sind} ; d) \all{ihr schlaft} (pas de changement à ihr !) ; e) \all{sie haben} ; f) \all{du siehst}.

2. a) \all{Meine Schwester liest ein Buch.} b) \all{Wir sind eine Großfamilie.} c) \all{Hast du Geschwister?} d) \all{Mein Großvater wohnt in Mbalmayo.}
\end{corrige}

\coursec{Commentaire et production guidée}

Vous savez maintenant lire le texte \all{Familien heute}, employer le vocabulaire de la famille et conjuguer correctement les verbes au présent avec le verbe en \cle{deuxième position} (V2). Il reste la dernière étape, la plus importante pour le Bac : \textbf{parler et écrire} sur le texte et sur votre propre famille. Dans cette section, vous apprendrez à commenter un texte en quatre étapes (\cle{Kommentar}), puis à rédiger et à présenter oralement un court paragraphe \all{Meine Familie}.

\courssub{La méthode du commentaire de texte en quatre étapes}

Au lycée camerounais, le commentaire de texte (\all{der Kommentar}) est un exercice régulier : on vous donne un texte allemand et on vous demande de le présenter, de l'analyser et de donner votre avis. La bonne nouvelle : la méthode est toujours la même, quel que soit le texte.

\begin{methode}[Commenter un texte en 4 étapes]
\begin{enumerate}
\item \textbf{Présenter le texte} : de quoi parle-t-il ? Quel est son thème, son type (article, lettre, dialogue) ? Phrase-modèle : \all{Der Text handelt von...} (Le texte traite de...).
\item \textbf{Dégager l'idée principale} : quelle est l'information essentielle, en une seule phrase ? Phrase-modèle : \all{Der Autor sagt, dass...} (L'auteur dit que...).
\item \textbf{Analyser les informations et le vocabulaire clé} : relever les \cle{mots-clés} du texte et expliquer ce qu'ils montrent (ici : \all{die Kernfamilie}, \all{allein erziehende Eltern}, \all{die Patchworkfamilie}).
\item \textbf{Donner son avis et comparer avec son propre contexte} : qu'en pensez-vous ? Et au Cameroun, comment est-ce ? Phrases-modèles : \all{Meiner Meinung nach...} (À mon avis...) et \all{In Kamerun ist es anders: ...} (Au Cameroun, c'est différent : ...).
\end{enumerate}
\end{methode}

\begin{popfigure}
\begin{tikzpicture}[node distance=0.55cm]
\node[draw=popblue, fill=popblueL, rounded corners, align=center, minimum width=6.5cm] (a) {\textbf{1. Présentation}\\ \all{Der Text handelt von...}};
\node[draw=popgreen, fill=popgreenL, rounded corners, align=center, minimum width=6.5cm, below=of a] (b) {\textbf{2. Idée principale}\\ \all{Der Autor sagt, dass...}};
\node[draw=poporange, fill=poporangeL, rounded corners, align=center, minimum width=6.5cm, below=of b] (c) {\textbf{3. Analyse}\\ vocabulaire clé : \all{Kernfamilie, Patchworkfamilie...}};
\node[draw=poppurple, fill=poppurpleL, rounded corners, align=center, minimum width=6.5cm, below=of c] (d) {\textbf{4. Avis personnel / Comparaison}\\ \all{Meiner Meinung nach... In Kamerun...}};
\draw[-{Stealth}, thick, popdark] (a) -- (b);
\draw[-{Stealth}, thick, popdark] (b) -- (c);
\draw[-{Stealth}, thick, popdark] (c) -- (d);
\end{tikzpicture}
\legende{Organigramme du commentaire de texte en quatre étapes}
\end{popfigure}

\courssub{Les phrases-modèles du commentaire (\cle{Redemittel})}

Ces phrases toutes faites (\all{Redemittel}) sont vos outils. Apprenez-les par cœur : elles fonctionnent avec n'importe quel texte.

\begin{center}
\begin{tabularx}{\linewidth}{|X|X|X|}
\hline
\rowcolor{popblueL} Deutsch & Français & Emploi \\
\hline
\all{Der Text handelt von den Familienformen.} & Le texte traite des formes de familles. & présentation \\
\hline
\all{Der Autor sagt, dass es viele Familien gibt.} & L'auteur dit qu'il y a beaucoup de familles. & idée principale \\
\hline
\all{Ein wichtiges Wort ist „Patchworkfamilie“.} & Un mot important est « famille recomposée ». & analyse \\
\hline
\all{Meiner Meinung nach sind alle Familien normal.} & À mon avis, toutes les familles sont normales. & avis personnel \\
\hline
\all{In Kamerun ist es anders: Die Großfamilie ist wichtig.} & Au Cameroun, c'est différent : la famille élargie est importante. & comparaison \\
\hline
\end{tabularx}
\end{center}

\begin{attention}[L'inversion après « Meiner Meinung nach »]
\all{Meiner Meinung nach} occupe la \textbf{première position} de la phrase : le verbe conjugué vient donc juste après (place 2), et le sujet passe derrière le verbe. On dit : \all{Meiner Meinung nach \textbf{sind} alle Familienformen normal.} et jamais \all{Meiner Meinung nach alle Familienformen sind normal.} C'est la même règle qu'avec \all{heute}, \all{in Kamerun}, \all{jetzt} : la place 1 est occupée, le verbe reste en place 2.
\end{attention}

\courssub{Commentaire modèle : observation et analyse}

Lisez d'abord ce commentaire modèle rédigé sur le texte \all{Familien heute}. Repérez les quatre étapes de la méthode.

\begin{textebox}[Kommentar : Familien heute in Deutschland]
Der Text handelt von den Familienformen in Deutschland. Er zeigt, dass es heute viele verschiedene Familien gibt. Die Kernfamilie mit Vater, Mutter und Kindern ist klassisch, aber es gibt auch allein erziehende Eltern und Patchworkfamilien. Die Großeltern wohnen oft nicht in der Familie. In Kamerun ist das anders: Die Großfamilie ist sehr wichtig, und die Großeltern leben oft bei den Kindern. Meiner Meinung nach sind alle Familienformen normal.
\end{textebox}

\textbf{Glossaire :} \all{die Familienform, -en} (la forme de famille) ; \all{verschieden} (différent, varié) ; \all{die Kernfamilie, -n} (la famille nucléaire) ; \all{allein erziehend} (parent isolé, adj.) ; \all{die Patchworkfamilie, -n} (la famille recomposée) ; \all{die Großfamilie, -n} (la famille élargie).

\textbf{Analyse du modèle :}
\begin{itemize}
\item \textbf{Étape 1} (présentation) : \all{Der Text handelt von den Familienformen in Deutschland.}
\item \textbf{Étape 2} (idée principale) : \all{Es gibt heute viele verschiedene Familien.}
\item \textbf{Étape 3} (analyse) : les trois types de familles sont cités — \all{Kernfamilie}, \all{allein erziehende Eltern}, \all{Patchworkfamilien} — avec un détail sur les grands-parents.
\item \textbf{Étape 4} (avis et comparaison) : \all{In Kamerun ist das anders...} puis \all{Meiner Meinung nach sind alle Familienformen normal.}
\end{itemize}

\begin{aretenir}[La structure gagnante]
Un bon commentaire de niveau Première tient en 6 à 8 phrases : 1 phrase de présentation, 1 phrase d'idée principale, 2 à 3 phrases d'analyse avec le vocabulaire du texte, 1 phrase de comparaison avec le Cameroun, 1 phrase d'avis personnel. Le verbe conjugué est toujours en deuxième position.
\end{aretenir}

\begin{exoresolu}[Rédiger un mini-commentaire]
Énoncé : Rédigez un commentaire de 4 phrases sur le texte \all{Familien heute} en suivant les quatre étapes et en utilisant les phrases-modèles.
\tcblower
Solution détaillée :
\begin{enumerate}
\item Présentation : \all{Der Text handelt von den Familien in Deutschland.} (Le texte traite des familles en Allemagne.)
\item Idée principale : \all{Der Autor sagt, dass es heute viele verschiedene Familienformen gibt.} (L'auteur dit qu'il existe aujourd'hui beaucoup de formes de familles différentes.)
\item Analyse : \all{Wichtige Wörter sind „Kernfamilie“, „allein erziehend“ und „Patchworkfamilie“.} (Les mots importants sont « famille nucléaire », « parent isolé » et « famille recomposée ».)
\item Avis et comparaison : \all{Meiner Meinung nach ist die Großfamilie in Kamerun sehr wichtig.} (À mon avis, la famille élargie est très importante au Cameroun.)
\end{enumerate}
Vérification : chaque phrase a son verbe conjugué en deuxième position ; les noms ont une majuscule ; les articles sont corrects (\all{der Text}, \all{die Großfamilie}).
\end{exoresolu}

\courssub{Production écrite guidée : « Meine Familie »}

Après le commentaire du texte, place à votre propre production. L'objectif : un paragraphe de \textbf{5 à 8 phrases} avec au moins 6 mots du lexique de la famille, 2 \cle{possessifs} différents, les verbes \all{sein} et \all{haben}, et le verbe toujours en deuxième position.

\begin{methode}[Rédiger « Meine Familie » en 5 étapes]
\begin{enumerate}
\item \textbf{Présenter la famille} : sa taille et son type. \all{Meine Familie ist eine Großfamilie. Wir sind acht Personen.}
\item \textbf{Nommer les membres avec des possessifs} : \all{Mein Vater heißt Paul und meine Mutter heißt Marie.}
\item \textbf{Dire où elle vit} : \all{Wir wohnen in Yaoundé.}
\item \textbf{Ajouter une particularité} : \all{Ich habe zwei Brüder und eine Schwester. Meine Großeltern wohnen auch bei uns.}
\item \textbf{Conclure par une phrase d'opinion} : \all{Ich liebe meine Familie.}
\end{enumerate}
\end{methode}

\begin{exemplebox}[Modèle rédigé et traduit]
\all{Meine Familie ist eine Großfamilie. Wir sind acht Personen. Mein Vater heißt Paul und meine Mutter heißt Marie. Ich habe zwei Brüder und eine Schwester. Meine Großeltern wohnen auch bei uns. Ich liebe meine Familie.}

« Ma famille est une famille élargie. Nous sommes huit personnes. Mon père s'appelle Paul et ma mère s'appelle Marie. J'ai deux frères et une sœur. Mes grands-parents habitent aussi chez nous. J'aime ma famille. »

\textbf{Points de grammaire à observer :} \all{mein Vater} (masculin, sans -e), \all{meine Mutter / meine Großeltern} (féminin et pluriel, avec -e) ; \all{wir sind} (sein), \all{ich habe} (haben) ; verbe en deuxième position dans chaque phrase.
\end{exemplebox}

\begin{center}
\begin{tabularx}{\linewidth}{|X|X|X|}
\hline
\rowcolor{popblueL} Étape & Phrase-modèle & Outil grammatical \\
\hline
Présentation & \all{Meine Familie ist... Wir sind... Personen.} & \all{sein} au présent \\
\hline
Membres & \all{Mein Vater heißt... Meine Mutter heißt...} & possessifs \all{mein/meine} \\
\hline
Habitat & \all{Wir wohnen in Yaoundé / in Douala.} & verbe faible \all{wohnen} \\
\hline
Particularité & \all{Ich habe... Geschwister.} & \all{haben} au présent \\
\hline
Opinion & \all{Ich liebe meine Familie.} & verbe en 2\textsuperscript{e} position \\
\hline
\end{tabularx}
\end{center}

\begin{attention}[Erreurs fréquentes des francophones]
\begin{itemize}
\item Ne dites pas \all{meine Vater} : \all{der Vater} est masculin, donc \all{\textbf{mein} Vater}. Mais \all{\textbf{meine} Mutter} (féminin) et \all{\textbf{meine} Großeltern} (pluriel).
\item N'oubliez pas la majuscule des noms : \all{die Familie}, \all{der Bruder}, \all{die Schwester}.
\item Faux ami : \all{die Geschwister} signifie « les frères et sœurs » (ensemble), pas « les sœurs ».
\item \all{bei uns wohnen} = « habiter chez nous » : la préposition \all{bei} + datif est obligatoire, on ne peut pas traduire « chez » mot à mot.
\end{itemize}
\end{attention}

\begin{exercice}[À vous d'écrire]
Rédigez votre paragraphe \all{Meine Familie} (5 à 8 phrases). Consignes obligatoires : au moins 6 mots du lexique de la famille, 2 possessifs différents, les verbes \all{sein} et \all{haben}, le verbe conjugué en deuxième position dans chaque phrase.
\end{exercice}

\begin{corrige}[Exercice « Meine Familie »]
Exemple de production attendue (à adapter à votre situation réelle) :

\all{Meine Familie ist eine Kernfamilie. Wir sind fünf Personen. Mein Vater heißt Jean und meine Mutter heißt Chantal. Ich habe einen Bruder und eine Schwester. Wir wohnen in Douala. Meine Schwester ist drei Jahre alt. Ich habe meine Familie sehr gern.}

Grille d'auto-évaluation — cochez chaque point :
\begin{itemize}
\item Le verbe conjugué est-il en deuxième position dans chaque phrase ?
\item Les articles sont-ils corrects (\all{der Vater}, \all{die Mutter}, \all{das Kind}) ?
\item Les possessifs sont-ils accordés (\all{mein} + masculin/neutre, \all{meine} + féminin/pluriel) ?
\item Ai-je utilisé au moins 6 mots du vocabulaire du thème ?
\item Ai-je employé \all{sein} et \all{haben} correctement conjugués ?
\end{itemize}
Si vous cochez les cinq cases, votre paragraphe est réussi.
\end{corrige}

\courssub{Production orale : présenter sa famille à la classe}

La dernière étape est orale : vous présentez votre famille à la classe en \textbf{une minute}, à partir de votre paragraphe écrit, puis vous répondez à deux questions de vos camarades.

\begin{experience}[Présentation orale en 1 minute]
\begin{enumerate}
\item Préparez votre paragraphe \all{Meine Familie} et lisez-le à voix haute deux fois.
\item Présentez-vous devant la classe sans lire mot à mot : regardez vos camarades, parlez lentement et articulez (le \all{ch} de \all{ich} se prononce « r' » doux [ç], le \all{w} se prononce « v » : \all{wir} = « vir »).
\item Deux camarades vous posent ensuite une question chacun, par exemple : \all{Wie heißt dein Vater?} (Comment s'appelle ton père ?), \all{Hast du Geschwister?} (As-tu des frères et sœurs ?), \all{Wo wohnt deine Familie?} (Où habite ta famille ?).
\item Répondez par une phrase complète : \all{Mein Vater heißt Paul.} / \all{Ja, ich habe zwei Brüder.}
\end{enumerate}
\end{experience}

\begin{methode}[Réussir sa prise de parole]
\begin{enumerate}
\item Mémorisez la structure, pas chaque mot : présentation, membres, habitat, particularité, opinion.
\item Commencez toujours par \all{Meine Familie ist...} : cela vous lance.
\item Si vous ne comprenez pas une question, dites : \all{Wie bitte?} (Pardon ?) ou \all{Ich verstehe nicht.} (Je ne comprends pas.)
\item Terminez par votre phrase d'opinion : \all{Ich liebe meine Familie.}
\end{enumerate}
\end{methode}

\begin{savaistu}[La famille au Bac]
Au baccalauréat camerounais, la production écrite sur la famille est un sujet très fréquent : « Décrivez votre famille », « Parlez des types de familles dans votre pays »... Savoir présenter sa famille en 6 à 8 phrases correctes — verbe en deuxième position, possessifs accordés, vocabulaire précis — garantit des points faciles. Entraînez-vous à réciter votre paragraphe \all{Meine Familie} : c'est un investissement pour l'examen !
\end{savaistu}

\newpage

\coursec{Activité d'intégration}

\coursec{Les types de familles}

\courssub{Texte support : Familien heute}

\begin{textebox}[Brief von Lukas — Lettre de Lukas]
Liebe Freunde in Kamerun,

ich heiße Lukas und ich bin sechzehn Jahre alt. Ich zeige euch heute meine Familie. Meine Familie ist nicht groß. Mein Vater heißt Thomas, er ist vierzig Jahre alt, und meine Mutter heißt Sabine. Meine Eltern sind verheiratet. Ich habe eine Schwester. Sie heißt Lena und sie ist dreizehn Jahre alt. Meine Großeltern wohnen in Österreich. Mein Großvater ist siebzig Jahre alt. Wir besuchen unsere Großeltern im Sommer. Meine Tante wohnt in München. Sie ist allein erziehend und hat einen Sohn. Das ist mein Cousin Max.

Und ihr? Wie ist eure Familie? Habt ihr Geschwister? Wo wohnen eure Großeltern?

Viele Grüße\\
Lukas
\end{textebox}

\begin{definition}[Glossaire de la lettre]
\begin{itemize}
\item \all{heißen} (er heißt) : s'appeler (verbe faible, voyelle longue \emph{ie})
\item \all{verheiratet} : marié(e) (participe passé utilisé comme adjectif)
\item \all{die Schwester, die Schwestern} : la sœur
\item \all{die Großeltern} (pluriel) : les grands-parents
\item \all{wohnen} : habiter (verbe faible)
\item \all{besuchen} : rendre visite à (verbe faible, transitif + accusatif)
\item \all{die Tante, die Tanten} : la tante
\item \all{allein erziehend} : qui élève seul(e) son enfant (parent isolé)
\item \all{der Cousin, die Cousins} : le cousin
\item \all{und ihr?} : et vous ? / et toi ? (elliptique pour \all{und wie ist es mit euch?})
\end{itemize}
\end{definition}

\courssub{Compréhension du texte}

\begin{exercice}[Questions de compréhension détaillée]
Lisez la lettre de Lukas et répondez en allemand par des phrases complètes.
\begin{enumerate}
\item \all{Wie heißt Lukas' Schwester?}
\item \all{Wo wohnen seine Großeltern?}
\item \all{Ist seine Tante verheiratet?}
\item \all{Wie alt ist Lukas' Großvater?}
\item \all{Was macht die Familie im Sommer?}
\end{enumerate}
\end{exercice}

\begin{corrige}[Exercice de compréhension]
\begin{enumerate}
\item \all{Lukas' Schwester heißt Lena.}
\item \all{Seine Großeltern wohnen in Österreich.}
\item \all{Nein, sie ist nicht verheiratet. Sie ist allein erziehend.}
\item \all{Sein Großvater ist siebzig Jahre alt.}
\item \all{Im Sommer besuchen sie ihre Großeltern.}
\end{enumerate}
\end{corrige}

\courssub{Lexique thématique : die Familie}

\begin{aretenir}[Le vocabulaire de la famille — Articles, pluriels, genres]
\begin{center}
\begin{tabularx}{\linewidth}{|X|X|X|}\hline
\rowcolor{popblueL} Deutsch (Nom + Article + Pluriel) & Français & Remarque \\ \hline
\all{die Familie, die Familien} & la famille & féminin \\ \hline
\all{der Vater, die Väter} & le père & Umlaut au pluriel \\ \hline
\all{die Mutter, die Mütter} & la mère & Umlaut au pluriel \\ \hline
\all{die Eltern} (pluriel) & les parents & \emph{pluraletantum} \\ \hline
\all{der Bruder, die Brüder} & le frère & Umlaut au pluriel \\ \hline
\all{die Schwester, die Schwestern} & la sœur & \\ \hline
\all{die Geschwister} (pluriel) & les frères et sœurs & \emph{pluraletantum}, neutre \\ \hline
\all{der Großvater, die Großväter} & le grand-père & composé : \all{Groß} + \all{Vater} \\ \hline
\all{die Großmutter, die Großmütter} & la grand-mère & composé : \all{Groß} + \all{Mutter} \\ \hline
\all{die Großeltern} (pluriel) & les grands-parents & \emph{pluraletantum} \\ \hline
\all{der Onkel, die Onkel} & l'oncle & masculin \\ \hline
\all{die Tante, die Tanten} & la tante & féminin \\ \hline
\all{der Cousin, die Cousins} & le cousin & masculin (emprunt fr.) \\ \hline
\all{die Cousine, die Cousinen} & la cousine & féminin \\ \hline
\all{der Neffe, die Neffen} & le neveu & \\ \hline
\all{die Nichte, die Nichten} & la nièce & \\ \hline
\all{verheiratet} & marié(e) & adjectif invariable \\ \hline
\all{ledig} & célibataire & adjectif invariable \\ \hline
\all{geschieden} & divorcé(e) & adjectif invariable \\ \hline
\all{verwitwet} & veuf/veuve & adjectif invariable \\ \hline
\all{allein erziehend} & parent isolé & \all{allein} + part. présent \\ \hline
\end{tabularx}
\end{center}
\end{aretenir}

\begin{exercice}[Entraînement aux articles et aux pluriels]
Complétez avec l'article défini correct (\all{der}, \all{die}, \all{das}) et formez le pluriel.
\begin{enumerate}
\item \all{... Großvater} $\rightarrow$ \all{... Großväter}
\item \all{... Tante} $\rightarrow$ \all{... Tanten}
\item \all{... Geschwister} (déjà pluriel)
\item \all{... Onkel} $\rightarrow$ \all{... Onkel}
\item \all{... Cousine} $\rightarrow$ \all{... Cousinen}
\item \all{... Eltern} (déjà pluriel)
\item \all{... Kind} $\rightarrow$ \all{... Kinder} (neutre, \emph{n}-déclinaison pluriel)
\end{enumerate}
\end{exercice}

\begin{corrige}[Articles et pluriels]
\begin{enumerate}
\item \all{der Großvater} — \all{die Großväter}
\item \all{die Tante} — \all{die Tanten}
\item \all{die Geschwister} (pluriel seulement)
\item \all{der Onkel} — \all{die Onkel}
\item \all{die Cousine} — \all{die Cousinen}
\item \all{die Eltern} (pluriel seulement)
\item \all{das Kind} — \all{die Kinder}
\end{enumerate}
\end{corrige}

\courssub{Grammaire 1 : L'ordre des mots et les articles / possessifs}

\begin{propriete}[Règle fondamentale : Le verbe en deuxième position (V2)]
Dans la phrase affirmative (déclarative), le verbe conjugué occupe \textbf{strictement la deuxième place}. La première place est occupée par un seul constituant (sujet, complément de temps, de lieu, objet, etc.). Si un complément commence la phrase, le sujet se place \emph{après} le verbe (inversion partielle).
\begin{center}
\begin{tabularx}{\linewidth}{|X|X|X|}\hline
\rowcolor{popgreenL} Position 1 (Un seul élément) & Position 2 (Verbe conjugué) & Suite de la phrase \\ \hline
\all{Ich} & \all{heiße} & \all{Lukas.} \\ \hline
\all{Mein Vater} & \all{heißt} & \all{Thomas.} \\ \hline
\all{In München} & \all{wohnt} & \all{meine Tante.} \\ \hline
\all{Im Sommer} & \all{besuchen} & \all{wir unsere Großeltern.} \\ \hline
\all{Meine Großeltern} & \all{wohnen} & \all{in Douala.} \\ \hline
\end{tabularx}
\end{center}
\emph{Attention :} Ne séparez jamais le verbe de sa position 2. \all{*Meine Großeltern in Douala wohnen} est \textbf{faux}.
\end{propriete}

\begin{propriete}[Articles définis et indéfinis au nominatif et accusatif]
Au niveau A1/A2, on maîtrise le nominatif (sujet) et l'accusatif (objet direct). Le genre (m/f/n) ne change pas, seul le masculin change à l'accusatif.
\begin{center}
\begin{tabular}{|l|l|l|l|l|}\hline
\rowcolor{poporangeL} Cas & Masculin & Féminin & Neutre & Pluriel \\ \hline
Nominatif (Sujet) & \all{der Vater} / \all{ein Vater} & \all{die Mutter} / \all{eine Mutter} & \all{das Kind} / \all{ein Kind} & \all{die Eltern} / \all{Eltern} \\ \hline
Accusatif (Objet) & \all{den Vater} / \all{einen Vater} & \all{die Mutter} / \all{eine Mutter} & \all{das Kind} / \all{ein Kind} & \all{die Eltern} / \all{Eltern} \\ \hline
\end{tabular}
\end{center}
\emph{Rappel :} \all{haben}, \all{besuchen}, \all{sehen}, \all{suchen} gouvernent l'accusatif. Ex. \all{Ich habe einen Bruder.} / \all{Ich besuche meine Großeltern.}
\end{propriete}

\begin{propriete}[Les adjectifs possessifs (mein, dein, sein, ihr, unser, euer, ihr/Ihr)]
Ils se déclinent \textbf{exactement comme l'article indéfini} (\all{ein, eine, ein}). Le genre et le nombre sont ceux du nom \emph{possédé} (celui qui suit).
\begin{center}
\begin{tabular}{|l|l|l|l|}\hline
\rowcolor{poppurpleL} & Masculin & Féminin & Neutre / Pluriel \\ \hline
mon / ma / mes & \all{mein Vater} & \all{meine Mutter} & \all{mein Kind} / \all{meine Eltern} \\ \hline
ton / ta / tes & \all{dein Bruder} & \all{deine Schwester} & \all{dein Buch} / \all{deine Geschwister} \\ \hline
son / sa / ses (à lui) & \all{sein Onkel} & \all{seine Tante} & \all{sein Auto} / \all{seine Cousins} \\ \hline
son / sa / ses (à elle) & \all{ihr Vater} & \all{ihre Mutter} & \all{ihr Kind} / \all{ihre Großeltern} \\ \hline
notre / nos & \all{unser Großvater} & \all{unsere Großmutter} & \all{unser Haus} / \all{unsere Onkel} \\ \hline
votre / vos & \all{euer Vater} & \all{eure Mutter} & \all{euer Kind} / \all{eure Tanten} \\ \hline
leur / leurs (formel) & \all{Ihr Name} & \all{Ihre Adresse} & \all{Ihr Kind} / \all{Ihre Eltern} \\ \hline
\end{tabular}
\end{center}
\emph{Note sur \all{euer} :} Au féminin, pluriel et neutre, le \emph{e} de la deuxième syllabe tombe : \all{eure Mutter}, \all{eure Eltern} (prononcé « oyr-e »).
\end{propriete}

\begin{attention}[Pièges fréquents pour francophones]
\begin{itemize}
\item \textbf{Ordre des mots :} Ne calquez pas l'ordre français (Sujet + Verbe + Complément). En allemand, le verbe \textbf{ne bouge pas} de la 2\`eme place. \all{*Meine Mutter kocht gut} (SVO) mais \all{Heute kocht meine Mutter gut} (Adv-V-S-O), jamais \all{*Heute meine Mutter kocht gut}.
\item \textbf{Majuscules :} \textbf{Tous} les noms (substantifs) prennent une majuscule : \all{die Familie}, \all{der Vater}, \all{das Kind}, \all{meine Großeltern}. Les adjectifs possessifs (\all{mein}, \all{dein}) et les pronoms \all{ich}, \all{du}, \all{er}, \all{wir}... restent en minuscule (sauf \all{Sie} formel et \all{Ihr} formel).
\item \textbf{Articles :} Ne dites jamais \all{*Ich habe Bruder}. Dites \all{Ich habe \textbf{einen} Bruder} (accusatif) ou \all{Ich habe \textbf{meinen} Bruder}.
\item \textbf{\all{allein erziehend} :} C'est un adjectif composé. Il ne s'accorde pas comme un participe passé français (\all{*allein erzogene}). Il reste invariable : \all{Sie ist allein erziehend.} / \all{Ein allein erziehender Vater} (déclinaison adjectivale standard si attribut épithète).
\end{itemize}
\end{attention}

\courssub{Grammaire 2 : Conjugaison au présent}

\begin{propriete}[Verbes auxiliaires : \all{sein} (être) et \all{haben} (avoir) — Présent de l'indicatif]
Ces deux verbes sont \textbf{irréguliers} et fondamentaux (auxiliaires du passé, verbes d'état/possession).
\begin{center}
\begin{tabular}{|l|l|l|}\hline
\rowcolor{popgreenL} Personne & \all{sein} & \all{haben} \\ \hline
\all{ich} & \all{bin} & \all{habe} \\ \hline
\all{du} & \all{bist} & \all{hast} \\ \hline
\all{er / sie / es} & \all{ist} & \all{hat} \\ \hline
\all{wir} & \all{sind} & \all{haben} \\ \hline
\all{ihr} & \all{seid} & \all{habt} \\ \hline
\all{sie / Sie} & \all{sind} & \all{haben} \\ \hline
\end{tabular}
\end{center}
\emph{Particularités :} \all{ich bin} (pas *\all{be}), \all{du bist}, \all{er ist} (radical \all{ist-}). \all{ich habe}, \all{du hast} (perte du \emph{b}), \all{er hat} (perte du \emph{b}), \all{wir/sie/Sie haben} (forme de l'infinitif).
\end{propriete}

\begin{propriete}[Verbes faibles (réguliers) — Ex. \all{wohnen}, \all{heißen}, \all{besuchen}, \all{suchen}]
Le radical ne change pas. Les terminaisons sont stables.
\begin{center}
\begin{tabular}{|l|l|}\hline
\rowcolor{popblueL} Personne & \all{wohnen} (habiter) \\ \hline
\all{ich} & \all{wohne} \\ \hline
\all{du} & \all{wohnst} \\ \hline
\all{er/sie/es} & \all{wohnt} \\ \hline
\all{wir} & \all{wohnen} \\ \hline
\all{ihr} & \all{wohnt} \\ \hline
\all{sie/Sie} & \all{wohnen} \\ \hline
\end{tabular}
\quad
\begin{tabular}{|l|l|}\hline
\rowcolor{popblueL} Personne & \all{heißen} (s'appeler) \\ \hline
\all{ich} & \all{heiße} \\ \hline
\all{du} & \all{heißt} \\ \hline
\all{er/sie/es} & \all{heißt} \\ \hline
\all{wir} & \all{heißen} \\ \hline
\all{ihr} & \all{heißt} \\ \hline
\all{sie/Sie} & \all{heißen} \\ \hline
\end{tabular}
\end{center}
\emph{Règle du \ss / ss :} Après voyelle longue (\all{ie}), on écrit \ss (\all{heißt}). Après voyelle courte, on écrit \emph{ss} (\all{müssen} $\rightarrow$ \all{er muss}).
\end{propriete}

\begin{propriete}[Verbes forts (irréguliers) — Alternance vocalique au présent (2\`eme et 3\`eme pers. sg.)]
Seules les personnes \all{du} et \all{er/sie/es} changent de voyelle radicale (souvent \emph{e} $\rightarrow$ \emph{ie} ou \emph{i}, \emph{a} $\rightarrow$ \emph{ä}). Les autres personnes sont régulières.
\begin{center}
\begin{tabular}{|l|l|l|l|}\hline
\rowcolor{poppinkL} Verbe & \all{ich} & \all{du} / \all{er/sie/es} & \all{wir} / \all{sie/Sie} \\ \hline
\all{lesen} (lire) & \all{lese} & \all{liest} / \all{liest} & \all{lesen} \\ \hline
\all{sehen} (voir) & \all{sehe} & \all{siehst} / \all{sieht} & \all{sehen} \\ \hline
\all{geben} (donner) & \all{gebe} & \all{gibst} / \all{gibt} & \all{geben} \\ \hline
\all{fahren} (conduire/aller) & \all{fahre} & \all{fährst} / \all{fährt} & \all{fahren} \\ \hline
\all{schlafen} (dormir) & \all{schlafe} & \all{schläfst} / \all{schläft} & \all{schlafen} \\ \hline
\all{treffen} (rencontrer) & \all{treffe} & \all{triffst} / \all{trifft} & \all{treffen} \\ \hline
\end{tabular}
\end{center}
\emph{Verbes du thème :} \all{wohnen}, \all{heißen}, \all{besuchen}, \all{haben}, \all{sein} sont les plus utilisés ici. \all{lesen}, \all{sehen} utiles pour « lire une lettre », « voir la famille ».
\end{propriete}

\courssub{Production guidée : „Unsere Familien“}

\courssub{Situation de communication}
Le lycée de Yaoundé participe à un échange scolaire avec un \all{Gymnasium} de Munich. Pour préparer la visite des correspondants allemands, la professeure organise une exposition intitulée \all{„Unsere Familien“}. Chaque élève réalise un panneau avec l'arbre généalogique de sa famille (réelle ou imaginaire) et un court texte de présentation. Les correspondants ont déjà envoyé l'exemple de Lukas (voir texte support).

\begin{methode}[Étapes de la production]
Suivez l'ordre ci-dessous. Travaillez d'abord seul, puis en binôme pour l'oral.
\begin{enumerate}
\item \textbf{Arbre généalogique.} Dessinez votre arbre sur une feuille A4 (paysage). Placez-vous au centre (\all{ich}). Ajoutez : parents (\all{Eltern}), frères/sœurs (\all{Geschwister}), grands-parents (\all{Großeltern}), oncles/tantes (\all{Onkel/Tanten}), cousins/cousines (\all{Cousins/Cousinen}). \textbf{Étiquetez chaque personne en allemand} : article + nom + prénom (ex. \all{der Großvater Pierre}, \all{die Tante Marie}).
\item \textbf{Production écrite : „Meine Familie“.} Rédigez un paragraphe de \textbf{6 à 8 phrases}. Utilisez la \textbf{Grille d'auto-évaluation} (ci-dessous) comme guide de rédaction.
\item \textbf{Auto-évaluation.} Relisez votre texte, cochez les cases, corrigez \emph{avant} de rendre la copie.
\item \textbf{Production orale par binômes.} Présentez votre arbre à votre voisin en 4--5 phrases (ne lisez pas, parlez !). Ensuite, posez-vous mutuellement \textbf{deux questions} (ex. \all{Hast du Geschwister?}, \all{Wo wohnen deine Großeltern?}, \all{Wie heißt dein Vater?}).
\item \textbf{Exposition.} Les panneaux les plus soignés (orthographe, présentation, contenu) sont affichés au mur pour la visite des correspondants.
\item \textbf{Variante « Experts » (rapides).} Écrivez deux phrases de comparaison avec la famille de Lukas : \all{Meine Familie ist groß, aber Lukas' Familie ist klein. Ich habe zwei Brüder, aber Lukas hat nur eine Schwester.}
\end{enumerate}
\end{methode}

\begin{methode}[Grille d'auto-évaluation — À cocher avant de rendre]
\begin{enumerate}
\item \square Chaque phrase contient un verbe conjugué en \textbf{deuxième position} (règle V2).
\item \square Tous les noms (substantifs) commencent par une \textbf{majuscule}.
\item \square Chaque nom a son \textbf{article} correct (\all{der}, \all{die}, \all{das}) selon son genre.
\item \square J'ai utilisé au moins \textbf{trois adjectifs possessifs différents} (\all{mein}, \all{meine}, \all{dein}, \all{sein}, \all{ihr}, \all{unser}, \all{eure}...), correctement accordés.
\item \square J'ai conjugué correctement \all{sein} (\all{bin}, \all{ist}, \all{sind}...) et \all{haben} (\all{habe}, \all{hat}, \all{haben}...).
\item \square J'ai employé au moins \textbf{cinq mots} du vocabulaire de la famille (ex. \all{Geschwister}, \all{Großeltern}, \all{verheiratet}, \all{allein erziehend}, \all{Cousin}...).
\item \square J'ai respecté l'accord des participes/adjectifs si nécessaire (\all{verheiratet}, \all{allein erziehend} invariables en prédicat).
\end{enumerate}
\end{methode}

\begin{exoresolu}[Meine Familie — Production modèle commentée]
Rédigez votre paragraphe (6--8 phrases) en vous inspirant du modèle ci-dessous.
\tcblower
\textbf{Texte modèle (élève de Première A, Yaoundé) :}

\all{Ich heiße Amina und ich bin sechzehn Jahre alt. Meine Familie ist groß. Mein Vater heißt Ibrahim und meine Mutter heißt Chantal. Meine Eltern sind verheiratet. Ich habe zwei Brüder und eine Schwester. Meine Großeltern wohnen in Douala. Mein Großvater ist fünfundsiebzig Jahre alt. Meine Tante wohnt in Yaoundé. Sie ist allein erziehend und hat eine Tochter. Wir besuchen unsere Großeltern im Dezember.}

\textit{« Je m'appelle Amina et j'ai seize ans. Ma famille est grande. Mon père s'appelle Ibrahim et ma mère s'appelle Chantal. Mes parents sont mariés. J'ai deux frères et une sœur. Mes grands-parents habitent à Douala. Mon grand-père a soixante-quinze ans. Ma tante habite à Yaoundé. Elle élève seule sa fille. Nous rendons visite à nos grands-parents en décembre. »}

\textbf{Commentaire linguistique détaillé :}
\begin{enumerate}
\item \all{Ich heiße Amina} : Verbe fort \all{heißen} (3\`eme pers. sg. \all{heißt}), verbe en position 2. \all{und ich bin} : coordination, \all{bin} (sein, 1\`ere pers. sg.) en position 2 de la proposition.
\item \all{Meine Familie ist groß} : Sujet \all{Meine Familie} (fém. sg.), verbe \all{ist} (sein, 3\`eme pers. sg.), adjectif \all{groß} (invariable en prédicat).
\item \all{Mein Vater heißt Ibrahim} : Possessif \all{Mein} (masc. sg., comme \all{ein}), verbe \all{heißt} (V2). \all{und meine Mutter heißt Chantal} : \all{meine} (fém. sg.).
\item \all{Meine Eltern sind verheiratet} : \all{Eltern} pluriel $\rightarrow$ possessif \all{Meine}, verbe \all{sind} (sein, pluriel), \all{verheiratet} adjectif invariant.
\item \all{Ich habe zwei Brüder und eine Schwester} : \all{habe} (haben, 1\`ere pers. sg.), objet accusatif \all{zwei Brüder} (masc. plur., pas d'article indéfini au pluriel après nombre), \all{eine Schwester} (fém. sg. accusatif = nominatif).
\item \all{Meine Großeltern wohnen in Douala} : Sujet pluriel \all{Meine Großeltern}, verbe \all{wohnen} (faible, 3\`eme pers. pl. = infinitif), complément de lieu \all{in Douala} (préposition \all{in} + accusatif pour le lieu « où » ? Non, \all{wohnen} + \all{in} + \textbf{Datif} pour le lieu « où » $\rightarrow$ \all{in Douala} (nom propre, pas d'article visible, mais Datif sous-jacent). \emph{Note :} Au niveau A1, on accepte \all{in Douala} sans marquer le cas visible).
\item \all{Mein Großvater ist fünfundsiebzig Jahre alt} : Nombre \all{fünfundsiebzig} (57), construction d'âge \all{ist ... Jahre alt}.
\item \all{Meine Tante wohnt in Yaoundé. Sie ist allein erziehend und hat eine Tochter.} : Enchaînement cohérent. \all{Sie} (pronom 3\`eme fém. sg. pour \all{die Tante}). \all{allein erziehend} adjectif composé invariant. \all{hat} (haben, 3\`eme sg.), \all{eine Tochter} (accusatif fém. sg.).
\item \all{Wir besuchen unsere Großeltern im Dezember.} : Sujet \all{Wir}, verbe \all{besuchen} (faible, 1\`ere pl.), objet \all{unsere Großeltern} (accusatif pluriel = nominatif), complément de temps \all{im Dezember} (\all{in dem}).
\end{enumerate}
\end{exoresolu}

\begin{exercice}[Pour aller plus loin — Interaction épistolaire]
Imaginez que vous répondez à Lukas.
\begin{enumerate}
\item Écrivez \textbf{deux questions} personnelles à lui poser (ex. sur ses loisirs, son école, son cousin...).
\item Répondez à \textbf{ses trois questions} de la fin de sa lettre : \all{Wie ist eure Familie?} \quad \all{Habt ihr Geschwister?} \quad \all{Wo wohnen eure Großeltern?}
\end{enumerate}
\end{exercice}

\begin{corrige}[Exercice « Pour aller plus loin »]
\textbf{Questions possibles pour Lukas :}
\begin{itemize}
\item \all{Wie alt ist dein Cousin Max?} (Quel âge a ton cousin Max ?)
\item \all{Was sind deine Hobbys?} (Quels sont tes loisirs ?)
\item \all{Gehst du in München zur Schule?} (Vas-tu à l'école à Munich ?)
\end{itemize}

\textbf{Réponses modèles aux questions de Lukas :}
\begin{enumerate}
\item \all{Meine Familie ist groß. Wir sind sechs Personen.} (Ma famille est grande. Nous sommes six personnes.)
\item \all{Ja, ich habe Geschwister: zwei Brüder und eine Schwester.} (Oui, j'ai des frères et sœurs : deux frères et une sœur.) \emph{Attention :} Question \all{Habt ihr...?} (verbe position 1) $\rightarrow$ Réponse \all{Ich habe...} (verbe position 2).
\item \all{Meine Großeltern wohnen in Douala.} (Mes grands-parents habitent à Douala.) \emph{Ordre V2 respecté.}
\end{enumerate}
\end{corrige}

\courssub{Bilan de la leçon}

Cette leçon a permis de valider les compétences du programme officiel (MINESEC, Allemand LV2, Première A) pour le Chapitre 1 « Vie familiale et sociale » :

\begin{itemize}
\item \textbf{Compréhension de l'écrit :} Lecture globale et détaillée d'une lettre authentique (niveau A1+/A2), extraction d'informations précises, réponse par phrases complètes.
\item \textbf{Lexique :} Maîtrise active du vocabulaire de la famille (noms, articles, pluriels, Umlautes), statuts matrimoniaux (\all{verheiratet}, \all{allein erziehend}), prépositions de lieu (\all{in} + Datif/Ville, \all{in} + \all{dem} = \all{im}).
\item \textbf{Grammaire :} Application rigoureuse de la règle V2 (ordre des mots), déclinaison des articles définis/indéfinis (Nominatif/Accusatif), déclinaison des possessifs (comme \all{ein}).
\item \textbf{Conjugaison :} Présent de l'indicatif des verbes irréguliers fondamentaux (\all{sein}, \all{haben}), verbes faibles (\all{wohnen}, \all{heißen}, \all{besuchen}), introduction aux verbes forts à alternance vocalique (\all{lesen}, \all{sehen}, \all{fahren}).
\item \textbf{Expression écrite :} Production d'un paragraphe cohérent (6--8 phrases), auto-évaluation guidée par critères explicites (méthodologie).
\item \textbf{Expression orale :} Interaction en binôme (présentation + questions/réponses), préparation à l'échange réel avec correspondants.
\item \textbf{Médiation / Culture :} Comparaison structures familiales Cameroun/Allemagne (famille nucléaire vs élargie), valorisation du contexte local (Yaoundé, Douala, décembre/saison sèche).
\end{itemize}

\begin{savaistu}[Les familles en Allemagne et au Cameroun : regards croisés]
En Allemagne (\all{Deutschland}), en Autriche (\all{Österreich}) et en Suisse (\all{Schweiz}), le modèle dominant est la \textbf{famille nucléaire} (\all{die Kernfamilie}) : parents et enfants (souvent 1 ou 2). Les grands-parents (\all{Großeltern}) habitent souvent dans une autre ville, voire un autre \all{Bundesland} (Land). Le terme \all{die Großfamilie} (famille élargie) désigne explicitement le cercle large (oncles, tantes, cousins), qui ne vit pas sous le même toit.

Au Cameroun, la \textbf{famille élargie} est la norme vécue au quotidien : oncles (\all{Onkel}), tantes (\all{Tanten}), cousins (\all{Cousins/Cousinen}) partagent souvent le même quartier, les mêmes repas, les mêmes devoirs de solidarité. Votre arbre généalogique camerounais sera donc souvent bien plus fourni que celui de Lukas — c'est une richesse culturelle à expliquer à vos correspondants avec fierté !

\textbf{Vocabulaire bonus pour l'échange :}
\begin{itemize}
\item \all{die Kernfamilie} : famille nucléaire (père, mère, enfants)
\item \all{die Großfamilie} : famille élargie
\item \all{das Familienfest} : fête de famille
\item \all{der Stammbaum} : arbre généalogique
\item \all{verwandt sein} : être parent (avec qqn) — \all{Wir sind verwandt.} (Nous sommes parents.)
\end{itemize}
\end{savaistu}

\newpage

\coursec{Méthodes et exercices résolus}

\courssub{Méthode 1 : Comprendre un texte sur les types de familles}

\begin{methode}[Lire et exploiter un texte sur la famille]
\begin{enumerate}
\item \textbf{Lecture globale.} Lisez le texte une première fois sans dictionnaire. Repérez le thème (ici : \all{Familien heute}), le type de texte (article, portrait, dialogue) et les personnages. Soulignez les mots de la famille lexicale : \all{die Familie, die Eltern, die Großeltern, die Geschwister, verheiratet, allein erziehend}.
\item \textbf{Repérage des mots transparents.} Beaucoup de mots allemands ressemblent au français ou à l'anglais : \all{die Familie, die Mutter, der Bruder, modern, traditionell}. Ils donnent le sens général.
\item \textbf{Lecture détaillée.} Pour chaque phrase, repérez d'abord le \cle{verbe conjugué} (toujours en 2\textsuperscript{e} position dans la phrase déclarative), puis le sujet, puis les compléments.
\item \textbf{Réponse aux questions.} Reformulez la question en phrase complète : la question \all{Wo wohnt Familie Mbarga?} appelle la réponse \all{Familie Mbarga wohnt in Yaoundé.} Ne répondez jamais par un seul mot.
\item \textbf{Résumé.} Gardez uniquement : qui ? où ? quelle situation familiale ? 2 ou 3 phrases maximum, avec vos propres mots.
\end{enumerate}
\end{methode}

\begin{exoresolu}[Compréhension de texte : Familie Mbarga]
\textbf{Énoncé.} Lisez le texte et répondez en allemand par des phrases complètes.

\all{Familie Mbarga wohnt in Yaoundé. Herr und Frau Mbarga sind verheiratet und haben drei Kinder: zwei Söhne und eine Tochter. Die Großeltern leben auch im Haus. Sie helfen den Kindern bei den Hausaufgaben. In Douala wohnt Frau Ngo Bell. Sie ist allein erziehend und hat eine kleine Tochter. Sie arbeitet als Lehrerin.}

Questions : 1) \all{Wo wohnt Familie Mbarga?} 2) \all{Wie viele Kinder haben Herr und Frau Mbarga?} 3) \all{Was machen die Großeltern?} 4) \all{Ist Frau Ngo Bell verheiratet?}
\tcblower
\textbf{Raisonnement.} 1) La question porte sur le lieu (\all{wo}) : on cherche le lieu dans la première phrase. 2) \all{Wie viele} demande un nombre : le texte dit \all{drei Kinder}. 3) \all{Was machen} demande une action : \all{sie helfen den Kindern}. 4) Question fermée : le texte dit \all{allein erziehend}, donc la réponse est négative.

\textbf{Réponses rédigées.}
\begin{enumerate}
\item \all{Familie Mbarga wohnt in Yaoundé.} « La famille Mbarga habite à Yaoundé. »
\item \all{Herr und Frau Mbarga haben drei Kinder.} « M. et M\textsuperscript{me} Mbarga ont trois enfants. »
\item \all{Die Großeltern helfen den Kindern bei den Hausaufgaben.} « Les grands-parents aident les enfants à faire leurs devoirs. »
\item \all{Nein, Frau Ngo Bell ist nicht verheiratet. Sie ist allein erziehend.} « Non, M\textsuperscript{me} Ngo Bell n'est pas mariée. Elle élève seule son enfant. »
\end{enumerate}
\textbf{Conclusion.} Chaque réponse reprend la structure de la question, avec le verbe conjugué en 2\textsuperscript{e} position et une majuscule à tous les noms.
\end{exoresolu}

\courssub{Méthode 2 : Employer le vocabulaire de la famille}

\begin{methode}[Apprendre et réutiliser le lexique thématique]
\begin{enumerate}
\item \textbf{Apprenez chaque nom avec son article et son pluriel} : \all{die Familie, -n ; der Vater, die Väter ; die Mutter, die Mütter ; das Kind, -er ; die Geschwister} (pluriel seulement).
\item \textbf{Classez par familles de mots} : \all{heiraten} (verbe) $\rightarrow$ \all{verheiratet} (adjectif) ; \all{erziehen} $\rightarrow$ \all{allein erziehend}.
\item \textbf{Mémorisez les expressions clés} : \all{eine große Familie haben} (avoir une famille nombreuse), \all{bei den Eltern wohnen} (habiter chez ses parents), \all{ein Einzelkind sein} (être enfant unique).
\item \textbf{Réutilisez le lexique dans vos propres phrases} : présentez votre famille en 3 ou 4 phrases en piochant dans le tableau.
\end{enumerate}
\end{methode}

\begin{definition}[Le vocabulaire de la famille : article, pluriel, traduction]
\begin{center}
\begin{tabularx}{\linewidth}{|l|l|X|}
\hline
\rowcolor{popblueL} Deutsch & Plural & Français \\
\hline
\all{die Familie} & \all{-n} & la famille \\
\hline
\all{der Vater} & \all{die Väter} & le père \\
\hline
\all{die Mutter} & \all{die Mütter} & la mère \\
\hline
\all{die Eltern} & (pluriel seulement) & les parents \\
\hline
\all{die Großeltern} & (pluriel seulement) & les grands-parents \\
\hline
\all{der Großvater / der Opa} & \all{die Großväter / die Opas} & le grand-père / le papi \\
\hline
\all{die Großmutter / die Oma} & \all{die Großmütter / die Omas} & la grand-mère / la mamie \\
\hline
\all{das Kind} & \all{-er} & l'enfant \\
\hline
\all{der Sohn} & \all{die Söhne} & le fils \\
\hline
\all{die Tochter} & \all{die Töchter} & la fille \\
\hline
\all{der Bruder} & \all{die Brüder} & le frère \\
\hline
\all{die Schwester} & \all{-n} & la sœur \\
\hline
\all{die Geschwister} & (pluriel seulement) & les frères et sœurs \\
\hline
\all{das Einzelkind} & \all{-er} & l'enfant unique \\
\hline
\all{der Onkel} & \all{-} & l'oncle \\
\hline
\all{die Tante} & \all{-n} & la tante \\
\hline
\all{verheiratet} & --- & marié(e) \\
\hline
\all{allein erziehend} & --- & qui élève seul(e) son enfant \\
\hline
\all{geschieden} & --- & divorcé(e) \\
\hline
\all{ledig} & --- & célibataire \\
\hline
\end{tabularx}
\end{center}
\end{definition}

\begin{exoresolu}[Vocabulaire : compléter un portrait de famille]
\textbf{Énoncé.} Complétez avec le mot qui convient : \all{Großeltern, Geschwister, verheiratet, allein erziehend, Kinder}.

\all{Meine Eltern sind seit zwanzig Jahren \dots\ Ich habe zwei \dots\ : einen Bruder und eine Schwester. Meine \dots\ wohnen im Dorf bei Bafoussam. Meine Tante ist \dots\ und hat zwei \dots\ .}
\tcblower
\textbf{Raisonnement.} 1) « mariés depuis vingt ans » $\rightarrow$ \all{verheiratet}. 2) « un frère et une sœur » = frères et sœurs $\rightarrow$ \all{Geschwister}. 3) « habitent au village » : personnes de la génération précédente $\rightarrow$ \all{Großeltern}. 4) La tante élève seule $\rightarrow$ \all{allein erziehend}. 5) Elle a deux enfants $\rightarrow$ \all{Kinder}.

\textbf{Texte complété.} \all{Meine Eltern sind seit zwanzig Jahren verheiratet. Ich habe zwei Geschwister: einen Bruder und eine Schwester. Meine Großeltern wohnen im Dorf bei Bafoussam. Meine Tante ist allein erziehend und hat zwei Kinder.}

« Mes parents sont mariés depuis vingt ans. J'ai deux frères et sœurs : un frère et une sœur. Mes grands-parents habitent au village près de Bafoussam. Ma tante élève seule ses deux enfants. »
\end{exoresolu}

\courssub{Méthode 3 : Construire une phrase simple allemande (ordre des mots)}

\begin{methode}[Placer le verbe en 2\textsuperscript{e} position]
\begin{enumerate}
\item \textbf{Règle d'or} : dans la phrase déclarative, le \cle{verbe conjugué} occupe toujours la \textbf{2\textsuperscript{e} position}. Une « position » = un groupe de mots (sujet, complément de temps, de lieu).
\item \textbf{Ordre de base} : sujet + verbe + compléments. \all{Die Familie wohnt in Yaoundé.}
\item \textbf{Inversion} : si un complément ouvre la phrase, le verbe reste 2\textsuperscript{e} et le sujet passe après : \all{In Yaoundé wohnt die Familie.} (et jamais \all{In Yaoundé die Familie wohnt}).
\item \textbf{Ordre des compléments} : temps avant lieu : \all{Die Großeltern kommen am Sonntag nach Yaoundé.}
\item \textbf{Vérifiez toujours} : comptez les positions (1 = premier groupe, 2 = verbe) avant d'écrire.
\end{enumerate}
\end{methode}

\begin{propriete}[La phrase déclarative allemande : verbe en 2\textsuperscript{e} position]
Dans une phrase déclarative affirmative, le verbe conjugué est toujours le \textbf{deuxième élément} de la phrase, quelle que soit la longueur du premier élément.

\begin{center}
\begin{tabularx}{\linewidth}{|X|l|X|X|}
\hline
\rowcolor{popblueL} Position 1 & Position 2 (verbe) & Suite & Traduction \\
\hline
\all{Die Familie} & \all{wohnt} & \all{in Yaoundé.} & La famille habite à Yaoundé. \\
\hline
\all{In Yaoundé} & \all{wohnt} & \all{die Familie.} & À Yaoundé habite la famille. \\
\hline
\all{Am Sonntag} & \all{besuchen} & \all{wir die Großeltern.} & Le dimanche, nous rendons visite aux grands-parents. \\
\hline
\all{Meine Eltern} & \all{haben} & \all{vier Kinder.} & Mes parents ont quatre enfants. \\
\hline
\end{tabularx}
\end{center}

\textbf{Exception à connaître} : dans la question sans mot interrogatif, le verbe est en \textbf{1\textsuperscript{re} position} : \all{Wohnst du in Douala?} « Habites-tu à Douala ? » ; dans la question avec mot interrogatif, le verbe est 2\textsuperscript{e} : \all{Wo wohnst du?}
\end{propriete}

\begin{exoresolu}[Thème : traduire en respectant l'ordre des mots]
\textbf{Énoncé.} Traduisez en allemand : 1) « Ma sœur habite à Douala. » 2) « Le dimanche, les grands-parents visitent la famille. » 3) « Mon frère a deux enfants. »
\tcblower
\textbf{Raisonnement.} 1) Sujet \all{meine Schwester} en 1\textsuperscript{re} position, verbe \all{wohnt} en 2\textsuperscript{e}, lieu \all{in Douala}. 2) La phrase commence par le temps « le dimanche » = \all{am Sonntag} : le verbe \all{besuchen} doit rester 2\textsuperscript{e}, donc inversion du sujet. 3) Verbe \all{haben} : 3\textsuperscript{e} personne du singulier \all{hat}.

\textbf{Traductions.}
\begin{enumerate}
\item \all{Meine Schwester wohnt in Douala.}
\item \all{Am Sonntag besuchen die Großeltern die Familie.}
\item \all{Mein Bruder hat zwei Kinder.}
\end{enumerate}
\textbf{Conclusion.} La phrase 2 illustre l'inversion : après un complément initial, le sujet suit immédiatement le verbe. C'est le point le plus sanctionné au Bac.
\end{exoresolu}

\courssub{Méthode 4 : Choisir l'article et le possessif corrects}

\begin{methode}[Articles définis et adjectifs possessifs]
\begin{enumerate}
\item \textbf{Apprenez le genre avec le nom} : \all{der Vater} (masculin), \all{die Mutter} (féminin), \all{das Kind} (neutre). Ne devinez jamais le genre d'après le français : « la famille » se dit bien \all{die Familie}, mais « le problème » se dit \all{das Problem}.
\item \textbf{Tableau des possessifs au nominatif} : \all{mein Vater, meine Mutter, mein Kind} ; \all{sein/ihr Bruder, seine/ihre Schwester, sein/ihr Kind}.
\item \textbf{Règle} : le possessif prend la terminaison de l'article indéfini : masculin et neutre sans terminaison au nominatif (\all{mein Vater, mein Kind}), féminin et pluriel avec \textbf{-e} (\all{meine Mutter, meine Kinder}).
\item \textbf{Attention à l'accusatif masculin} : \all{Ich besuche meinen Vater.} (le possessif prend \textbf{-en}).
\end{enumerate}
\end{methode}

\begin{propriete}[Les adjectifs possessifs : bases et terminaisons]
\textbf{Bases} : \all{mein} (mon), \all{dein} (ton), \all{sein} (son, pour un possesseur masculin ou neutre), \all{ihr} (son/sa, pour une possesseure féminine ; leur), \all{unser} (notre), \all{euer} (votre), \all{Ihr} (votre, politesse).

\textbf{Terminaisons au nominatif et à l'accusatif} (comme l'article indéfini \all{ein}) :

\begin{center}
\begin{tabular}{|l|l|l|l|}
\hline
\rowcolor{popblueL} & Masculin (\all{der Vater}) & Féminin (\all{die Mutter}) & Neutre (\all{das Kind}) \\
\hline
Nominatif & \all{mein Vater} & \all{meine Mutter} & \all{mein Kind} \\
\hline
Accusatif & \all{meinen Vater} & \all{meine Mutter} & \all{mein Kind} \\
\hline
Datif & \all{meinem Vater} & \all{meiner Mutter} & \all{meinem Kind} \\
\hline
\end{tabular}
\end{center}

Au pluriel : nominatif et accusatif \all{meine Kinder}, datif \all{meinen Kindern} (le nom prend aussi \textbf{-n} au datif pluriel).
\end{propriete}

\begin{exoresolu}[Grammaire : articles et possessifs]
\textbf{Énoncé.} Complétez avec le possessif correct (\all{mein-, sein-, ihr-}) : 1) \all{Paul liebt \dots\ Familie.} 2) \all{\dots\ Mutter kocht heute.} (je) 3) \all{Anna besucht \dots\ Großeltern.} 4) \all{Wir spielen mit \dots\ Kind.} (il)
\tcblower
\textbf{Raisonnement.} 1) Paul = il $\rightarrow$ \all{sein-} ; \all{die Familie} féminin accusatif $\rightarrow$ \all{seine}. 2) \all{die Mutter} féminin nominatif $\rightarrow$ \all{meine}. 3) Anna = elle $\rightarrow$ \all{ihr-} ; \all{die Großeltern} pluriel accusatif $\rightarrow$ \all{ihre}. 4) \all{das Kind} neutre, après \all{mit} (datif) $\rightarrow$ \all{seinem}.

\textbf{Solutions.} 1) \all{Paul liebt seine Familie.} 2) \all{Meine Mutter kocht heute.} 3) \all{Anna besucht ihre Großeltern.} 4) \all{Wir spielen mit seinem Kind.}

\textbf{Conclusion.} Trois questions à se poser : qui possède (choix de la base \all{mein/sein/ihr}) ? quel est le genre du nom possédé ? quel est son cas dans la phrase ?
\end{exoresolu}

\courssub{Méthode 5 : Conjuguer au présent (verbes faibles, forts, sein, haben)}

\begin{methode}[Le présent de l'indicatif]
\begin{enumerate}
\item \textbf{Verbes faibles} (réguliers) : radical + terminaisons \textbf{-e, -st, -t, -en, -t, -en}. \all{wohnen : ich wohne, du wohnst, er wohnt, wir wohnen, ihr wohnt, sie wohnen.}
\item \textbf{Verbes forts} : même schéma, mais souvent changement de voyelle à la 2\textsuperscript{e} et 3\textsuperscript{e} personne du singulier : \all{sprechen : du sprichst, er spricht ; lesen : du liest, er liest ; fahren : du fährst, er fährt.}
\item \textbf{sein} (irrégulier) : \all{ich bin, du bist, er/sie/es ist, wir sind, ihr seid, sie/Sie sind.}
\item \textbf{haben} : \all{ich habe, du hast, er/sie/es hat, wir haben, ihr habt, sie/Sie haben.}
\item \textbf{Réflexe} : identifiez d'abord le sujet, puis vérifiez si le verbe est fort (voyelle qui change) ou irrégulier.
\end{enumerate}
\end{methode}

\begin{propriete}[Tableaux de conjugaison au présent]
\textbf{Verbe faible} \all{wohnen} (habiter) et \textbf{verbe fort} \all{sprechen} (parler, e $\rightarrow$ i) :

\begin{center}
\begin{tabular}{|l|l|l|}
\hline
\rowcolor{popblueL} Personne & \all{wohnen} (faible) & \all{sprechen} (fort) \\
\hline
\all{ich} & \all{wohne} & \all{spreche} \\
\hline
\all{du} & \all{wohnst} & \all{sprichst} \\
\hline
\all{er/sie/es} & \all{wohnt} & \all{spricht} \\
\hline
\all{wir} & \all{wohnen} & \all{sprechen} \\
\hline
\all{ihr} & \all{wohnt} & \all{sprecht} \\
\hline
\all{sie/Sie} & \all{wohnen} & \all{sprechen} \\
\hline
\end{tabular}
\end{center}

\textbf{Verbes irréguliers} \all{sein} (être) et \all{haben} (avoir) :

\begin{center}
\begin{tabular}{|l|l|l|}
\hline
\rowcolor{popblueL} Personne & \all{sein} & \all{haben} \\
\hline
\all{ich} & \all{bin} & \all{habe} \\
\hline
\all{du} & \all{bist} & \all{hast} \\
\hline
\all{er/sie/es} & \all{ist} & \all{hat} \\
\hline
\all{wir} & \all{sind} & \all{haben} \\
\hline
\all{ihr} & \all{seid} & \all{habt} \\
\hline
\all{sie/Sie} & \all{sind} & \all{haben} \\
\hline
\end{tabular}
\end{center}

\textbf{Changements de voyelle fréquents des verbes forts} (2\textsuperscript{e} et 3\textsuperscript{e} personnes du singulier uniquement) : e $\rightarrow$ i (\all{sprechen, helfen : er hilft}), e $\rightarrow$ ie (\all{lesen, sehen : er sieht}), a $\rightarrow$ ä (\all{fahren, schlafen : er schläft}).
\end{propriete}

\begin{exoresolu}[Conjugaison : verbes au présent]
\textbf{Énoncé.} Conjuguez au présent : 1) \all{ich (wohnen)} 2) \all{er (haben)} 3) \all{wir (sein)} 4) \all{du (sprechen)} 5) \all{sie (lesen)} 6) \all{ihr (kommen)}
\tcblower
\textbf{Raisonnement.} 1) Verbe faible : \all{ich wohne}. 2) \all{haben} irrégulier : \all{er hat} (pas de \all{habt}). 3) \all{sein} : \all{wir sind}. 4) Verbe fort, changement e $\rightarrow$ i : \all{du sprichst}. 5) Verbe fort, e $\rightarrow$ ie : \all{sie liest}. 6) Verbe faible : \all{ihr kommt}.

\textbf{Solutions.} 1) \all{ich wohne} 2) \all{er hat} 3) \all{wir sind} 4) \all{du sprichst} 5) \all{sie liest} 6) \all{ihr kommt}

\textbf{Phrase d'application.} \all{Meine Familie ist groß. Wir haben vier Kinder und wir wohnen in Douala.} « Ma famille est nombreuse. Nous avons quatre enfants et nous habitons à Douala. »
\end{exoresolu}

\courssub{Méthode 6 : Produire un court paragraphe sur sa famille}

\begin{methode}[Rédiger une présentation de famille (5 à 8 lignes)]
\begin{enumerate}
\item \textbf{Plan en 4 points} : présentation générale (\all{Meine Familie ist groß/klein}), les parents (nom, métier), les frères et sœurs, une particularité (grands-parents, parent isolé, famille recomposée).
\item \textbf{Structures à réutiliser} : \all{Ich habe \dots\ Geschwister. / Mein Vater arbeitet als \dots\ / Meine Eltern sind verheiratet. / Meine Großeltern wohnen in \dots}
\item \textbf{Variez les débuts de phrase} (inversion) : \all{In Yaoundé wohnt meine Familie. Am Wochenende besuchen wir die Großeltern.}
\item \textbf{Relisez} : verbe en 2\textsuperscript{e} position, majuscule aux noms, Umlaute, accord du possessif.
\end{enumerate}
\end{methode}

\begin{exoresolu}[Production écrite : ma famille]
\textbf{Énoncé.} Rédigez un texte de 5 à 6 phrases pour présenter votre famille (type de famille, membres, lieu d'habitation).
\tcblower
\textbf{Démarche.} On suit le plan : 1) phrase d'ouverture, 2) les parents, 3) les frères et sœurs, 4) les grands-parents, 5) phrase de conclusion. On vérifie chaque verbe (2\textsuperscript{e} position) et chaque majuscule.

\textbf{Production modèle.}

\all{Meine Familie ist groß und wohnt in Yaoundé. Meine Eltern sind verheiratet und haben vier Kinder. Ich habe zwei Brüder und eine Schwester. Mein Vater arbeitet als Lehrer und meine Mutter ist Verkäuferin. Am Sonntag besuchen wir unsere Großeltern im Dorf. Ich liebe meine Familie sehr.}

« Ma famille est nombreuse et habite à Yaoundé. Mes parents sont mariés et ont quatre enfants. J'ai deux frères et une sœur. Mon père travaille comme enseignant et ma mère est vendeuse. Le dimanche, nous rendons visite à nos grands-parents au village. J'aime beaucoup ma famille. »

\textbf{Conclusion.} Le texte contient : verbes conjugués correctement (\all{ist, sind, haben, habe, arbeitet, besuchen, liebe}), une inversion (\all{Am Sonntag besuchen wir}), des possessifs accordés (\all{meine, unsere}) et le vocabulaire du thème.
\end{exoresolu}

\begin{attention}[Les 5 erreurs qui coûtent des points au Bac]
\begin{enumerate}
\item \textbf{La place du verbe.} Après un complément initial, le sujet doit suivre le verbe : écrivez \all{Am Sonntag besuchen wir die Großeltern} et jamais \all{Am Sonntag wir besuchen\dots} Le verbe conjugué est toujours en 2\textsuperscript{e} position.
\item \textbf{La majuscule des noms.} Tous les noms prennent la majuscule : \all{die Familie, das Kind, der Bruder}. Écrire \all{die familie} est une faute d'orthographe sanctionnée.
\item \textbf{Le genre et l'article.} Ne calquez pas le français : « la mère » se dit \all{die Mutter}, mais « le problème » se dit \all{das Problem}. Apprenez chaque nom avec son article.
\item \textbf{Les faux amis et calques.} \all{allein erziehend} signifie « qui élève seul son enfant » (parent isolé), pas « seul » en général. « J'ai 15 ans » se traduit \all{Ich bin 15 Jahre alt} (avec \all{sein}), et non avec \all{haben}.
\item \textbf{La conjugaison de sein et haben.} Les formes \all{er ist} et \all{er hat} sont irrégulières : jamais \all{er seint} ni \all{er habt}. De même \all{du hast} (sans \textbf{b}) et \all{du bist}.
\end{enumerate}
\end{attention}

\newpage

\coursec{Exercices}

\coursec{ALA01 : Les types de familles}

\courssub{Texte support : Familien heute}

\begin{textebox}[Familien heute — Texte de compréhension]
In Deutschland und in Kamerun gibt es viele verschiedene Familienformen. Die klassische Kernfamilie besteht aus Vater, Mutter und ihren gemeinsamen Kindern. Doch heute leben viele Menschen in anderen Konstellationen. Eine Alleinerziehende oder ein Alleinerziehender betreut die Kinder ohne Partner. Oft sind die Eltern geschieden oder waren nie verheiratet. Eine Patchworkfamilie entsteht, wenn getrennte Eltern neue Partner mit eigenen Kindern in die Familie bringen. Dann gibt es Stiefgeschwister und Halbgeschwister. In Kamerun ist die Großfamilie traditionell sehr wichtig: Großeltern, Onkel, Tanten und Cousins wohnen oft zusammen oder in der Nähe und unterstützen sich gegenseitig. Auch in Deutschland ziehen Großeltern manchmal zu ihren Enkeln, um bei der Betreuung zu helfen. Egal welche Form: In jeder Familie zählen Liebe, Verantwortung und Zusammenhalt.
\end{textebox}

\courssub{Compréhension du texte}

\begin{exercice}[Compréhension globale et détaillée]
Réponds en allemand par des phrases complètes.
\begin{enumerate}
\item Nenne drei verschiedene Familienformen, die im Text vorkommen.
\item Was ist eine Patchworkfamilie? Erkläre mit eigenen Worten.
\item Welche Rolle spielen die Großeltern in Kamerun und in Deutschland?
\item Finde im Text das Adjektiv, das „verantwortungsvoll“ bedeutet.
\item Übersetze den letzten Satz ins Französische.
\end{enumerate}
\end{exercice}

\courssub{Lexique thématique : die Familie}

\begin{definition}[Vocabulaire de base — die Familie]
Voici les mots essentiels du thème. \textbf{Apprends chaque nom avec son article et son pluriel.}
\end{definition}

\begin{center}
\begin{tabularx}{\linewidth}{|X|X|X|}
\hline
\rowcolor{popblueL} \textbf{Deutsch (Singulier)} & \textbf{Pluriel} & \textbf{Français} \\
\hline
die Familie & die Familien & la famille \\
der Vater & die Väter & le père \\
die Mutter & die Mütter & la mère \\
die Eltern & (kein Singular) & les parents \\
das Kind & die Kinder & l'enfant \\
der Sohn & die Söhne & le fils \\
die Tochter & die Töchter & la fille \\
der Bruder & die Brüder & le frère \\
die Schwester & die Schwestern & la sœur \\
die Geschwister & (kein Singular) & les frères et sœurs \\
der Großvater / der Opa & die Großväter / die Opas & le grand-père \\
die Großmutter / die Oma & die Großmütter / die Omas & la grand-mère \\
die Großeltern & (kein Singular) & les grands-parents \\
der Onkel & die Onkel & l'oncle \\
die Tante & die Tanten & la tante \\
der Cousin / die Cousine & die Cousins / die Cousinen & le cousin / la cousine \\
\hline
\end{tabularx}
\end{center}

\begin{definition}[États civils et types de familles]
\end{definition}

\begin{center}
\begin{tabularx}{\linewidth}{|X|X|X|}
\hline
\rowcolor{poporangeL} \textbf{Deutsch} & \textbf{Français} & \textbf{Remarque} \\
\hline
verheiratet & marié(e) & Participe II adjectivé \\
ledig & célibataire & \\
geschieden & divorcé(e) & \\
verwitwet & veuf/veuve & \\
allein erziehend & parent isolé / mono-parental & Adjectif invariable (souvent) \\
die Patchworkfamilie & la famille recomposée & Composé : Patchwork + Familie \\
der Stiefvater / die Stiefmutter & le beau-père / la belle-mère & Stief- = beau- (par alliance) \\
der Stiefbruder / die Stiefschwester & le demi-frère / la demi-sœur (par alliance) & \\
der Halbbruder / die Halbschwester & le demi-frère / la demi-sœur (un parent commun) & Halb- = demi- (par sang) \\
das Enkelkind & le petit-enfant & Neutre ! \\
\hline
\end{tabularx}
\end{center}

\begin{aretenir}[Familles de mots — Wortfamilien]
\begin{itemize}
\item \textbf{eltern} : die Eltern (parents) — elterlich (parental) — die Elternzeit (congé parental).
\item \textbf{verwandt} : verwandt (parenté) — der Verwandte / die Verwandte (le/la parent) — die Verwandtschaft (la parenté).
\item \textbf{heiraten} : heiraten (se marier) — die Hochzeit (le mariage) — der Ehemann / die Ehefrau (le mari / la femme).
\item \textbf{erziehen} : erziehen (élever) — die Erziehung (l'éducation) — der Erzieher / die Erzieherin (l'éducateur/trice).
\end{itemize}
\end{aretenir}

\courssub{Grammaire 1 : L'ordre des mots et les articles/possessifs}

\begin{propriete}[L'ordre des mots dans la phrase simple (Position 2)]
En allemand, le verbe conjugué occupe \textbf{toujours la deuxième position} (Position 2) dans la phrase principale affirmative. Le sujet peut être en position 1 (ordre normal) ou après le verbe (inversion) si un complément (temps, lieu, objet) commence la phrase.
\begin{center}
\begin{tabular}{|l|l|l|l|}
\hline
\rowcolor{popgreenL} \textbf{Position 1} & \textbf{Position 2 (Verbe)} & \textbf{Sujet (si pas en 1)} & \textbf{Reste de la phrase} \\
\hline
Meine Familie & wohnt &  & in Yaoundé. \\
In Yaoundé & wohnt & meine Familie & . \\
Heute & besucht & meine Tante & uns. \\
\hline
\end{tabular}
\end{center}
\cle{Règle d'or :} Le verbe ne bouge pas de la 2e place. Si un élément précède le sujet, le sujet glisse juste après le verbe (inversion).
\end{propriete}

\begin{attention}[Piège francophone : l'inversion]
En français, on dit : „Aujourd'hui, ma tante nous rend visite.“ (Sujet avant verbe).
En allemand : \all{Heute besucht meine Tante uns.} (Verbe avant sujet).
\emph{Jamais :} „Heute meine Tante besucht uns.“
\end{attention}

\begin{propriete}[Articles définis et indéfinis (Nominatif / Accusatif)]
Pour cette leçon, nous avons besoin du Nominatif (sujet) et de l'Accusatif (objet direct).
\end{propriete}

\begin{center}
\begin{tabular}{|l|c|c|c|c|}
\hline
\rowcolor{poppurpleL} \textbf{Cas} & \textbf{Maskulin} & \textbf{Féminin} & \textbf{Neutre} & \textbf{Pluriel} \\
\hline
Nom. défini & der Vater & die Mutter & das Kind & die Eltern \\
Akk. défini & \textbf{den} Vater & die Mutter & das Kind & die Eltern \\
\hline
Nom. indéfini & ein Vater & eine Mutter & ein Kind & (keine) Eltern \\
Akk. indéfini & \textbf{einen} Vater & eine Mutter & ein Kind & (keine) Eltern \\
\hline
\end{tabular}
\end{center}

\begin{aretenir}[Accusatif : seul le masculin change !]
\cle{der $\to$ den} / \cle{ein $\to$ einen}. Féminin, neutre, pluriel : inchangés.
\end{aretenir}

\begin{propriete}[Possessifs (Nominatif / Accusatif)]
Les possessifs se déclinent comme l'article indéfini (\cle{ein, eine, ein}).
\end{propriete}

\begin{center}
\begin{tabular}{|l|c|c|c|c|}
\hline
\rowcolor{poppinkL} \textbf{Personne} & \textbf{Maskulin} & \textbf{Féminin} & \textbf{Neutre} & \textbf{Pluriel} \\
\hline
ich (mein) & mein / \textbf{meinen} & meine & mein & meine \\
du (dein) & dein / \textbf{deinen} & deine & dein & deine \\
er/es (sein) & sein / \textbf{seinen} & seine & sein & seine \\
sie (ihr) & ihr / \textbf{ihren} & ihre & ihr & ihre \\
wir (unser) & unser / \textbf{unseren} & unsere & unser & unsere \\
ihr (euer) & euer / \textbf{euren} & eure & euer & eure \\
sie/Sie (ihr/Ihr) & ihr / \textbf{ihren} & ihre & ihr & ihre \\
\hline
\end{tabular}
\end{center}

\begin{exemplebox}[Exemples d'utilisation]
\all{Mein Vater wohnt in Douala.} (Nominativ, masc.) \\
\all{Ich besuche meinen Vater.} (Akkusativ, masc. $\to$ \textbf{-en}) \\
\all{Meine Mutter ist Lehrerin.} (Nominativ, fém.) \\
\all{Ich liebe meine Mutter.} (Akkusativ, fém. $\to$ inchangé) \\
\all{Unser Haus ist groß.} (Nominativ, neutre) \\
\all{Wir verkaufen unser Haus.} (Akkusativ, neutre $\to$ inch
\end{exemplebox}


\newpage

\coursec{Corrigés des exercices}

\begin{corrige}[Exercice 1 — Compréhension globale et détaillée]
\textbf{Rappel de la consigne :} répondre en allemand par des phrases complètes.
\begin{enumerate}
\item \textbf{Frage :} Nenne die im Text genannten Familienformen.\\
\all{Im Text kommen die Kernfamilie, die Familie mit einem alleinerziehenden Elternteil, die Patchworkfamilie und die Großfamilie vor.}\\
« Dans le texte apparaissent la famille nucléaire, la famille monoparentale, la famille recomposée et la famille élargie. »\\
\emph{Justification :} les quatre formes sont explicitement nommées dans le texte (\all{Kernfamilie}, \all{Alleinerziehende} / \all{alleinerziehende Eltern}, \all{Patchworkfamilie}, \all{Großfamilie}). \cle{Compréhension détaillée}

\item \textbf{Frage :} Was ist eine Patchworkfamilie?\\
\all{Eine Patchworkfamilie entsteht, wenn getrennte Eltern neue Partner haben und diese Partner ihre eigenen Kinder mit in die Familie bringen.}\\
« Une famille recomposée naît quand des parents séparés ont de nouveaux partenaires et que ceux-ci amènent leurs propres enfants dans la famille. »\\
\emph{Point de grammaire :} dans la subordonnée introduite par \all{wenn}, le verbe conjugué (\all{haben}) va à la fin : \all{..., wenn getrennte Eltern neue Partner haben.} \cle{Subordonnée / Verbe en fin}

\item \textbf{Frage :} Welche Rolle spielen die Großeltern in Kamerun und in Deutschland?\\
\all{In Kamerun wohnen die Großeltern oft mit der Familie zusammen oder in der Nähe und unterstützen die Familie. In Deutschland ziehen sie manchmal zu ihren Enkeln und helfen bei der Betreuung.}\\
« Au Cameroun, les grands-parents habitent souvent avec la famille ou à proximité et la soutiennent. En Allemagne, ils emménagent parfois chez leurs petits-enfants et aident à la garde. » \cle{Comparaison culturelle}

\item \textbf{Frage :} Finde im Text das Adjektiv, das „verantwortungsvoll“ bedeutet.\\
Le texte ne contient pas d'adjectif correspondant ; il contient le \textbf{nom} \all{die Verantwortung} (la responsabilité). L'adjectif dérivé est \all{verantwortungsvoll} (ou \all{verantwortlich}). Réponse attendue : \all{Das Wort im Text ist „Verantwortung“ (ein Nomen); das Adjektiv dazu ist „verantwortungsvoll“.} \cle{Dérivation / Faux ami}

\item \textbf{Frage :} Übersetze den letzten Satz ins Französische.\\
\all{Egal welche Form: In jeder Familie zählen Liebe, Verantwortung und Zusammenhalt.}\\
« Quelle que soit sa forme : dans chaque famille, ce qui compte, c'est l'amour, la responsabilité et la solidarité. »\\
\emph{Points de vigilance :} \all{zählen} = « compter » au sens de « avoir de l'importance » (et non « compter » = \all{zählen} nombres) ; \all{der Zusammenhalt} = « la cohésion, la solidarité » ; \all{egal} dans \all{egal welche Form} = « quelle que soit » (ne pas traduire par « égal » → \cle{faux ami partiel}).
\end{enumerate}
\textbf{Barème indicatif (10 pts) :} 2 pts par question (1 pt pour le contenu, 1 pt pour la correction grammaticale : verbe en position 2, majuscules aux noms, cas).
\end{corrige}

\begin{corrige}[Exercice 2 — Vocabulaire : articles et pluriels]
\textbf{Consigne type :} complète avec l'article défini (\all{der}, \all{die}, \all{das}) et donne le pluriel.
\begin{enumerate}
\item \all{die Familie} — Pluriel : \all{die Familien}
\item \all{der Vater} — Pluriel : \all{die Väter} (\cle{Umlaut} !)
\item \all{die Mutter} — Pluriel : \all{die Mütter} (\cle{Umlaut} !)
\item \all{das Kind} — Pluriel : \all{die Kinder}
\item \all{der Bruder} — Pluriel : \all{die Brüder} (\cle{Umlaut} !)
\item \all{die Schwester} — Pluriel : \all{die Schwestern}
\item \all{das Enkelkind} — Pluriel : \all{die Enkelkinder}
\item \all{die Tante} — Pluriel : \all{die Tanten}
\item \all{der Onkel} — Pluriel : \all{die Onkel}
\item \all{der Sohn} — Pluriel : \all{die Söhne} (\cle{Umlaut} !)
\end{enumerate}
\textbf{Points de vigilance :}
\begin{itemize}
\item Les noms de personnes en \all{-er} prennent souvent l'\cle{Umlaut} au pluriel (\all{Vater $\to$ Väter}, \all{Mutter $\to$ Mütter}, \all{Bruder $\to$ Brüder}, \all{Sohn $\to$ Söhne}).
\item \all{die Eltern} (les parents) et \all{die Geschwister} (les frères et sœurs) n'existent qu'au pluriel (\emph{pluralia tantum}).
\item \all{das Kind} est neutre (\cle{genre neutre}), même s'il désigne une fille ou un garçon.
\item Majuscule obligatoire à tous les noms : \all{die Familie}, \all{der Vater}, \all{das Kind}.
\end{itemize}
\end{corrige}

\begin{corrige}[Exercice 3 — L'ordre des mots : remets la phrase en ordre]
\textbf{Consigne type :} forme des phrases correctes ; attention à la position du verbe (règle de la \cle{position 2}).
\begin{enumerate}
\item (wohnt / in Yaoundé / meine Familie) $\to$ \all{Meine Familie wohnt in Yaoundé.} « Ma famille habite à Yaoundé. »
\item (besucht / heute / uns / meine Tante) $\to$ \all{Heute besucht uns meine Tante.} (ordre standard : complément de temps — verbe — objet pronom — sujet) « Aujourd'hui, ma tante nous rend visite. »
\item (am Sonntag / die Großeltern / kommen / zu Besuch) $\to$ \all{Am Sonntag kommen die Großeltern zu Besuch.} « Dimanche, les grands-parents viennent en visite. »
\item (in Deutschland / leben / viele Patchworkfamilien) $\to$ \all{In Deutschland leben viele Patchworkfamilien.} « En Allemagne vivent beaucoup de familles recomposées. »
\item (mein Bruder und ich / spielen / am Nachmittag / Fußball) $\to$ \all{Am Nachmittag spielen mein Bruder und ich Fußball.} « L'après-midi, mon frère et moi jouons au football. »
\end{enumerate}
\textbf{Justification grammaticale :} le verbe conjugué reste en \textbf{position 2}. Quand un complément de temps ou de lieu ouvre la phrase (\all{Heute}, \all{Am Sonntag}, \all{In Deutschland}, \all{Am Nachmittag}), le sujet passe \textbf{après} le verbe : c'est l'\cle{inversion}.\\
\textbf{Faute typique du francophone :} *\all{Heute meine Tante besucht uns.} — à proscrire absolument (verbe en position 3).
\end{corrige}

\begin{corrige}[Exercice 4 — Articles et possessifs au Nominativ et Akkusativ]
\textbf{Consigne type :} complète avec l'article ou le possessif correct (terminaison).
\begin{enumerate}
\item \all{Der Vater arbeitet in Douala.} (Nominativ, masculin, inchangé)
\item \all{Ich besuche \textbf{meinen} Vater.} (Akkusativ masculin : \all{mein $\to$ meinen})
\item \all{\textbf{Die} Mutter kocht heute.} (Nominativ féminin)
\item \all{Wir lieben \textbf{unsere} Mutter.} (Akkusativ féminin : inchangé)
\item \all{\textbf{Das} Kind spielt im Garten.} (Nominativ neutre)
\item \all{Sie hat \textbf{ein} Kind.} (Akkusativ neutre : inchangé)
\item \all{\textbf{Meine} Geschwister wohnen in Bafoussam.} (Pluriel : \all{meine})
\item \all{Er besucht \textbf{seinen} Onkel.} (Akkusativ masculin : \all{sein $\to$ seinen})
\end{enumerate}
\end{corrige}

\begin{propriete}[Règle de l'Akkusatif]
À l'accusatif, \textbf{seul le masculin change} :
\all{der $\to$ den}, \quad \all{ein $\to$ einen}, \quad \all{mein/dein/sein/ihr/unser/euer/ihr $\to$ meinen/deinen/seinen/ihren/unseren/euren/ihren}.
Féminin (\all{die/eine/meine}), neutre (\all{das/ein/mein}) et pluriel (\all{die/-/meine}) restent identiques au nominatif.
\end{propriete}


\begin{corrige}[Exercice 5 — Conjugaison au présent : verbes faibles, forts, sein et haben]
\textbf{Consigne type :} conjugue le verbe entre parenthèses au présent.
\begin{enumerate}
\item \all{Ich \textbf{wohne} in Yaoundé.} (wohnen, \cle{verbe faible} : ich wohne)
\item \all{Du \textbf{kommst} aus Kamerun.} (kommen, \cle{verbe faible} : du kommst)
\item \all{Er \textbf{spricht} Deutsch und Französisch.} (sprechen, \cle{verbe fort} : \all{e $\to$ i} aux 2e et 3e pers. du sg.)
\item \all{Wir \textbf{haben} drei Kinder.} (haben, \cle{verbe irrégulier} : wir haben)
\item \all{Ihr \textbf{seid} eine große Familie.} (sein, \cle{verbe irrégulier} : ihr seid)
\item \all{Meine Großeltern \textbf{lesen} gern.} (lesen, \cle{verbe fort} ; 3e pers. pl. = forme de base)
\item \all{Sie \textbf{ist} allein erziehend.} (sein : sie ist)
\item \all{Du \textbf{hast} viele Cousins.} (haben : du hast)
\end{enumerate}
\textbf{Tableau de référence :}
\begin{center}
\begin{tabular}{|l|l|l|l|l|}
\hline
\rowcolor{popgreenL} \textbf{Personne} & \textbf{wohnen} (faible) & \textbf{sprechen} (fort) & \textbf{sein} (irrég.) & \textbf{haben} (irrég.) \\
\hline
ich & wohne & spreche & bin & habe \\
du & wohnst & \textbf{sprichst} & bist & hast \\
er/sie/es & wohnt & \textbf{spricht} & ist & hat \\
wir & wohnen & sprechen & sind & haben \\
ihr & wohnt & sprecht & seid & habt \\
sie/Sie & wohnen & sprechen & sind & haben \\
\hline
\end{tabular}
\end{center}
\textbf{Points de vigilance :}
\begin{itemize}
\item Les \cle{verbes forts} à voyelle radicale \all{-e-} la changent en \all{-i-} (ou \all{-ie-}) à \all{du} et \all{er/sie/es} : \all{sprechen $\to$ spricht}, \all{lesen $\to$ liest}, \all{essen $\to$ isst}, \all{geben $\to$ gibt}, \all{treffen $\to$ trifft}.
\item \all{sein} et \all{haben} sont \cle{irréguliers} et doivent être appris par cœur (tableau ci-dessus).
\item Attention à \all{du hast} (pas *\all{du habst}) et \all{er hat} (pas *\all{er habt}).
\end{itemize}
\end{corrige}

\begin{corrige}[Exercice 6 — Production écrite guidée : ma famille]
\textbf{Consigne type :} rédige un court paragraphe (5 à 8 phrases) sur ta famille en utilisant le vocabulaire et la grammaire de la leçon.

\textbf{Proposition de rédaction modèle :}

\all{Meine Familie ist nicht sehr groß. Mein Vater heißt Paul und meine Mutter heißt Chantal. Ich habe einen Bruder und eine Schwester. Wir wohnen in Yaoundé. Meine Großeltern leben im Dorf, aber sie besuchen uns oft. Am Sonntag essen wir zusammen. Ich liebe meine Familie, denn bei uns zählen Liebe und Zusammenhalt.}

« Ma famille n'est pas très grande. Mon père s'appelle Paul et ma mère s'appelle Chantal. J'ai un frère et une sœur. Nous habitons à Yaoundé. Mes grands-parents vivent au village, mais ils nous rendent souvent visite. Le dimanche, nous mangeons ensemble. J'aime ma famille, car chez nous l'amour et la solidarité comptent. »

\textbf{Barème indicatif (10 pts) :}
\begin{itemize}
\item Respect de la consigne et du nombre de phrases (5--8) : 2 pts
\item Vocabulaire du thème (au moins 5 mots de la famille : \all{Vater, Mutter, Bruder, Schwester, Großeltern, Onkel, Tante, Kind...}) : 2 pts
\item Verbe en position 2 dans chaque phrase principale : 2 pts
\item Conjugaison correcte au présent (\all{sein, haben}, verbes faibles/forts) : 2 pts
\item Articles et possessifs corrects (cas), majuscules à tous les noms : 2 pts
\end{itemize}

\textbf{Points de vigilance (check-list avant de rendre) :}
\begin{itemize}
\item \cle{Majuscules} : \all{der Vater}, \all{die Familie}, \all{das Dorf}, \all{der Sonntag}, \all{die Liebe}, \all{der Zusammenhalt}.
\item Conjonctions de coordination (\all{aber, denn, und, oder}) : le verbe reste en position 2 : \all{..., aber sie besuchen uns oft.} / \all{..., denn bei uns zählen Liebe...}
\item Ordre des mots : si la phrase commence par un complément (\all{Am Sonntag}), inversion : \all{Am Sonntag essen wir...} (pas *\all{Am Sonntag wir essen...}).
\item Accusatif masculin : \all{ich habe \textbf{einen} Bruder}, \all{ich besuche \textbf{meinen} Onkel} (terminaison \textbf{-en} sur article/possessif).
\item \all{allein erziehend} (séparé) ou \all{alleinerziehend} (collé, adjectif) — les deux existent, \all{alleinerziehend} est plus courant comme adjectif attributif.
\end{itemize}
\end{corrige}

\newpage

\coursec{Fiche bilan}

\begin{popfigure}
\begin{tikzpicture}[mindmap, grow cyclic, every node/.style={concept, align=center, font=\small},
  root concept/.append style={concept color=popblue, font=\bfseries\small, minimum size=2.6cm},
  level 1/.append style={level distance=3.4cm, sibling angle=51.4},
  level 1 concept/.append style={minimum size=2.1cm}]
\node[root concept, align=center]{Die Familie\\Types de familles}
  child[concept color=poporange]{ node[align=center]{Lexique\\die Familie,\\die Eltern,\ldots} }
  child[concept color=popgreen]{ node[align=center]{Ordre des mots\\Verbe en\\2\up{e} position} }
  child[concept color=poppurple]{ node[align=center]{Articles\\der, die, das\\ein, eine} }
  child[concept color=popgold]{ node[align=center]{Possessifs\\mein, dein,\\sein, ihr} }
  child[concept color=poppink]{ node[align=center]{Présent\\verbes faibles\\et forts} }
  child[concept color=popblue!60]{ node[align=center]{sein / haben\\bin, ist, sind\\habe, hat} }
  child[concept color=popgreen!60]{ node[align=center]{Production\\résumé,\\court texte} };
\end{tikzpicture}
\legende{Carte mentale de la leçon : les types de familles}
\end{popfigure}

\begin{bilan}
\textbf{1. Lexique clé (à connaître avec article et pluriel)}
\begin{itemize}\setlength{\itemsep}{1pt}
  \item \all{die Familie, -n} : la famille ; \all{der Familienstand} : l'état civil
  \item \all{die Eltern} (plur.) : les parents ; \all{der Vater, die Väter} : le père ; \all{die Mutter, die Mütter} : la mère
  \item \all{die Großeltern} (plur.) : les grands-parents ; \all{der Großvater / die Großmutter} : le grand-père / la grand-mère
  \item \all{die Geschwister} (plur.) : les frères et sœurs ; \all{der Bruder, die Brüder} ; \all{die Schwester, -n}
  \item \all{das Kind, -er} : l'enfant ; \all{der Sohn, die Söhne} ; \all{die Tochter, die Töchter}
  \item \all{verheiratet} : marié(e) ; \all{geschieden} : divorcé(e) ; \all{ledig} : célibataire
  \item \all{allein erziehend} : parent isolé ; \all{die Alleinerziehende, -n} : la mère (ou le père) seul(e)
  \item \all{die Großfamilie, -n} : la famille élargie ; \all{die Kleinfamilie / Kernfamilie} : la famille nucléaire
  \item \all{die Patchworkfamilie, -n} : la famille recomposée ; \all{das Einzelkind, -er} : l'enfant unique
\end{itemize}

\textbf{2. Grammaire à connaître par cœur}
\begin{center}
\begin{tabularx}{\linewidth}{|X|X|}
\hline
\rowcolor{popblueL} \textbf{Règle} & \textbf{Exemple} \\
\hline
Le verbe conjugué est toujours en \cle{2\up{e} position} dans la phrase déclarative. & \all{Meine Familie \textbf{wohnt} in Yaoundé.} \\
\hline
Inversion sujet--verbe si un autre élément ouvre la phrase. & \all{\textbf{Heute} \textbf{besucht} mein Bruder uns.} \\
\hline
Articles définis : \all{der} (masc.), \all{die} (fém.), \all{das} (neutre) ; indéfinis : \all{ein, eine, ein}. & \all{der Vater, die Mutter, das Kind} \\
\hline
Possessifs : \all{mein, dein, sein, ihr, unser, euer, ihr/Ihr} ; ils se comportent comme \all{ein}. & \all{\textbf{Meine} Mutter ist Lehrerin.} \\
\hline
Présent des verbes faibles : radical + \all{-e, -st, -t, -en, -t, -en}. & \all{ich wohne, du wohnst, er wohnt} \\
\hline
Verbes forts : changement de voyelle à la 2\up{e}/3\up{e} pers. du singulier (\all{e$\to$i/ie}, \all{a$\to$ä}). & \all{ich spreche, du \textbf{sprichst}, er \textbf{spricht}} ; \all{du \textbf{lädst}, er \textbf{lädt}} \\
\hline
\all{sein} : bin, bist, ist, sind, seid, sind ; \all{haben} : habe, hast, hat, haben, habt, haben. & \all{Wir \textbf{sind} eine Großfamilie. Sie \textbf{hat} zwei Kinder.} \\
\hline
\end{tabularx}
\end{center}

\textbf{3. Méthodes express}
\begin{itemize}\setlength{\itemsep}{1pt}
  \item \textbf{Comprendre un texte} : 1) lire le titre et anticiper ; 2) lecture globale (qui ? où ? quoi ?) ; 3) lecture détaillée en soulignant le vocabulaire du thème ; 4) répondre par des phrases complètes en allemand.
  \item \textbf{Répondre aux questions} : reprendre les mots de la question et conjuguer correctement le verbe en 2\up{e} position.
  \item \textbf{Produire un court paragraphe} : 4 à 6 phrases simples (sujet + verbe + compléments), au présent, avec possessifs et vocabulaire de la famille.
\end{itemize}
\end{bilan}

\begin{aretenir}[Checklist avant le Bac]
\begin{itemize}\setlength{\itemsep}{2pt}
  \item[$\square$] Je connais l'article et le pluriel des noms de la famille (\all{der Vater, die Väter}).
  \item[$\square$] Je sais décrire les types de familles : \all{Kleinfamilie, Großfamilie, Patchworkfamilie, allein erziehend}.
  \item[$\square$] Je place toujours le verbe conjugué en 2\up{e} position.
  \item[$\square$] J'inverse sujet et verbe après un complément en tête de phrase (\all{Heute kommt\ldots}).
  \item[$\square$] Je choisis le bon article : \all{der / die / das}, \all{ein / eine / ein}.
  \item[$\square$] J'emploie correctement les possessifs : \all{mein Vater, meine Mutter, mein Kind}.
  \item[$\square$] Je conjugue sans erreur les verbes faibles au présent (\all{wohnen, machen, lernen}).
  \item[$\square$] Je maîtrise les changements de voyelle des verbes forts (\all{sprechen$\to$spricht, lesen$\to$liest, fahren$\to$fährt}).
  \item[$\square$] Je conjugue \all{sein} et \all{haben} par cœur.
  \item[$\square$] Je mets une majuscule à tous les noms allemands.
  \item[$\square$] Je sais répondre à des questions de compréhension par des phrases complètes.
  \item[$\square$] Je sais rédiger un court paragraphe (4 à 6 phrases) sur ma famille.
\end{itemize}
\end{aretenir}

\findecours

\end{document}
